\documentclass[11pt]{article}
\usepackage[T1]{fontenc}
\usepackage[utf8]{inputenc}
\usepackage{lmodern,amsmath,amssymb,amsthm,mathtools}
\usepackage[margin=1in]{geometry}
\usepackage{booktabs,longtable,array,microtype,enumitem}
\usepackage[hidelinks]{hyperref}
\setlist{itemsep=3pt,topsep=5pt}
\setlength{\emergencystretch}{2em}
\newtheorem{theorem}{Theorem}[section]
\newtheorem{proposition}[theorem]{Proposition}
\newtheorem{lemma}[theorem]{Lemma}
\newtheorem{corollary}[theorem]{Corollary}
\theoremstyle{definition}
\newtheorem{definition}[theorem]{Definition}
\theoremstyle{remark}
\newtheorem{remark}[theorem]{Remark}
\newcommand{\Q}{\mathbb Q}
\newcommand{\Z}{\mathbb Z}
\newcommand{\N}{\mathbb N}
\newcommand{\F}{\mathbb F}
\newcommand{\Cset}{\mathcal C}
\providecommand{\mathscr}[1]{\mathcal{#1}}
\newcommand{\Sel}{\operatorname{Sel}}
\newcommand{\Sha}{\operatorname{Sha}}
\newcommand{\ord}{\operatorname{ord}}
\newcommand{\Norm}{\operatorname{N}}
\newcommand{\Tr}{\operatorname{Tr}}
\newcommand{\Trd}{\operatorname{Trd}}
\newcommand{\Close}{\operatorname{Close}}
\newcommand{\Test}{\operatorname{Test}}
\newcommand{\Val}{\operatorname{Val}}
\newcommand{\tw}[1]{2\mathbin{\uparrow\!\uparrow}#1}
\title{$\mathbb{Z}$ is Diophantine in $\mathbb{Q}$}
\author{}
\date{November 4, 2026}
\begin{document}
\maketitle
\begin{abstract}
We construct an integer-coefficient polynomial in
$292\,890\,231\,905$ variables whose rational zero projection is exactly
$\Z$. Its degree is less than $2\mathbin{\uparrow\!\uparrow}7$, and the
absolute value of every coefficient is less than
$2\mathbin{\uparrow\!\uparrow}8$.

Well-known consequences of $\Z$ being Diophantine in $\Q$ include a refutation of Mazur's conjecture that the topological closure of the rational points of a variety over $\Q$ has finitely many components, and existence of a universal rational Diophantine polynomial. Moreover, the set $\operatorname{HTP}(\Q)$ of integer polynomials with a rational zero is many-one equivalent to the halting set $0'$.
\end{abstract}
\tableofcontents

% Begin 01-specification-model.tex
\section{The main theorem and proof structure}\label{sec:main}
This construction builds on OpenAI's preprint
\emph{Hilbert's tenth problem over the rational numbers}~\cite{openaihtp}.

We use $\N=\{0,1,2,\ldots\}$. For a reduced rational number $a/b$, with
$b>0$, put $h(a/b)=\log\max(|a|,b)$. Define $\tw1=2$ and
$\tw{(r+1)}=2^{\tw r}$; every tower is associated to the right.

\begin{theorem}[A polynomial definition of the integers]\label{thm:target}
There exists one polynomial
\[
 P\in\Z[X_0,X_1,\ldots,X_{292890231904}]
\]
such that
\begin{align*}
 \deg P&<\tw7,\\
 |[X^\alpha]P|&<\tw8\qquad\text{for every multi-index }\alpha,\\
 q\in\Z&\ \Longleftrightarrow\
 \exists(w_1,\ldots,w_{292890231904})\in\Q^{292890231904},\quad
 P(q,w_1,\ldots,w_{292890231904})=0\quad(q\in\Q).
\end{align*}
All conclusions concern the same polynomial. Here $q\in\Z$ means the
ordinary embedded integers. The ambient variable count does not assert that
every variable occurs essentially.
\end{theorem}

The file \nolinkurl{ordinary-integer-definition.json} specifies the
polynomial $P_*$. The argument first represents integer indices by points
on an elliptic curve over $\Q(\sqrt2)$. A separate value relation identifies
the ordinary rational value of an index. The point coordinates and the
integer index have distinct roles throughout.

Theorem~\ref{thm:target} refutes Mazur's conjecture that the closure of
the rational points of a variety in its real locus has finitely many
connected components. It also gives a universal rational Diophantine
polynomial, with
$\operatorname{HTP}(\Q)$ many-one equivalent to $0'$.
Section~\ref{sec:four-corollaries} proves these three consequences.

The final composition uses the following precise statements.
\begin{description}[style=nextline,leftmargin=1.8cm]
\item[$\mathsf G$] Canonical point codes form exactly the infinite cyclic
group $2E(\Q(\sqrt2))=\Z P_0$. The membership polynomial types arbitrary inputs;
raw addition describes the group law on canonical codes.
\item[$\mathsf M$] The actual guarded multiplication template has rational
witnesses exactly for triples with indices $a,b,ab$.
\item[$\mathsf H$] The six-contact height theorem holds with the constants
used by the precision computation and, separately, by multiplication.
\item[$\mathsf{PC}$] The explicit pole polynomial is soluble for every
input in Theorem~\ref{hyp:PC}.
\item[$\mathsf C$] The actual five-input natural polynomial and its
rational translation have the exact computation and substitution semantics
of Section~\ref{sec:compiler}, including the power template.
\item[$\mathsf B$] The templates, fixed constructors and private witness
substitutions specify the same saved expression $P_*$. The quantitative
syntax calculation applies to that expression.
\end{description}

\begin{lemma}[Assembly]\label{lem:assembly}
Assume $\mathsf G,\mathsf M,\mathsf H,\mathsf{PC},\mathsf C,\mathsf B$.
Then $P_*$ satisfies Theorem~\ref{thm:target}.
\end{lemma}
\begin{proof}
Section~\ref{sec:bridge} proves the exact ordinary-value relation using
one shared shift and disjoint private witness blocks. The literal
substitutions in $\mathsf B$ identify its existential polynomial with
$P_*$. Section~\ref{sec:size} proves the stated bounds for that same
expression.
\end{proof}

\begin{proof}[Proof of Theorem~\ref{thm:target}]
Section~\ref{sec:model}, including Proposition~\ref{prop:generator},
proves $\mathsf G$ with its canonical membership and addition formulas.
Theorem~\ref{thm:height} proves $\mathsf H$ for each of the two sets of
constants used below. Sections~\ref{sec:rc}--\ref{sec:limits}, with
Appendices~\ref{app:even-source}--\ref{app:endpoints}, prove the restricted
converse in Theorem~\ref{hyp:RC}. The simultaneous descent in
Section~\ref{sec:pole} then proves $\mathsf{PC}$. Using these results,
Section~\ref{sec:mult} proves $\mathsf M$ on all rational input codes,
including the finite false-triple guard.

Section~\ref{sec:compiler} proves $\mathsf C$ for the fixed five-input
natural polynomial and its rational model translation. Its witness
variables remain natural before that translation. Sections~\ref{sec:syntax}
and~\ref{sec:size} and Appendix~\ref{app:literal} prove $\mathsf B$, including
the fixed scalars and private witness substitutions. All six statements
therefore hold, and Lemma~\ref{lem:assembly} applies to the specified $P_*$.
\end{proof}

\section{A finite specification of the polynomial}\label{sec:syntax}
The saved specification is a $23\,180\,023$-byte JSON expression with SHA256
\begin{center}\small\ttfamily
a86790bbf1095d6b1cea9c66067ad8432dd35adc320254eecde21f9f98574714.
\end{center}
It contains 23 library templates and seven named recipes. The root recipe
is \texttt{integers}; its only parameter is the rational variable at index
zero. Every other variable is existentially quantified over $\Q$. Sources
S01--S03 fix the construction and its semantics. Stable paths and hashes
are in the companion \texttt{SOURCE\_MAP.json}.

\subsection{Literal operations and witness scope}
A scalar node is a variable, integer constant, sum, product, or a call to
an earlier template. A call substitutes its parameter expressions and a
private consecutive witness block into the callee's polynomial. Shared
nodes represent repeated subexpressions, not new equations. A recipe
sums squares of constraint polynomials or multiplies branch polynomials.
Over $\Q$,
\[
 \sum_i f_i^2=0\ \Longleftrightarrow\ \bigwedge_i f_i=0,
 \qquad\prod_i f_i=0\ \Longleftrightarrow\ \bigvee_i f_i=0.
\]
Private witness blocks can be moved to the front because the rational
domain is nonempty. A whole inactive product-call block may receive
arbitrary values. This does not discard fixed equations inside an active
callee. A rational constant $p/q$, $q\ne0$, is represented by a fresh
$c$ and the globally conjoined equation $qc=p$. In particular the cached
coordinates of $P_0$ satisfy
\[
 8g_0=1641,\qquad4g_1=-15465,\qquad32g_2=87483.
\]
These equations remain active outside internal disjunctions in the
multiplication, signed-encoding and exception templates. Inverses are
equations $uv=1$, never partial evaluation operations.

Strict positivity has the polynomial expression
\[
 s>0\ \Longleftrightarrow\
 \exists z_1,z_2,z_3,z_4\in\Q,\quad s\sum_i z_i^2=1.
\]
Indeed, if $1/s=a/b>0$, express $ab$ as four integer squares and divide
their coordinates by $b$. The reverse direction follows from the order
on $\Q$. Absolute value has the exact expression
\[
 a=|x|\quad\Longleftrightarrow\quad
 \exists z_1,\ldots,z_4\in\Q,\quad a=\sum_i z_i^2,\quad a^2=x^2.
\]

\subsection{Fixed finite constructors}
Three node types abbreviate finite polynomials. A
\texttt{fixed\_integer} inserts a fixed integer expression;
\texttt{psi\_component} inserts one coordinate of a division polynomial;
\texttt{finite\_guard\_product} inserts a product over a fixed box of
integers. Their bounds never depend on the free rational input.

For example, on $y^2=x^3+Ax$, write $\psi_n=f_n$ for odd $n$ and
$\psi_n=2yf_n$ for even $n$. The bases are
\begin{align*}
 f_0&=0,&f_1=f_2&=1,\\
 f_3&=3x^4+6Ax^2-A^2,&
 f_4&=2(x^6+5Ax^4-5A^2x^2-A^3).
\end{align*}
For $m\ge3$,
$f_{2m}=f_m(f_{m+2}f_{m-1}^2-f_{m-2}f_{m+1}^2)$.
For $m\ge2$, put $L=f_{m+2}f_m^3$, $R=f_{m-1}f_{m+1}^3$, and
$D_0=16(x^3+Ax)^2$. Then $f_{2m+1}=D_0L-R$ for even $m$ and
$L-D_0R$ for odd $m$. Negative indices satisfy $\psi_{-n}=-\psi_n$.
This is a terminating recursion.

The fixed prime catalogue consists of the first 22 primes above 324 in
each unit residue class modulo $W=4\prod_{p\le324}p$. Dirichlet's theorem
ensures termination. Put $N_s=(\prod_{p\le324}p)H_{\rm cat}$, where
$H_{\rm cat}$ is the catalogue product. Assign inverse abscissa zero
at the identity. For the normalized representation
$1/x([m]P_0)=(a+bs)/c$ with $a,b,c\in\Z$, $c>0$, and
$\gcd(a,b,c)=1$, define $M_0$ to be the least $m\ge1$ such that
$N_s^4\mid a,b$ and $\gcd(c,N_s)=1$ (S31).
This is a fixed integer, with a successful index bounded in
Section~\ref{sec:size}.
The fixed point coordinates in the guard have nonzero denominators by
the infinite order of $P_0$. Thus finite expansion gives one integer
polynomial without requiring a monomial expansion. Its arithmetic zero
set remains a separate issue. In particular, bounding the integer $T$
does not prove that it excludes all false multiplication triples.

\section{The elliptic model of integer indices}\label{sec:model}
Let
\[
 K_0=\Q(s),\quad s^2=2,\quad d=75-53s,\quad E:y^2=x^3-d^2x,
\]
and fix
\[
 P_0=\left(\frac{1641-1160s}{8},\frac{-123720+87483s}{32}\right).
\]
The identity has code $(1,0,0,0,0)$. An affine point has code
$(0,x_0,x_1,y_0,y_1)$ with $x=x_0+x_1s$, $y=y_0+y_1s$.
Linear independence of $1,s$ makes this unique. Write $\iota(n)$ for
the code of $[n]P_0$.

\begin{proposition}[Generator certificate]\label{prop:generator}
The point $P_0$ has infinite order and $2E(K_0)=\Z P_0$.
\end{proposition}
\begin{proof}
We give the argument of S05--S07, making its finite certificate explicit.
The ring of integers $\Z[s]$ has class number one because its Minkowski
bound is $\sqrt2<2$. It has four unit squareclasses. The sign vectors
of $-1,1+s$ are independent, so every totally positive unit is a square.

The $x$-squareclass map in the two-isogeny descent on $E$ has image
contained in the eight classes generated by $-1,1+s,d$. Away from $(d)$,
odd $v(x)$ makes the minimum of $3v(x)$ and $v(x)+2v(d)$ odd, impossible
for $y^2$. The torsion points give $1,-1,d,-d$. Put $\epsilon=1+s$,
$\delta=3-s$. The change $(x,y)\mapsto(\epsilon^4x,\epsilon^6y)$ takes
$E$ to $C:y^2=x^3-\delta^2x$, and $R=(1+s,2s)$ provides the remaining
unit class. Thus the first image has order eight.

On $E':y^2=x^3+4d^2x$, a nonzero $x$ is positive in both real embeddings.
Its valuations are even. At $(s)$ the two terms have valuations $3m$
and $m+4$, with equality only at $m=2$. At $(d)$ the only odd case
would require $u^2=-4$ in $\F_7$. Thus the second image is trivial.
The exact two-isogeny sequence gives $[E(K_0):2E(K_0)]=8$; its possible
extra kernel term vanishes since a nonkernel rational two-torsion point
maps to the dual kernel. Mordell--Weil gives rank one.

The good reductions at $(17,s-6)$ and $(23,s-5)$ have 20 and 24 points.
Prime-to-residue-characteristic torsion injects, so the torsion group
has order at most four, hence is exactly the displayed full two-torsion.
The $x$-class of $R$ has opposite real signs, whereas all torsion classes
have equal signs. Its free index is therefore odd.

The duplication forms on $C$ are
\[
 F_2(a,b)=(a^2+\delta^2b^2)^2,\qquad
 G_2(a,b)=4ab(a^2-\delta^2b^2).
\]
At either real embedding their maximum lies between
$\max(|a|,|b|)^4$ and $676\max(|a|,|b|)^4$. For primitive $a,b$ their
common ideal has valuations at most six at $(s)$ and four at $(\delta)$,
and none elsewhere. For the latter bound use that $-1$ is nonsquare
modulo seven. For the former, a local unit is $A+Bs$ with $A$ odd, and
$v_s((A+Bs)^2+1)$ is two or three according to the parity of $B$.
Hence
\[
 4h_x(T)-\log392\le h_x(2T)\le4h_x(T)+\log676.
\]
With $2\widehat h(T)=\lim_m4^{-m}h_x(2^mT)$ this gives
$h_x(T)\le2\widehat h(T)+\log392/3$. Direct doubling gives
$x(2R)=(57-28s)/8$, with height $\log41$. If $R=mS+T$ for torsion
$T$ and $|m|\ge3$, it follows that
\[
 h_x(S)\le\frac{\log(41^3\cdot676)}{108}+\frac{\log392}{3}<\log9.
\]
The strict inequality is the exact integer comparison
$41^3\cdot676\cdot392^{36}<9^{108}$.

A nonrational $x\in K_0$ in this height range has primitive polynomial
$aX^2+bX+c$ with $1\le a\le81$, $|b|\le162$, $|c|\le81$, and
discriminant $2k^2$, $1\le k\le162$. Both roots are included.
Rational possibilities have coprime numerator of absolute value at most
nine and denominator from one to nine. The exact finite square test of
$x^3-\delta^2x$ gives precisely
\[
 0,\quad\delta,\quad-\delta,\quad1+s,\quad23-17s,
 \quad2+4s,\quad(5-4s)/2.
\]
The certificate covers 7,004 primitive quadratics, their 14,008 roots,
and 111 rational candidates. To test whether $u+vs$ is a square, take
both signs of a rational square root $t$ of $u^2-2v^2$, test
$y_0^2=(u+t)/2$, then set $y_1=v/(2y_0)$ if $y_0\ne0$.
The $y_0=0$ case requires $v=0$ and $y_1^2=u/2$. Substitution verifies
each result. This is a necessary and sufficient test, including zero.
The last four listed coordinates are those of $R+T$ for $T\in C[2]$.
Every surviving nontorsion point therefore has height $\widehat h(R)$,
contradicting $\widehat h(R)=m^2\widehat h(S)$. This finite computation
is an explicit part of the proof, with a recorded independent replay;
it is not replaced by a rank oracle.

Thus $2C(K_0)=\Z(2R)$. The inverse coordinate change sends $2R$ to
$P_0$, as asserted.
\end{proof}

\subsection{Polynomial membership and addition}
Apart from the identity branch, membership in $2E(K_0)$ is expressed by
choosing $X,Y,H\in K_0$, $u\in\Q$, and imposing
\[
 Y^2=X^3-d^2X,\quad u\Norm(Y)=1,\quad2YH=3X^2-d^2,
 \quad x=H^2-2X,\quad y=H(X-x)-Y.
\]
Every nonidentity double has a preimage with nonzero ordinate.
Each field equation is two rational equations. Membership is an all-input
statement. Raw addition has the narrower exact contract
\[
 \operatorname{Add}_{\rm raw}(\iota(a),\iota(b),R)
             \quad\Longleftrightarrow\quad R=\iota(a+b).
\]
Its branches cover both identities, opposite points, distinct abscissas
and doubling with nonzero ordinate. Every later use supplies canonical
input codes. Multiplication first imposes membership of all three roots
and is sound on arbitrary rational tuples. In the natural predicate,
four initially untyped tuples enter multiplication calls, so are typed
before their squared values are added. Lagrange's theorem then gives
exactly $\iota(\N)$. This uses multiplication completeness; the generator
proposition alone does not supply the natural predicate.

% End 01-specification-model.tex


% Begin 02-height-pole-multiplication.tex
\section{The six-contact height estimate}\label{sec:height}
Set
\[
 \Cset=\{0,-15,-20,-24,-30,-60\},\qquad
 F(u)=u(u+15)(u+20)(u+24)(u+30)(u+60),
\]
and $G(u)=7u^2+300u+3600$. For a finite prime set $S$ containing all
primes at most 324 and $u\notin\Cset$, define
\[
 M_S(u)=\prod_{\substack{p\notin S:\ v_p(u-c)>0\text{ odd}\\
                          \text{for some }c\in\Cset}}p.
\]
Negative valuations of $u$ do not enter the product. Put
\begin{align*}
 F_0&=324!,\qquad D=F_0\,3^{256F_0^4},\qquad G_0=116064000000,\\
 E_0&=10(DG_0)^2,\qquad Q_R=R^{RE_0},\\
 B_{\rm cm}&=D^2\,3^{21000D^2}+1,\\
 C_{\rm cm}&=2B_{\rm cm}+4096(1028+5D+84000D^2),\\
 H_6(S;R)&=2\cdot10^{13}\max(1,C_{\rm cm},Q_R^{17}),\\
 &\hspace{1em}R\ge\max(D,\max S,10^{12}).
\end{align*}
When the second argument is omitted, we use
$H_6(S)=H_6(S;\max(D,\max S,10^{12}))$.
The proof follows the corrected geometric route in S09, with the
constants of S08. It does not use the older alternative modularity route.

\begin{theorem}[Six-contact estimate]\label{thm:height}
For every such $S,R$ and every $u\in\Q\setminus\Cset$,
\[
 h(u)\le H_6(S;R)M_S(u)^{24}.
\]
More precisely the external arithmetic-geometric inputs below give
$h(u)\le C(S)M_S(u)^{24}$, where
\[
 C(S)=10^{13}\max((8D)^{288D},6^{100})\Bigl(\prod_{p\in S}p\Bigr)^{24}.
\]
\end{theorem}
\begin{proof}
Take $B=(2,-3)$, the maximal order $\mathcal O=\Z[i,(1+j)/2]$, and
$\mu=2j+ij$. With its positive sign choice, the form
$\Trd(\mu^{-1}x\overline y)$ has matrix
\[
 \begin{pmatrix}0&0&-1&1\\0&0&1&-2\\1&-1&0&0\\-1&2&0&0\end{pmatrix}
\]
and determinant one. This is the principal-polarization datum whose
normalizer quotient is $\mathbf P^1_j$, with branch values
$0,\infty,-16/27$ and indices $4,6,2$ \cite{babagranath,linyang}.
The full $\mathcal O$-linear level-three cover $Y_3$ has total degree
at most $4|\mathcal O/3\mathcal O|=324$, retaining all determinant
components. Here the normalizer quotient contributes four. Central
ideles act trivially after multiplying by a rational scalar and choosing
its sign to give residue one modulo three.

For an integral model apply \cite[Theorem 2.1, Corollary 2.3]{buzzard}
first at $V_1(5)$. This level has full determinant image, is maximal at
the division primes, and uses $5\ge4$ prime to the discriminant six.
Add full five-marking and full $\mathcal O$-linear three-marking, then
forget five-marking by a free finite quotient. The three-marking sheaf
is nonempty because an $\mathcal O_3$-stable lattice in the rank-one
$B_3$-module is a principal fractional left ideal. Freeness follows
because a polarized automorphism has finite order, while a finite-order
element congruent to one modulo three on $T_3$ is the identity.
The universal abelian surface descends. The auxiliary level contributes
no factor to the final degree.

Every component has genus at least two, by its negative orbifold Euler
characteristic, so the automorphisms extend to its smooth proper model.
At $p>324$ the action is tame. A stabilizer is cyclic and its tangent
character is faithful; a nontrivial fixed-point equation has invertible
linear term. Unique lifting prevents branch collisions. The three
ordered branch sections fix the integral $j$-coordinate, so no further
denominator prime is discarded.

Pull back by $J(u)=-57600F(u)^2/G(u)^6$. If
\[
 H(u)=(u^2-60u-1800)(13u^2+480u+3600)(11u^2+600u+7200),
\]
direct calculation gives $G^6-97200F^2=H^2$. The pullback indices are
$2,6,2$. Tame normalization replaces $e$ by $e/\gcd(e,a)$ and leaves
$2,1,1$. The discriminants and resultants have prime support at most 13,
so this persists at $p>324$. The normalized cover of $\mathbf P^1_u$
has degree at most 324, not $12\cdot324$, and is branched only at the
six contacts. Locally it has equation $u-c=z^2$ after unramified
extension. Positive even contact valuation is unramified; negative
valuation specializes to the unbranched point at infinity. This is
exactly the support specified by $M_S(u)$.

The splitting field of a fiber has degree at most $F_0$. The lattice
of ordered homomorphisms among its conjugate abelian surfaces has rank
at most $16F_0^2$. A finite subgroup of $\operatorname{GL}_r(\Z)$
injects modulo three. Adjoining its field of definition gives a normal
field $L$ of degree at most $D$, ramified only in
$S\cup\operatorname{supp}M_S(u)$. It defines all endomorphisms and all
ordered conjugacy homomorphisms.

The geometric endomorphism algebra is either $B$ or $M_2(k)$ for an
imaginary quadratic field $k$, by taking the right centralizer of the
complex structure. In the first case Skolem--Noether corrects conjugacy
isomorphisms to $B$-linear isogenies over $L$. These are in particular
compatible $\Q(\sqrt2)$-linear isogenies. Thus the normal-field,
endomorphism and compatible-isogeny hypotheses, including condition $(*)$,
of \cite[Eq. (7.1), Section 7.1.4, Proposition 7.13]{vkk} hold.
With $P_S=\prod_{p\in S}p$ and $M=M_S(u)$, this gives
\[
 h_F(A)\le(8D)^{288D}(P_SM)^{24}
\]
in the non-CM case.

In the CM case retain the degree bound $d_L=[L:\Q]\le D$.
The maximal-order reduction of \cite[Section 7.3.3, Eq. (7.17)]{vkk}
constructs, over the same field, an isogenous product
$A'=\prod A_i^{e_i}$ with maximal CM endomorphism rings. All geometric
endomorphisms were already defined over $L$, so the reduction needs
no further field extension. Its intermediate CM-type estimate gives
$h_F(A_i)\le6^{96}(P_SM)^{10}$, since the CM fields are unramified
outside the discriminant support of $L$. As $\sum e_i\le2$, put
$H_0=2\cdot6^{96}(P_SM)^{10}$; then $h_F(A')\le H_0$.
The isogeny estimate \cite[Eq. (7.13)]{vkk}, in dimension two, is
\[
 h_F(A)\le h_F(A')+\tfrac12\log\kappa(A'),\qquad
 \kappa(A')=\bigl(14^{32}d_L
             \max(h_F(A'),\log d_L,1)\bigr)^{16}.
\]
Since $\log14<3$, $D\ge2$, and $H_0\ge1$, it follows that
\[
 h_F(A)\le9H_0+16D+768\le(8D)^{288D}(P_SM)^{24}.
\]
For the last inequality use $18\cdot6^{96}<16^{64}$ and
$16D+768\le(8D)^4$, then absorb their sum in $(8D)^{65}$.
This argument bounds the original stable height while retaining the
field degree. It does not use the degree-free final conclusion of
Proposition 7.14, nor the historical alternative CM bound in S08.
The prescribed integer $C_{\rm cm}$ remains unchanged. Its former
degree-free interpretation is unnecessary: the corrected estimate is
absorbed by $Q_R^{17}$. Later comparisons between those integers are
numerical comparisons, not uses of the discarded height claim.

The compatible theta normalizations in \cite{linyang,lny} give
\[
 j=3779136\frac{(\prod z_m)^4}{(\sum z_m^8)^5},\qquad
 h(j)\le40h([z_m])+22,
\]
with ten even theta nulls; the other six vanish. The comparison in
\cite[Corollary 1.3(2)]{pazuki} yields
$h([z_m])\le10000\cdot2^{24}\max(1,h_F(A))$.
The change to the Deligne metric adds $\log\pi$ in dimension two and
has been included in this constant. A finite extension for a theta
structure changes neither absolute height. The coefficient norms of
$-57600F^2,G^6$ are below $4000^6$. The rational-root bound gives
$h(u)<h(j)+74$, after removing any initial power of $X$ or dropped
leading term. Since $40\cdot10000\cdot2^{24}+96<10^{13}$, this proves
the bound with $C(S)$.

Finally $P_S\le R^R$, $8D\le R^2$, and $6^{100}\le R^{576D}$.
Thus
\[
 C(S)\le10^{13}R^{576D+24R}\le10^{13}R^{600R}
 \le10^{13}Q_R^{17}\le H_6(S;R)/2.
\]
This also proves the assertion for every permitted larger $R$.
\end{proof}

The final value argument uses $S_{\rm ht}=\{p:p\le324\}$. Multiplication
uses $S_{\rm mul}=S_{\rm ht}\cup\mathcal H\cup\operatorname{supp}(c_*)$,
where $\mathcal H$ is the fixed catalogue. A common numerical $R$ may
bound both sets. It does not identify them.

\section{The pole predicate and its converse input}\label{sec:pole}
Let $S_0=\{2,3,5,7,13\}$. With input $B$ and eleven witnesses
$b_0,\ldots,b_5,e,z,r_1,r_2,w$, define
\begin{align*}
 m&=(1-2b_0)(1+b_1)(1+2b_2)(1+4b_3)(1+6b_4)(1+12b_5),\\
 A&=mB,\qquad H=A(A-B^2),\\
 T_1&=A(A-e)(A+3e),\qquad T_2=A(Az^2-eH)(Az^2+3eH),\\
 Q_{\rm pole}&=B\Bigl(\sum_{i=0}^5(b_i^2-b_i)^2+(eHzw-1)^2\\
 &\hspace{4em}+(r_1^2-T_1)^2+(r_2^2-T_2)^2\Bigr).
\end{align*}
Write $\Phi(B)$ for rational solubility. For $B\ne0$, $m$ is a signed
divisor of 2730, and the equations give points $(A,r_1)\in E_e$ and
$(Az^2,zr_2)\in E_{eH}$ with $eHz\ne0$, where
$E_\ell:y^2=x(x-\ell)(x+3\ell)$. The polynomial has total degree 75.

\begin{lemma}[One-sided obstruction]\label{lem:pole}
If $B\ne0$ and $\Phi(B)$, then $v_p(B)$ is not negative and odd for
any $p\notin S_0$.
\end{lemma}
\begin{proof}
If $k=v_p(B)<0$ were odd, then $v_p(A)=k$, $v_p(H)=3k$.
Exactly one of $e,eH$ has even valuation. Square scaling makes its
parameter a unit, preserving the parity of the abscissa's valuation.
For a unit parameter the cubic has valuation $v_p(x)$ if $v_p(x)>0$
and $3v_p(x)$ if $v_p(x)<0$. Its being a square forces even valuation.
But both points have abscissa in the squareclass $A$, a contradiction.
\end{proof}

\begin{theorem}[Pole completeness, $\mathsf{PC}$]\label{hyp:PC}
For every rational $B$ with square reduced denominator,
$0<B<1/2$ and $v_p(B)\ge4$ for every $p\in S_0$ imply $\Phi(B)$.
\end{theorem}
\begin{theorem}[Restricted converse, $\mathsf{RC}$]\label{hyp:RC}
For every positive squarefree $\ell\equiv11\pmod{12}$,
\[
 \operatorname{corank}_{\Z_2}\Sel_{2^\infty}(E_\ell/\Q)=1
 \quad\Longrightarrow\quad\ord_{s=1}L(E_\ell,s)=1.
\]
\end{theorem}
Here $c_4=208\ell^2$, $\Delta=2304\ell^6$, and $j=35152/9$.
The nonintegral rational $j$ shows these curves are non-CM. They are
different from the CM generator curve. Theorem~\ref{hyp:RC} is proved in
Sections~\ref{sec:rc}--\ref{sec:limits}, with its coefficient and
normalization arguments in Appendices~\ref{app:even-source}--\ref{app:endpoints}.
We now prove Theorem~\ref{hyp:PC} from this restricted converse.

\subsection{The exact descent reduction}
The signed divisors of 2730 realize all real, $2$-adic and $3$-adic
squareclass signatures. After $-1,2,3$ fix sign and valuation parities,
the remaining unit vectors of $5,7,13$ are
$(0,1,1),(1,0,0),(0,1,0)$ and form a basis. Choose positive $m$ making
$A=mB$ square at two and a nonsquare unit squareclass at three.
Let $\beta,\theta$ be the positive squarefree representatives of $A,H$.
Then $\beta\equiv17\pmod{24}$, $\theta\equiv1\pmod{24}$, with disjoint
supports, $\beta>1$, and $\theta=1$ permitted. Each is a local square
at the support of the other. At a prime of $\beta$, the relevant odd
valuation of $B$ is positive, so $1-B/m$ is a square unit. At a prime
of $\theta$, $v(B)$ must be zero, since otherwise $v(H)$ is even;
there $B/m\equiv1$, so $A$ is a square.

The simultaneous construction in S11--S12 gives $n,n\theta$ positive
squarefree and congruent to eleven modulo twelve, both with two-Selmer
dimension three and with extra Kummer class $(\beta,1)$. Its matrix
conditions can be stated explicitly. For support primes $q_i$, set
$s_i=[-1/q_i]$, $\lambda_i=[3/q_i]$, and
\[
 (Kv)_i=\sum_{j\ne i}[q_j/q_i](v_j+v_i).
\]
One arranges $\ker K=\langle\mathbf1,\mathbf1_{\operatorname{supp}\beta}\rangle$
and $\lambda\notin\operatorname{im}K$ on both supports. Add $d-1$
prime slots of type $(s,\lambda)=(0,0)$ and one slot $q_0$ of type
$(1,1)$, where $d=\omega(\beta)$. The cross-block has columns
$e_i+e_d$. After the rank-one correction $K^\sharp=K+\lambda a^t$,
deleting $q_0$ leaves
\[
 \begin{pmatrix}D_1&C_1\\C_1^t&0\end{pmatrix},
 \qquad\mathbf1^tD_1\mathbf1=1.
\]
This is nonsingular: $C_1^tx=0$ makes $x$ constant, and the first
equation forces that constant and then $y$ to vanish. Adding the
support of $\theta$ preserves this block. Free $q_0$-symbols toggle
the new diagonal entries; the determinant's full-product coefficient
is one. Thus a single Boolean assignment works on both supports.
Reciprocity, CRT and Dirichlet realize it by distinct auxiliary primes.

In Kummer coordinates $(x+3\ell,x-\ell)$, a positive-second-coordinate
Selmer representative is $(3^\delta\prod q_i^{u_i},\prod q_i^{v_i})$.
The necessary local equations are
\[
 Ku+\operatorname{diag}(\lambda)v=\delta\lambda,\quad
 (K+\operatorname{diag}(s))v=0,\quad(\lambda+s)^tv=0.
\]
The final equation is the condition at three. The matrix properties
force $v=0$, $\delta=0$, $u\in\ker K$, so there are at most eight
classes in total. The extra class is actually local at every place,
including two, and independent of the two torsion classes, giving
dimension three. This sufficiency check for the particular class is
essential; arbitrary parity conditions at two would not suffice.

\begin{lemma}[Dimension three gives corank one]\label{lem:corank}
For an elliptic curve with full rational two-torsion,
$\dim_{\F_2}\Sel_2=3$ implies $\operatorname{corank}_{\Z_2}\Sel_{2^\infty}=1$.
\end{lemma}
\begin{proof}
The finite Kummer sequence gives $r+\dim\Sha[2]=1$.
Write $\Sha[2^\infty]=(\Q_2/\Z_2)^t\oplus F$ with $F$ finite.
The Cassels--Tate pairing on the quotient by the maximal divisible
subgroup is nondegenerate and alternating for an elliptic curve
\cite{poonenstoll}. Its finite primary cyclic factors are paired,
so $\dim F[2]$ is even. Thus $F=0$ and $r+t=1$. The restriction of
the pairing to $F[2]$ need not be nondegenerate and is not used.
\end{proof}

Applying Theorem~\ref{hyp:RC} to both $n$ and $n\theta$, the forward rank-one
theorem \cite{gz,kolyvagin} gives rank one and finite $\Sha$, hence
$\Sha[2]=0$.
The Kummer class is then represented by actual rational points. Their
abscissas have squareclass $\beta$, by the cubic equation. Write
$H=\theta c^2$ and rescale the second point to $E_{nH}$. If the first
abscissa is $Aq^2$, divide both parameters by $q^2$ and both coordinates
by $q^2,q^3$. The abscissas become $A,Az^2$, giving precisely the
eleven witnesses, with $w=(eHz)^{-1}$. If $\theta=1$, one can reuse
a point on the same curve. This proves Theorem~\ref{hyp:PC}, with both
restricted-converse applications in their stated domain.

\section{Multiplication and the finite false-triple guard}\label{sec:mult}
The actual multiplication template first forces all three input codes
to be members of $2E(K_0)$. These membership calls are conjuncts inside
the height certificate; they are not assumed caller preconditions.
Thus even untyped rational inputs, if accepted, have unique indices.
The template has 35,134 rational witnesses. Its construction is S03
and its mathematical specification is S13.

For indices $n_1,n_2,n_3$, each test group chooses $H=[h]P_0$,
$U_i=[z_i]P_0$, and sets $V_0=H-[c_*]P_0$,
$V_i=U_i-[c_*n_i]P_0$, where $c_*=4M_0$ is fixed. Use
\begin{align*}
 L(z)={}&16575z-6630d^2z^5+11050d^4z^9-25500d^6z^{13}\\
        &+68250d^8z^{17}-198900d^{10}z^{21}+612612d^{12}z^{25}.
\end{align*}
This is 16575 times the degree-25 truncation of the normalized formal
logarithm. The twelve tests are four inverse abscissas at $H,U_i$,
their four conjugate counterparts at $V_0,V_i$, three differences of
the corresponding $L$-ratios, and the product of the first two unprimed
ratios minus the third. Identity inverse abscissas are zero; all
denominators are guarded.

For a slope $t\notin\Cset\cup\{1\}$, put
\[
 R_0=\eta F(t)/(t-1)^8,\quad\eta^2=1,\quad
 a=1-R_0q^2\ne0,\quad R=\ell_0^2a^2R_0.
\]
Denote the twelve unscaled test values in group $g$ by $D_{gj}$.
Their normalized values are
\[
 C_{gj}=\frac{D_{gj}}{R_0^{k_j}(t-1)^{k_j+1}},\qquad
 k_j=\begin{cases}1,&1\le j\le8,\\25,&9\le j\le12.\end{cases}
\]
For six independent groups set
$W_g=\prod_j(1+RC_{gj}^2)$ and
$\Theta=\Tr\sum_gW_g^{-1}$, require them nonzero, and impose
$\Phi(R/\Theta^2)$. The pole input is nonzero, so the prefactor $B$
in $Q_{\rm pole}$ cannot make these constraints vacuous.

At a positive odd contact prime outside $S_{\rm mul}$, each factor
valuation is the minimum of zero and an odd integer. No cancellation
occurs. The trace and Lemma~\ref{lem:pole} force one common group and
place where the first eight tests have valuation at least one and
the last four at least 25. The first tests put the points in the
formal group. At least one of $h,h-c_*$ is a unit at the prime, so its full-log
ratio is integral. The truncated logarithm is the full logarithm times
a unit congruent to one in valuation at least 28. Its ratio is therefore
integral. The difference test transfers integrality to the other truncated
ratio, then to the other full-log ratio. Only now do the relative errors
give absolute congruences modulo $p^{28}$. Multiplying the two ratio congruences
by their integer denominators gives
$n_i\equiv L(z(U_i))/L(z(H))\pmod{p^{25}}$.
The product test then implies $p^{25}\mid n_1n_2-n_3$.

To specify the height certificate, write the abscissa of each input as
$x_{i0}+x_{i1}s$, using $x_{i0}=x_{i1}=0$ at the identity. Choose
$a_{i0},a_{i1},b_i\in\Q$ with $b_i\ne0$ and
$b_i x_{ij}=a_{ij}$. Express each of these nine rational quantities
as a sum of three slopes. Impose the preceding slope predicate on
all 27 slopes, with private witnesses, and impose membership on the
three input codes. This conjunction is the height certificate.
The multiplication predicate conjoins it with the fixed scalar
graphs and the inverse of the finite guard below.

\begin{proposition}[Finite-box deduction]\label{prop:box}
Assuming these local test identities and the height estimate for
$S_{\rm mul}$ with constant $A_{\rm ht}$, every false triple accepted
by the unguarded certificate has $\max_i|n_i|<T=(48A_{\rm ht})^{13}$.
\end{proposition}
\begin{proof}
Put $f=n_1n_2-n_3\ne0$. For each of the 27 slopes,
$M_{S_{\rm mul}}(t)^{25}\mid f$, so $h(t)\le A_{\rm ht}|f|^{24/25}$.
The nine rational homogeneous quantities for the three abscissas
are sums of three slopes. The sum and quotient height bounds give
$h(x_j)\le6A_{\rm ht}|f|^{24/25}+2\log3$ for each rational component.
The fixed coordinate change and the duplication bounds yield
\[
 2n_i^2\widehat h(P_0)\le12A_{\rm ht}|f|^{24/25}+B_0,
\]
where $2\widehat h(P_0)>1$ and
$B_0=4\log3+\tfrac32\log2+2\log(1+\sqrt2)+\tfrac13\log676<10$.
For $N=\max|n_i|$, $|f|\le2N^2$, so
$N^2<24A_{\rm ht}N^{48/25}+10$.
If $N\ge5$, this implies $N^{2/25}<48A_{\rm ht}$, hence
$N<(48A_{\rm ht})^{25/2}\le T$. The case $N\le4$ is immediate.
\end{proof}

For each false triple $(a,b,c)\in[-T,T]^3$, write its fifteen code
coordinates as $A_j/B_j$ with $B_j>0$, put $D_{abc}=\prod B_j$, and
form
\[
 g_{abc}(X)=\sum_{j=1}^{15}
       (D_{abc}X_j-(D_{abc}/B_j)A_j)^2.
\]
The guard requires the product of all these factors to have an inverse.
A false accepted triple makes its own factor zero. At a true triple
every factor is positive by injectivity of codes, and the inverse exists.
Thus the sound direction uses the generator and corrected all-$S$
height theorem but does not use pole completeness.

For completeness, S14 constructs, for every integer $A$ and bound $B_0$,
integers $u_1,u_2,v$ with $u_3=Av-u_1-u_2$ such that all 22 forms
$v,u_j-bv$, $b\in\Cset\cup\{1\}$, are signed catalogue multiples of
primes above $B_0$. At each such prime the conjugate reduction orders
have gcd four and are prime to that prime. The finite-complexity prime
asymptotic \cite{greentao,moebius,gtz} is used after fixing the catalogue
and initial CRT data, then the sieve cutoff $L$. Twenty-two catalogue
primes per unit class allow distinct choices for the twenty-two
occurrences. At one catalogue prime the distinguished zero plane avoids
at most $42p<p^2$ points; a lift modulo $p^2$ makes its divided value
different from zero and minus one. The fixed affine system is therefore
locally admissible. The asymptotic gives a lower count
$c_A(\log L)^{-44}X^3/(\log X)^{22}$ for fixed sieve cutoff $L$.
Each pair of forms has a two-by-two coefficient minor equal to
$\pm1$ or a nonzero contact difference of absolute value at most 61.
Thus the local conditions are nondegenerate uniformly in $A$;
other minors may involve $A$.
A Gaussian-pair sieve and the 21-form sieve on each nonzero fiber bound
the loss by $C_AX^3/(L(\log X)^{22})+o_A(X^3/(\log X)^{22})$.
Choose $L$ first and then $X$ with $A,L$ fixed. A uniform prime-asymptotic
error as $L$ varies is neither asserted nor needed. The corrected
all-orders inverse theorem is among the external inputs.

Apply this construction to the nine integral homogeneous quantities
of a true triple. Each slope has at most six odd-contact primes.
At one such prime with reduction orders $d_1,d_2$, solve
\[
 h\equiv0\pmod{\operatorname{lcm}(M_0,d_1)},\qquad
 h\equiv c_*\pmod{d_2}.
\]
Compatibility follows from
\[
 \gcd(\operatorname{lcm}(M_0,d_1),d_2)
                  \mid M_0\gcd(d_1,d_2)=4M_0=c_*.
\]
Choose a solution larger than $c_*$. Each contact prime uses its own
private group within its own slope. No common shift across different
primes is imposed, and selected prime values need not be distinct.
The saved library allocates six private groups per slope and twenty-seven
private slope calls. Set $U_i=[n_i]H$ and assign one group to each prime. Fixed local
clearing at the old primes can be checked as follows.
Put $S_b=S_{\rm ht}\cup\mathcal H$, $B_p=4$ on $S_{\rm ht}$ and
$B_p=1$ on the remaining catalogue primes, $C_p=14B_p$, and
$D_p=51C_p$. The selected integral forms give
$|v_p(R_0)|,|v_p(t-1)|\le C_p$, while the small formal subgroup
in the definition of $M_0$ gives $v_p(D_{gj})\ge0$ at $S_b$.
Consequently $v_p(C_{gj})\ge-D_p$. The fixed normalization has
$v_p(\ell_0)=E_p=52C_p+v_p(12)+4$.

Weak approximation chooses $q$ small at $S_b$, integral at the
selected primes, and arbitrarily close to a real root of
$1-R_0q^2$, without equality; choose $\eta$ to make $R_0>0$.
Then $a$ is a unit at the specified finite primes. At $p\in S_b$,
\[
 v_p(RC_{gj}^2)\ge 2E_p-C_p-2D_p>v_p(12),\qquad
 v_p(\Theta-12)>v_p(12),
\]
so $v_p(R/\Theta^2)\ge103C_p+8\ge4$.
At a selected odd-contact prime, $v_p(R)=1$. Each of the twelve
summands of the trace of the six $W_g^{-1}$ has residue zero or one;
the assigned group and place contributes one. Since $p>324$, their
sum is a unit and the pole input has valuation one. At every other
prime outside $S_b$, its valuation is even. Thus its reduced
denominator is square. In both real embeddings each $W_g$ is positive;
as $a\to0$, $W_g\to1$, $\Theta\to12$, and $R/\Theta^2\to0$
through positive values. This verifies exactly $\mathsf{PC}$'s
hypotheses, including all nonzero guards. Negative indices, zero indices and
coincident points are handled by the group-law branches. Therefore
Theorem~\ref{hyp:PC} supplies the required pole witnesses. This proves
the full multiplication statement $\mathsf M$, including completeness
for every integer triple $(n_1,n_2,n_1n_2)$ and soundness on arbitrary
rational input codes.

% End 02-height-pole-multiplication.tex


% Begin 03-restricted-converse.tex
\section{The restricted converse and its dependency order}\label{sec:rc}
For a positive squarefree $\ell\equiv11\pmod {12}$ put
\[
 E_\ell:y^2=x(x-\ell)(x+3\ell).
\]
Its invariants are $c_4=208\ell^2$, $\Delta=2304\ell^6$ and
$j=35152/9$. Thus it has full rational two-torsion and is not a
CM curve, since a rational CM $j$-invariant is an integer.
The curve in this section is different from the fixed CM curve used
to represent integer indices.

The assertion needed in Theorem~\ref{hyp:RC} is
\[
 \operatorname{corank}_{\Z_2}\Sel_{2^\infty}(E_\ell/\Q)=1
 \quad\Longrightarrow\quad \ord_{s=1}L(E_\ell,s)=1.                 \tag{RC}
\]
We prove this assertion in this section and the next two sections.
Appendices~\ref{app:even-source}--\ref{app:endpoints} give the fixed-form,
theta, nonvanishing and normalization constructions used in this deduction.
Fix $E=E_\ell$ and a modular parametrization
$\pi:X_0(N)\to E$ with $\pi(\infty)=0$. Fix the integral Tate lattice
and the local squareclasses allowed at the primes of $2N$.
Constants depending on this data are chosen before any conductor,
companion field or binary dimension varies. Quotienting by a rational
point of order two gives a rational two-isogeny. The Selmer parity
theorem \cite[Theorem 1.8]{dd} therefore gives root number $-1$
under the left side of (RC). This application does not replace Selmer
corank by Mordell--Weil rank and does not assume finite $\Sha$.

Here is the dependency order of the proof. A fixed-form Euler-system
calculation gives the even coefficient bound. A varying-companion
arithmetic estimate, proved by the two limits below, then gives the
genus-point lower bound. Only after that lower bound do we normalize
the odd coefficients. Their finite-stage forms and the even forms give
one common binary family. Finally a fixed-field determinant argument
excludes a zero derivative at the base. The even estimate does not use
that final determinant argument on Heegner classes.

\subsection{The even coefficient calculation}
Use normalized central values $L_*(h)$ with the real and imaginary
periods fixed by $\pi$ and the integral lattice. Only their valuations
are used; every subsequent comparison uses this same normalization.
For every twist whose two-primary Tate--Shafarevich group is finite, let
$s(h)=\operatorname{length}_{\Z_2}\Sha(E^{(h)}/\Q)[2^\infty]$.

The fixed-form construction uses one integral symbol combination with
nonzero projections to both real signs. Kato's integral maps, with
$j_*$ cohomology taken modulo each power of two before passage to the
limit, give classes whose inertia restrictions vanish integrally.
Indeed the low-degree Leray edge map gives this vanishing before
rationalization. Restriction of transfer is the Mackey sum of the
corresponding upper inertia restrictions. Hence unnormalized transfer
preserves the vanishing even when the intervening extension is ramified.
No character projector divided by a growing group order is used.

For a finite binary character group, let $F$ be the least common
multiple of the odd parts of all its character conductors, including
the base character. Write $Z_F$ for the compatible Iwasawa class
obtained by pushing the retained-smoothed Kato classes at levels
$F2^n$ through the integral coefficient map.
At a character $\chi$ its unused-prime factors are
\[
 E_q(\chi)=1-\frac{a_q\chi(q)}q\sigma_q^{-1}
                 +\frac1q\sigma_q^{-2}.                           \tag{3.1}
\]
They are the actual determinants of the unused good local terms.
Active ramified good primes have rational inertia $-1$, so have no such
rational local term. The fixed bad and dyadic corrections form a factor
$F_E$ with nonzero central value.

The rational divisibility input is Kato's Theorem 13.4(2),
together with the nonzero-class input in 13.5(2)
\cite{kato}. We use the height-one primes not containing two.
Rational irreducibility is preserved on a quadratic subgroup of the
non-CM open image. That subgroup contains a nontrivial unipotent,
which supplies the required one-dimensional fixed space. Oddness
supplies the two real eigenspaces, and purity supplies the common
Frobenius weight. This application does not use residual irreducibility
at two or the stronger hypotheses of Theorem 13.4(3).

The smoothing constants $c,d_s$ can be chosen to satisfy
\[
 c,d_s\equiv1\pmod {3N_{\rm odd}F},\qquad
 c,d_s\equiv1+2^e\pmod {2^{e+1}},
\]
where $e\ge\max\{3,v_2(N)\}$ is fixed by the original dyadic data.
In weight two and
cohomological degree one the two smoothing exponents are one.
Kato's explicit reciprocity formula therefore contributes valuation
$2e$, independent of $F$. Quadratic Gauss and square-root generators at
odd conductor differ by dyadic units. The numerical sizes of $c,d_s$
may grow; their relevant valuation does not.

For completeness, the integral calculation at two uses the positive
complex. Finite continuous cochains are filtered unions of finite free
modules and hence flat. Splitting the residual complex into cohomology
and disks, then using the nilpotent perturbation calculation of
Section~\ref{sec:limits}, gives finite free models with their maps.
In the fixed real cyclotomic tower the ordinary two-primary class group
vanishes. One proves this inductively: only the prime above two ramifies
in each quadratic step; the abelianization of a possible Hilbert
two-class tower is generated by that inertia of order two. A finite
two-group with cyclic abelianization is cyclic, and inertia excludes a
larger cyclic extension. The unit-signature calculation is separate.
A norm-minus-one unit in positive layers, and $-1$ at the base, has
signature vector a unit in the characteristic-two cyclic group algebra.
For a scalar coefficient the resulting residual dimensions are
\[
 \dim H^1=2^n+s_n-1,\qquad \dim H^2=s_n-1.
\]
For each fixed support $s_n$ is bounded. Universal coefficients give
generic residual rank one in degree one; the two trivial elliptic
factors give rank two. Unit cancellations thus identify the positive
complex over $\mathscr O[G]$ with $\mathscr O[G]^2[-1]$.
This last rank is independent of support, although preliminary models
need not have a uniform size.

A generator of the fixed real invariant lattice gives a connecting class $a$.
Its central image need not be saturated.
Twice an ordinary class lifts because the relevant real $H^1$ is killed
by two. Wedge $a$ with a lift of $2F_E Z_F$ in the original determinant
line. This produces a common scalar $u^{\rm ev}$.
The first assertion below holds for the full family. The central valuation
identity is asserted only at a constituent of Selmer corank zero with
$L_*(h)\ne0$:
\[
 2^{C_{\rm ev}}u^{\rm ev}\in\Lambda[G],\qquad
 v_2u^{\rm ev}_h(0)=v_2L_*(h)-2\omega(h)-s(h)+O_E(1).               \tag{3.2}
\]
Here $\Lambda=\Z_2[[t]]$ and
$\mathscr O=\widehat{\Lambda_{(2)}}$. The error and $C_{\rm ev}$ are
fixed before the binary group and support. To justify membership in
$\Lambda[G]$, first use rational characterwise divisibility and Fourier
inversion for membership in $\Lambda[1/2][G]$, without estimating the
Fourier denominator. Intersect coefficientwise with the integral
$\mathscr O[G]$ result. The explicit coefficient construction also covers field discriminant $4h$,
with its dyadic character retained before the cofinal norm limit, as in
Appendix~\ref{app:even-source}. At a corank-zero constituent the same coordinate
vanishes centrally if and only if the central value vanishes.
At such a nonzero corank-zero constituent, keeping the real saturation
and torsion terms gives the local formula
\[
 v_2u(0)=v_2\exp^*z-s-\sum_v g_v+2g+1-r_1-e_2.
\]
At an active good odd prime $g_v=2$, since $A^{I_v}=E[2]$.
Unused terms cancel (3.1); all other terms belong to the fixed local
list. This proves (3.2) with the indicated active-prime coefficient.

\subsection{Nonvanishing with all support restrictions}\label{sec:nonvan}
A signed unit filter prescribes a real sign and unit squareclasses at
a finite set $S$ containing $2,3$ and the bad primes. Let $B$ be any
additional finite set of forbidden primes. If its root number is even,
there is a squarefree parameter in this filter, coprime to $B$, with
at most five prime factors and nonzero central value. In an odd filter
there is a simple-zero parameter, with no asserted prime-count bound.

Here is the reduction to the precise classical statements.
Choose a signed prime $q_0$ in the prescribed unit classes by
Dirichlet's theorem, avoiding $S\cup B$. Apply Hoffstein--Luo
\cite{hl} in the even case, or Bump--Friedberg--Hoffstein
\cite{bfh} in the odd case, to the actual curve $E^{(q_0)}$.
Require their positive discriminant $d$ to split at
$S\cup B\cup\{|q_0|\}$ and at the conductor primes. Then $q_0d$ has
the prescribed original classes. Hoffstein--Luo gives at most four
prime factors of $d$, so at most five in total.
The BFH clause used in the odd case is the simple-zero clause.

The root-number check includes ramified dyadic quadratic characters.
At the primes in $S$ equal squareclasses give isomorphic local curves.
At a good odd prime newly ramified in the parameter, the root factor is
$(-1/p)$ \cite[Proposition 3.3]{desjardins}. Its product outside $S$
is fixed by the real sign and the dyadic unit class. Thus unit classes
$3$ or $7$ modulo eight are allowed; the field discriminant may be four
times the parameter. We do not invoke an odd-discriminant root formula
on those parameters.

\paragraph{The extra upper-bound argument.}
Nonvanishing alone does not bound the valuation of $L_*$, and (3.2)
by itself gives a lower bound. The needed upper bound uses a separate
even-symbol selection and transfer argument, which we now give.

For the even coefficient system described below, a unit symbol at a
parameter $d$ gives
\[
       L_*(d)\ne0,\qquad v_2L_*(d)\le2\omega(d)+C_0.              \tag{3.3a}
\]
This follows from its same-filter coefficient-square formula and its
fixed normalization. That even normalization uses the lower bound
from (3.2), not the as-yet undefined odd minimum.
Call a prime list partitionable if a permitted contraction forest
takes its symbol to a unit seed. The indicated iterated finite
difference is nonzero, so its terminal polynomial is nonzero.

For a partitionable parameter $d$ already known to have Selmer
corank zero, keep all its prime factors stationary. The one-symbol
address construction gives arbitrarily large binary cubes based at
$d$, with no old prime removed, whose other addresses satisfy (3.3a).
The fixed-form common scalar in (3.2) has one uniformly cleared
denominator on every such cube. Choose a two-adic precision larger
than the nonzero-address scalar valuation bound and its fixed
clearing exponent. Its sign-evaluation coefficient is a Boolean
function with the bounded bit degrees proved in
Section~\ref{sec:limits}. A cube dimension larger than their sum
therefore supplies a nonzero address congruent to the base at this
precision. If the base scalar were zero or had larger valuation than
the bound, so would that nonzero address, a contradiction.
The corank-zero central comparison in (3.2) consequently gives
\[
 L_*(d)\ne0,\qquad
             v_2L_*(d)-2\omega(d)-s(d)\le C_1.                  \tag{3.3b}
\]
The constants are fixed by the even system. The dimension needed
may depend on $d$; the valuation bound does not. This is the
cyclotomic missing-address argument on the fixed-form Kato scalar,
not the later Heegner transfer.

Here is why the actual rough companion is partitionable.
Prepare an even symbol for every possible total type at the finite
fixed support. There are finitely many. In our full rational
two-torsion case the alphabet is all split. A least nonempty unit
seed has no dependence on internal edge bits: an edge derivative
would be a smaller nonempty unit configuration. Its individual
labels determine it. The negative-discriminant parity at two makes
the total label nonzero, so this seed is nonempty.
A list with the same total label and enough copies of the seed
labels contracts to it: keep the first $r-1$ seed anchors and merge
all remaining vertices into one nonempty cluster. Its label is
forced to be the last seed label. Thus the list is partitionable,
regardless of the number or labels of the remaining vertices.

Before choosing the rough factor, insert bounded dummy supplies
for every prepared total type, including the required anchors and
reserves. Realize their labels by distinct primes. For a varying $h$,
place in $B$ the entire fixed support, the support of $h$, every
dummy prime and every other already fixed forbidden prime. If a
reduction prime has already been fixed, include it too.
Prescribe the total $k'$ to be a negative odd fundamental
discriminant, square at every place of $2Nh$ and outside the
extra-unit cases. Absorb $h$ and the dummy product into the curve
and use the preceding nonvanishing theorem in its even filter.
It adds at most five new primes. For the actual resulting total
type, the already chosen dummy supply makes $hk'$ partitionable.
Its analytic nonvanishing gives Selmer corank zero by the forward
rank-zero theorem. Apply (3.3b), obtaining
\[
 \omega(k')\le A,\qquad
 v_2L_*(hk')\le 2\omega(h)+2\omega(k')+s(hk')+C_E.                \tag{3.3}
\]
Here $A$ is the fixed dummy count plus five. This is the argument
of the partitionable-even and rough-companion lemmas, with their
stationary primes and order of choices retained.
Neither prime sizes nor a uniform analytic threshold are asserted.
Excluding a later coefficient prime does not enlarge the old
Selmer support.

\subsection{The genus lower bound precedes normalization}
Fix a negative genus partner $h_*$ with $L_*(h_*)\ne0$.
For a known simple-zero $h$, let $P_h^{\rm gen}$ be the unaveraged
conductor-one genus sum over $\Q(\sqrt{hh_*})$.
Let $P_{h,k'}^{\rm rc}$ be the unaveraged conductor-$|h|$ ring-class
sum over $\Q(\sqrt{k'})$, weighted by $\chi_h$.
They are different points in different native lattices.
Write $j_{\rm gen}(h)$ and $j_{\rm rc}(h,k')$ for their index valuations.

The explicit height formula \cite[Theorem 1.1]{cst} applies twice.
All level primes split, conductors are prime to $N$, extra-unit fields
are excluded, and the parametrization is the same. Its conductor factor
is $c\sqrt{|\operatorname{disc}K|}$ and its sum is unaveraged.
The two Rankin products have the common factor $L(E^{(h)},s)$.
Each rational free line has index at most two in its native rank-one
lattice: conjugation fixes the free class, and the norm makes twice the
class rational. Thus
\[
 2j_{\rm rc}(h,k')-2j_{\rm gen}(h)
   =v_2L_*(hk')-v_2L_*(h_*)+\epsilon,\qquad |\epsilon|\le2,        \tag{3.4}
\]
up to the one fixed initial period convention.

The arithmetic estimate of Section~\ref{sec:limits}, for each fixed $A$,
is
\[
 2j_{\rm rc}(h,k')\ge
 4\omega(h)+s(h)+s(hk')-C_R(E,A),\qquad \omega(k')\le A.           \tag{3.5}
\]
Combining (3.3)--(3.5) cancels $s(hk')$ and gives
\[
 2j_{\rm gen}(h)\ge 2\omega(h)+s(h)-C.                            \tag{3.6}
\]
This calculation uses no odd normalization.

\subsection{A fixed nonzero witness and actual odd forms}
Choose one positive squarefree $D_w$ in the odd filter containing one,
after adjoining the primes of $h_*$ and their compatible unit tests.
Section~\ref{sec:nonvan} supplies $L'(E^{(D_w)},1)\ne0$.
The explicit height formula makes $P_{D_w}^{\rm gen}$ nontorsion.
Its unwanted $h_*$ component has rank zero. One fixed multiplier $L$
kills the bounded-degree torsion and the projection denominator, so
$Q=LP_{D_w}^{\rm gen}$ is a rational nontorsion point on $E^{(D_w)}$.

Let $C$ be complex conjugation on its Tate module and choose a nonzero
integral row $b$ with $bC=b$. The Kummer translation image of $Q$ is a
full lattice. Indeed, if it vanished, its rational cocycle would descend
to the Tate image. A nonidentity central scalar $z$ gives
\[
 (z-1)c(g)=(g-1)c(z),
\]
so that descended cocycle is a coboundary, contrary to Kummer
injectivity. Irreducibility makes a nonzero translation span the Tate
space, and intersection with any fixed finite kernel preserves full
rank. Twisting preserves open image by restriction to an index-at-most-two
subgroup.

Choose one such translation $v$ with $2b(v)\ne0$.
Multiplying conjugation by this translation gives square translation
$(1+C)v$, detected by $2b(v)$ with one fixed valuation $c_w$.
The element is a single actual Galois element on the compositum.
It is trivial on $\sqrt{D_w}$ and nontrivial on $\sqrt{h_*}$ because
$D_w>0>h_*$. Its restrictions therefore already agree on intersections.
Finite Chebotarev, with a place chosen above the resulting prime, gives
arbitrarily advanced stages $(M_\theta,\ell_0)$ such that
\[
 \rho(\operatorname{Fr}_{\ell_0})\equiv C\pmod {2^{M_\theta}},
 \quad (D_w/\ell_0)=1,\quad(h_*/\ell_0)=-1,
 \quad v_2\lambda_{M_\theta}(\overline Q)=c_w<M_\theta.            \tag{3.7}
\]
Here $\lambda_{M_\theta}=b\kappa_{M_\theta}$ is the actual finite-field
Kummer map on $E(\F_{\ell_0^2})$. It has Frobenius sign $+1$.
All reductions use that same place. This is a fixed-witness statement,
not a bound for arbitrary points or arbitrary reduction primes.

The weighted ternary theta construction gives actual holomorphic
forms at these stages. At original level primes the root weight uses
the summand selected by the cyclic level structure. At two, the
integral spectral projectors of $(1+x)/2$ specify that summand; the
operator itself is not in general idempotent. The nilpotent weight at
$h_*$ and the division-algebra residue-line weight at $\ell_0$ have the
prescribed conjugation and uniformizer signs. Their product is even.
The root-character field is fixed before $\ell_0$ and $M_\theta$.

Appendix~\ref{app:odd-source} gives these functions, their exact
transformation rules, and the direct extraction using the form $Q/m$.
Poisson summation for the resulting compactly supported local functions
gives the scalar transformation law and holomorphy at every cusp.
A fixed stabilizer multiple, for example one divisible by
$|\operatorname{GL}_4(\F_3)|$, clears the finite automorphism groups.
For the Hecke recurrence, a split or irreducible nonzero reduction has
$1+(-n/p)$ invariant lines, a nonzero nilpotent has one, and zero has
$p+1$. The unique integral primitive nilpotent neighbor supplies the
remaining square term. Tracking the scaling characters gives the
recurrence with the same character convention as the coefficient
comparison.

For an admissible $D$ with $(D/\ell_0)=1$, the oriented embedding
count \cite[Theorem 3.1]{cornutjetchev}, with changed set $\{\ell_0\}$, gives
\[
 a_{M_\theta}(D)
   =c_0\,2^{\omega(N)}\zeta_{M_\theta}(D)
       \lambda_{M_\theta}(\overline{LP_D^{\rm gen}})
       \pmod {2^{M_\theta}},                                    \tag{3.8}
\]
where $c_0\ne0$ is fixed and $\zeta$ is a permitted root-of-unity phase.
For admissible $D$ with $(D/\ell_0)\ne1$, including $\ell_0\mid D$,
the local Schwartz support instead gives $a_{M_\theta}(D)=0$.
Formula (3.8) is asserted only on the first support.
There is no division by a class number. Formula (3.7) makes $D_w$
visible with fixed upper valuation in (3.8). It remains on every
stage's support, including the extra test $(D_w/\ell_0)=1$.

After phase removal, the coefficients used for valuation are integers.
Equation (3.6) gives a uniform lower bound for
$v_2a_{M_\theta}(D)-\omega(D)$ on visible bounded-weight sequences.
Take the finite integer values constant on a set belonging to the
chosen nonprincipal ultrafilter. This set of values is nonempty by
$D_w$, is bounded below, and hence has an attained minimum $\delta$.
If the resulting lower inequality failed for a fixed weight bound on
an ultrafilter-large set, choosing a violating index and then one of
finitely many smaller integer values would contradict minimality.
For any fixed finite operator budget, the coarse lower bound handles
the large weights and stabilization handles the bounded weights.
Consequently every required division is valid at one sufficiently
advanced actual stage. No infinite-precision modular form or
arbitrary-ring base-change assertion is needed.

\subsection{One common family and the final order of choices}
The even source is a fixed algebraic vector in the full good-Hecke
eigenspace. Waldspurger's Corollary 2, including its same-dyadic-squareclass
removal of the auxiliary condition, compares the actual squarefree
Fourier functionals \cite{waldspurger}. Local upper-unipotent projections
retain a nonzero initial functional; smoothness supplies lower-unipotent
invariance. The algebraic linear conditions choose a vector over a
number field. Galois conjugation followed by the matching finite
character twist shows that each normalized coefficient ratio is a
rational number times a two-power root of unity. This supplies the
valuation and nonvanishing comparison on the whole filter.
Only flat field extension of the algebraic model is used.
The comparison is proved at the original support $S_0$ before any extra
coprimality sieve $B$. The sieve is applied last, and every later good
recurrence prime avoids the full support $S_0\cup B$.

For both coefficient systems the integral operators have the required
omission and contraction rules. At a new prime $V_p$ raises the level
exponent only to one, with ramified quadratic character. The recurrence
makes the relevant differences divisible by the integer two before
residue reduction. Multiplying by ordinary theta, then taking a fresh
trace, removes the integral unary corrections. Two fresh traces vanish.
The matrix trace identity gives $t(gu^2)=t(g)$, hence elementary
quadratic labels. At a new odd level prime the weight-two local
representation theorem \cite{carayol} has one unramified and one
ramified line; its Euler factor identifies the $U_p$ eigenvalue.
The diagonal-matrix calculation gives the two inertia differences.
These identities prove omission and primary isolation for the same
original form and depth, without choosing another minimum.

If the intrinsic odd limit symbol is a unit at the empty configuration,
choose an actual stage carrying that value with $M_\theta>\delta_o$.
The index $D=1$ lies in every moving support and then satisfies
$v_2a_{M_\theta}(1)=\delta_o<M_\theta$. The visibility implication
(O.8) gives $\ord_{s=1}L(E,s)=1$, which is exactly (RC).
No final companion or determinant bound is needed in this case.
For the rest of the final-family and missing-address argument, assume
that the empty configuration is not a unit. The bounded-weight witness
then supplies a least nonempty odd seed, as in
Appendix~\ref{app:endpoints}.

\paragraph{Fixing the final companion.}
Use a separate even symbol on negative discriminants square at $2N$,
excluding the fixed odd support, the ramification support of $E[8]$,
and the exceptional imaginary fields. Choose a unit $k$ of this
symbol. Then $L(E^{(k)},1)\ne0$. Fix $k$ and $K=\Q(\sqrt{k})$
before choosing the binary dimension. At least one prime ramified
in $K$ is outside all the old fixed label extensions and torsion
ramification. Choose every later moving theta prime $\ell_0$ away
from $k$. The intrinsic odd symbol continues to omit the primes of $k$.

At any such stage, the normal closure of the old odd label extension
is unramified outside its old support and $\ell_0$. The fresh ramified
prime of $K$ lies outside this set, so its intersection with $K$ is
$\Q$. Thus each old odd label still has a lift split in $K$.
The added stationary tests at the primes of $k$ may have relations;
include those relations in the common alphabet, rather than requiring
independent arbitrary prescriptions. Stabilize its separate projection
maps on the same ultrafilter. This preserves the old odd minimum.
Fresh ramification also makes $K$ independent of $\Q(E[8])$ and of
the old fields $\Q(i)$ and $\Q(\sqrt{\ell})$.

The forward theorem at this known nonzero central value gives
$s_2(E^{(k)}/\Q)=0$. The sum-and-difference isogeny over $\Q$ gives
\[
 s_2(E/K)=s_2(E/\Q)+s_2(E^{(k)}/\Q)=1.
\]
Indeed its kernel and cokernel on the relevant restriction and
corestriction maps are killed by a fixed power of two. Tensoring the
compact Selmer modules, equivalently the Pontryagin duals of the discrete
Selmer groups, with $\Q_2$ gives the direct sum and hence this identity
of coranks.
No integral direct-sum assertion is used.

The $K$-split even seed is nonempty. Otherwise its $p^3$ coefficient
and a fresh trace would force $a_p(E^{(k)})\equiv p+1\pmod8$ for every
good split $p$. Choose $p\equiv1\pmod4$ with
$(\ell/p)=-1$ and $(k/p)=1$. The three rational points of order two
cannot be doubled, since their necessary square differences include
$-\ell,\ell,-4\ell$. The two-primary reduction group therefore has
order exactly four, contradicting that congruence.

We record the finite combinatorial realization used here. Each label
space is a finite elementary two-group; the common alphabet projects
onto each separately. Product surjectivity is unnecessary. Take
sufficient repeated copies of every label, with total label zero.
Keeping the seed anchors, contract the unused copies into the last
anchor. Minimal positive seed size makes the internal bits irrelevant,
so each of the two terminal polynomials is nonzero. Grounded Laplacian
minors express the constant and pair-contraction operators, including
singular matrices. A cyclic singular network isolates its target
radical and a path network isolates its selected pair coefficient.
In the product of the two top forms, taken in the ordinary polynomial
ring before Boolean reduction, a variable of exponent two uses one
singular program in each factor; exponent one uses one varying program.
The total degree kills the other contractions and leaves coefficient
one. Thus one choice works for both forms. CRT and Dirichlet realize
the prescribed forward residue bits by distinct primes; quadratic
reciprocity supplies the reverse bits.

In this remaining branch, for every sufficiently large $b$ there is one
finite set $Q$ of distinct good odd primes split in $K$ and outside the
fixed support, and nonzero linear forms $\lambda_q:\F_2^b\to\F_2$,
such that the family
\[
 h_x=\prod_{q\in Q}(q^*)^{\lambda_q(x)}
 \quad(x\in\F_2^b),\qquad q^*=(-1)^{(q-1)/2}q,
\]
has $h_0=1$, all fixed local types unchanged, and, for $x\ne0$,
\[
 \ord_{s=1}L(E^{(h_x)},s)=1,\quad L_*(h_xk)\ne0,\quad
 j_{\rm gen}(h_x)\le\omega(h_x)+C_o,\quad
 v_2L_*(h_xk)\le2\omega(h_x)+C_e.                                \tag{3.9}
\]
The constants are independent of $b$. By (3.4), now with fixed $k$,
$2j_{\rm rc}(h_x,k)\le4\omega(h_x)+C_J$.

The order is fixed data and $D_w$, then the genus lower bound and odd
minimum, then the final $K=\Q(\sqrt k)$ and common alphabet, then
determinant precision $M$, then $b$ and its finite operator budget,
then a sufficiently advanced actual theta stage $M_\theta$ and its
prime realization. In particular $M_\theta$ is not $M$.
The finite graph construction, its local conditions and the coefficient
tests all occur at that same final stage. Appendix~\ref{app:endpoints}
spells out the absolute five-prime completion, normalization constants,
and the final use of the one corrected coordinate.

% End 03-restricted-converse.tex


% Begin 03a-arithmetic.tex
\section{Actual classes, local conditions and central comparison}
\label{sec:arithmetic}
Let $K=\Q(\sqrt k)$ be imaginary, split at $2N$, and different from
$\Q(i)$ and $\Q(\sqrt{-3})$. Active primes are good odd primes split
in $K$. Derivative primes are good inert primes outside all current
supports. The actual finite coefficients are
\[
 A_i=(\Z/2^i\Z)[t]/t^i,\quad
 M_i=(T_2E\otimes A_i)(\chi_h\Psi_i).
\]
For a fixed binary group $G$, use $A_i[G]$ and the universal quadratic
character. The character $\Psi_i$ is finite, trivial at $2N$, ramified
at a moving conductor prime $r_i$, and has prescribed odd scalar
exponents at the finite active list. The prime $r_i$ varies with
precision. There is no assertion that one fixed odd prime supports
an infinitely ramified $\Z_2$ character.

\subsection{One multiplier before any conductor varies}
Choose $n\ge2$ with $1+2^nM_2(\Z_2)$ contained in the Tate image, and
put $u=1+2^n$. A scalar lift $g$ has Tate action $uI$.
For complex conjugation $\tau$ set
\[
 z=g^2\tau g^2\tau^{-1}.
\]
The element $g^2$ fixes every quadratic field. Conjugation acts by
inversion on every ring-class quotient of an imaginary quadratic field,
so $z$ fixes each such ring-class field and every rational quadratic
character. Its Tate action is $u^4I$, with
$v_2(u^4-1)=n+2$. Thus the relevant invariant torsion has one bounded
exponent, independent of $K$ and conductor.

The required matrix-span bound is equally uniform. Squares of
preimages in the open subgroup give, inside every quadratic subgroup,
the matrices
\[
 I,\quad I+(2^{n+1}+2^{2n})E_{11},\quad
 I+2^{n+1}E_{12},\quad I+2^{n+1}E_{21}.
\]
Their span has fixed finite index in $M_2(\Z_2)$.
These facts are the non-CM input. They are never applied to the CM
generator curve.

Put $e=n+2$. In inflation--restriction for the complete conductor
extension, both the obstruction and ambiguity groups with invariant
coefficients are killed by $2^e$. Lift $2^e$ times an invariant Kummer
class, then multiply its lift by $2^e$. The resulting descended class
is independent of the lift, is additive, and restricts to $2^{2e}$
times the original class. Apply this once to the complete set of
derivative primes, not once per prime.

At an original bad place, these conductor extensions are unramified.
The first cohomology of the identity component vanishes under finite
unramified cyclic descent: Lang's theorem handles the connected
special fiber and the normal-basis calculation handles each additive
formal quotient. Completeness passes through the formal filtration.
The component-group order therefore kills the remaining obstruction.
Only finitely many local quadratic types occur at $2N$.
Together with the fixed cuspidal correction of $\pi$, their component
orders and $2^{2e}$ give one integer $L$. It is fixed before
$K,h,r_i$, derivative labels and binary dimension vary.

\subsection{Norms, derivatives and retained boundaries}
Let $H_c$ be the ring-class field and $y_c$ the consistently oriented
CM point mapped by $\pi$. For a new prime not dividing $6cND_K$ the
unaveraged norm identities are
\[
 \operatorname{Tr}y_{c\ell}=a_\ell y_c\quad(\ell\text{ inert}),\qquad
 \operatorname{Tr}y_{cq}=(a_q-\sigma_q-\sigma_q^{-1})y_c
       \quad(q\text{ split}).                                  \tag{4.1}
\]
The arbitrary-old-conductor version is
\cite[Theorem 2.15]{ccl}; no squarefree-inert restriction on $c$ is
inserted. The unit factor is one by the exclusions on $K$.

At an inert prime the reduction identity is
$\overline y_{c\ell}=\operatorname{Fr}_\ell\overline y_c$
for a compatible place. Locally $\mathcal O_c\otimes\Z_\ell$ is the
maximal unramified quadratic order. Its $\ell+1$ raising isogenies
reduce to the supersingular Frobenius quotient, and prime-to-$\ell$
level structure is transported by Frobenius. The parametrization
commutes with reduction. This proves the same identity after all old
class-group translates.

For $D_\ell=\sum_{j=1}^{\ell}j\sigma^j$,
\[
 (\sigma-1)D_\ell=(\ell+1)-\operatorname{Tr}.
\]
When the chosen two-power divides $a_\ell$ and $\ell+1$, (4.1)
gives the invariant class to which the preceding descent applies.
The finite and singular coordinates satisfy, at retained precision,
\[
 f=(a_\ell F-(\ell+1))\overline Q,\qquad
 s=(F-a_\ell)f+\varepsilon,\qquad 2^e\varepsilon=0.
\]
Here $F^2-a_\ell F+\ell=0$. Taking the initial precision at least
$e$ larger, and $F$ congruent to the chosen conjugation matrix $J$
with one additional bit, gives
\[
       f(P_{I+\ell})=0,\qquad s(P_{I+\ell})=Jf(P_I).               \tag{4.2}
\]
The vanishing follows also from
$\overline{D_\ell y_{c\ell}}=\ell(\ell+1)F\overline y_c/2$.
The extra bit pays for this division by two.

Weighted transfer is the Shapiro map followed by
$v\otimes[\delta]\mapsto\theta(\delta)v$. It is a finite sum without
division by a class number. Local Mackey decomposition transports the
places and Frobenius, so (4.1)--(4.2) hold term by term before
weighting, also over the universal binary group ring.
Whenever equality in the Selmer complex is asserted, retain an actual
boundary witness at the same finite stage.

\subsection{The derived Kummer condition}
For the actual local twist $B=E^{(h)}$ over $F=K_v$, put
\[
 U_F=B(F)^\wedge_2[-1]\longrightarrow R\Gamma(F,T_2B).
\]
The completion includes finite two-primary torsion. Derived reduction
modulo $2^a$ has $B(F)[2^a]$ in degree zero and
$B(F)/2^aB(F)$ in degree one. At $v\mid2N$, $\Psi_i$ is trivial,
so derived coefficient extension defines the scalar $A_i$ condition
for this actual twist. In the universal binary family every $h_x$
has the same local squareclasses at $2N$ as $h_0=1$. Hence every
sign character evaluates the universal character $\varepsilon$
trivially on $G_{K_v}$, and therefore $\varepsilon|_{G_{K_v}}=1$
in $G$, before reduction modulo two. The $A_i[G]$ Kummer condition
and its null-homotopy are consequently derived scalar extensions
of the fixed local data. Omitting the scalar twist, the universal
local-squareclass condition, or degree-zero torsion changes this argument.

The local finite dualities \cite[I.2.3, I.3.2--I.3.5]{milne} make the
Kummer map adjoint to $H^1(F,B[2^a])\to H^1(F,B)[2^a]$.
The latter has Kummer image as kernel, so the Kummer image is its exact
annihilator. To obtain the integral chain statement, write
$M=B(F)^\wedge_2$. The restricted cup map
$M[-1]\otimes^LM[-1]\to\Z_2[-2]$ has homotopy class determined by
$\operatorname{Hom}(M\otimes M,\Z_2)$. Finite isotropy makes every
value divisible by every power of two, hence zero.
A projective resolution therefore supplies a null-homotopy.
The possible ambiguity group in degree minus one vanishes.
Reduction modulo two, including the torsion in degree zero, gives the
perfect quotient-to-dual comparison; derived Nakayama lifts it.
The global Selmer duality theorem
\cite[Sections 5.1, 5.2.4, 5.4.3 and 6.3.4]{nekovar} is applied
to these actual local maps and homotopies.

For clarity the Selmer complex is the cone
\[
 C_U=\operatorname{Cone}\left(
 R\Gamma(G_{K,S},M_i)\oplus\bigoplus_{v\in S}U_v
 \longrightarrow\bigoplus_{v\in S}R\Gamma(K_v,M_i)\right)[-1].
                                                                    \tag{4.3}
\]
The compatible trace in (4.3) is coefficient-valued. A top-coefficient
functional used to test a finite Frobenius ring's perfection is not
used as a coefficient-reduction-compatible arithmetic trace.

\subsection{Actual support comparisons}
At a tame split prime the local free model has
\[
 d^0(v)=(Av,Cv),\quad d^1(v,w)=-Cv+Dw,\quad
 A=\Phi-1,\ C=g-1,\ D=\Phi-N_q(g),
\]
\[
 N_q(g)=\frac{q+1}{2}+\frac{q-1}{2}g,\qquad DC=CA.                 \tag{4.4}
\]
At augmentation $g=1$, this is the direct sum of the unramified
complex $[\Phi-1]$ in degrees zero and one and the singular quotient
$[\Phi-q]$ in degrees one and two.

The support-change theorem \cite[Proposition 7.8.8(ii)]{nekovar}
retains arbitrary old local conditions and puts the unramified
condition at new places. Its finite-module proof applies at two.
Use it only when the actual coefficient representation is unramified
at the removed place. There are five uses here:
universal augmentation at $Q$; an unused scalar prime $q\nmid h$;
characteristic-two residual removal after $\chi_h=1$; the moving
prime after $t=0$; and a new good inert derivative prime outside
$h,r_i$ and the old support. In the last case the coefficient module
is unramified and the chosen finite local plane is the unramified
condition. It is not applied to an active ramified scalar prime
or to $r_i$ before central specialization. Replacing the unramified
condition by the full condition then gives the singular quotient
by a cone of Selmer cones. In particular the old complex is an actual
subcomplex of an upper triangular model, not just a complex with the
same Euler characteristic.

The simultaneous character construction solves the finite ideal
relation lattice before choosing $r_i$. Prescribed odd valuations at
the active primes are compatible with the Kummer radicals; the
cyclotomic tower is unramified there. On the Tate side, commutators
of a fixed open subgroup give a fixed open subgroup in the allowed
coset. The functions $\operatorname{tr}\pm2$ are not identically zero
on any such open coset. There are only finitely many cosets at the
fixed level, giving a bound $c_E$ independent of the active list.
At each coefficient stage Chebotarev realizes the single compatible
target. Thus the same $\Psi_i$ has odd deformation exponent at every
active prime and bounded central moving-prime determinants.

\subsection{The same functional at every localization}
After the inner limit, put
$\mathscr O=\widehat{\Z_2[[t]]_{(2)}}$.
In (4.4), the residual active Frobenius is $(1+t)^aI$ with $a$ odd.
Thus $A,D$ are invertible over $\mathscr O[G]$.
The formulas
\[
 h^1(v,w)=A^{-1}v,\qquad h^2(u)=(0,D^{-1}u)
\]
contract the complex; the identities follow from $DC=CA$.
Combining these local comparisons with the retained cup map gives
a closed functional $\lambda$, with $2^a\lambda$ integral for one
fixed exponent $a$. This is a statement after the inner limit and
base change, not an inverse of $t$ in $A_i$.

For every height-one prime $\mathfrak p\ne(2)$ of $\Lambda$, construct
a representative separately. At an active scalar prime $C=-2$ is
invertible. At $r_i$, $C=t$ times a unit is invertible away from
$(t)$. There use
\[
 h^1(v,w)=C^{-1}w,\qquad h^2(u)=(-C^{-1}u,0).
\]
At $(t)$ the bounded central Frobenius determinants permit the first
contraction by adjugates. This includes primes such as $(t+2)$:
after inverting two, $t$ is a unit there. All representatives invert
the same generic comparison and compose with the same cup map.
Their equality in the generic derived category, not
$\mathscr O$-integrality alone, identifies the generic functional.

For a constituent whose generic cohomology has dimensions $(1,1)$
in degrees one and two and vanishes otherwise, use the determinant
line
\[
 \mathcal D=\det\nolimits^{-1} C
       \simeq H^1(C)\otimes H^2(C)^*.
\]
The class-functional coordinate $u$ is the coordinate of
$[z]\otimes[\lambda]\in\mathcal D$ in the retained integral
determinant basis. Here $z$ is the actual class and $\lambda$ is
the closed functional just constructed. For a generically acyclic
constituent set $u=0$. The primitive and relaxed subscripts
specify the corresponding local conditions on $C$.

\subsection{Central identity, torsion and the primitive coordinate}
The split norm relation gives an actual finite-stage boundary
\[
 z_i(0)-Lm_i\,j_i(b_h)=dv_i,\qquad
 m_i=a_{r_i}-2\chi_h(r_i).                                      \tag{4.5}
\]
The equality in global cohomology lifts to the Selmer cone because
the local degree-zero maps are onto and the degree-one maps inject.
Retain $v_i$. The uniform trace bound bounds $v_2(m_i)$ and the local
determinants by $c_E$; their adjugate identities multiplied by a fixed
power of two survive the limit.

For the ordinary compact Kummer complex $C_K$, universal coefficients
and the local degree-zero identifications give
\[
 H^1(C_K)_{\rm tors}=E^{(h)}(K)[2^\infty],\qquad
 0\to H^1(C_K)\otimes\Q_2/\Z_2\to\Sel_{2^\infty}
       \to H^2(C_K)_{\rm tors}\to0.
\]
The last finite group is Selmer modulo its maximal divisible subgroup.
Self-duality identifies the degree-three torsion with the degree-one
torsion and gives a unimodular pairing of the free degree-one and
degree-two parts. Thus their rational dimensions equal full Selmer
corank, including a possible divisible part of $\Sha$.

At known simple zeros, forward Gross--Zagier and Kolyvagin
\cite{gz,kolyvagin} give rank one and finite $\Sha$.
For restriction of scalars and the product of the two twists, the
sum/difference isogenies are dual and compose to $[2]$.
Their free-lattice indices satisfy $IJ=2$, with regulator ratio
$2/I^2$. At odd primes the differential lattices agree; at two the
pullback determinants have bounded total valuation.
Milne's isogeny ratio \cite[proof of I.7.3, (7.3.1)]{milne} follows
from duality without assuming the BSD leading-term equality.
It applies here with finite $\Sha$, canonical heights and pullback
measures.

At every active $p\mid h$, $K$ splits, so the isogeny local factors
cancel exactly. At either split place the full/Kummer quotient
has length
$\dim_{\F_2}E[2]^{\operatorname{Fr}_p}=2$, since all of $E[2]$
is rational. The two places therefore contribute exactly $-4$. At a companion prime $p\mid k$ the
ramified twist has component factor four, while the restriction of
scalars has connected special fiber. Hence its error $t_p$ is at most
two. The result is
\[
 v_2u_{\rm prim}(0)
 =2j_{\rm rc}(h,k)-4\omega(h)-s(h)-s(hk)
                    -\sum_{p\mid k}t_p+O_{E,\pi}(1).             \tag{4.6}
\]
Only the companion sum can be absorbed when $\omega(k)$ is bounded.
No class number or unbounded degree-two torsion length is absorbed.

At an unused prime the exact point factor is $e_q=a_q-z-z^{-1}$.
The second factor comes from the same closed functional, as follows.
Let $i:C_P\to C_R$ be the primitive-to-relaxed map, and let
$j:C_S\to C_{P^\perp}$ be the paired strict-to-primitive map.
All old orthogonal conditions and their homotopies are retained.
At $q$ these are the maps from the unramified condition to the full
condition and from zero to the unramified condition. Naturality of
the cup pairing \cite[Sections 6.4.1--6.4.3]{nekovar} gives
\[
             \beta_R(v,iw)=\beta_P(jv,w).
\]
The actual norm identity is $z_R=e_q i(z_P)$, with its retained
boundary. The generic paired strict lift satisfies
$jz_S=e_q^c z_{P^\perp}$. Conjugation sends $z$ to $z^{-1}$, so
$e_q^c=e_q$. Consequently
\[
                 \lambda_R\circ i=e_q\lambda_P.
\]
Thus the class and its paired functional each contribute one
$e_q$; no integral self-duality of the primitive condition is assumed.
The two quotient determinants satisfy
\[
 \delta_q=\det(z\operatorname{Fr}_q-q)
              \det(z^{-1}\operatorname{Fr}_q-q)=q^4l_q.
\]
In the original integral determinant bases, the support triangle
therefore gives
\[
       l_q u_{\rm rel}=e_q^2\varepsilon_q u_{\rm prim},
             \qquad \varepsilon_q\in\Lambda^\times.             \tag{4.7}
\]
To see the unit assertion, retain the finite mapping-cone bases.
Their comparison is a sign; replacing them by the retained homotopy
equivalent free models contributes only determinants of unit disks.
These are $A_i$-units. Their limits and inverse limits are
$\Lambda$-units, and $q^{-4}$ is one also. An arbitrary
$\mathscr O$-unit would not suffice.

The universal correction $\eta$ is integral over $\mathscr O[G]$:
$l_q/e_q^2=1+2v_q$, and $1+v_q(1+g_q)$ selects that factor precisely
at an inactive character. Put $\widetilde u=\eta u_{\rm rel}$.
It already is one common family; no interpolation of the individual
$\varepsilon_x$ is required. Formula (4.7) makes
$\widetilde u_x=\varepsilon_xu_{{\rm prim},x}$.
The central zero and valuations therefore transfer unchanged.
The multiplier $\eta_x$ may itself have a pole at $t=0$.
Only the product $\widetilde u_x$, identified with the regular
primitive coordinate times a $\Lambda$-unit, is specialized there.

Finally, suppose the unknown base derivative vanishes.
The conductor-one height formula makes its point torsion.
Use a torsion annihilator and (4.5) to retain its central boundary.
The central rational dimensions in degrees one and two are $(1,1)$
by full Selmer corank. This does not assert that the central $H^1$
vanishes: the actual point class is zero because the point is torsion.
The generic complex is acyclic or has dimensions $(1,1)$.
In the acyclic case the specified class-functional coordinate is zero.
In the latter case equality of central and generic dimensions rules
out a nonunit disk over $\Lambda_{(t)}$: each such disk would add
a central class in both adjacent degrees, including its Tor
contribution. After cancellation, determinant base change therefore
identifies the same coordinate with the central class paired with
the central functional. The retained boundary makes this zero.
Thus $u_{{\rm prim},0}(0)=\widetilde u_0(0)=0$.
No finite-$\Sha$ assertion at the unknown base enters this argument.

% End 03a-arithmetic.tex


% Begin 03b-limits.tex
\section{The two limits and determinant transfer}\label{sec:limits}
This section uses the actual diagrams of Section~\ref{sec:arithmetic}.
Their cohomology ranks alone would not suffice: maps to original
cochains, local comparisons, cup products, classes and boundary
witnesses are retained together. The same retained data include the
integral duality maps, their homotopy inverses, and the homotopies
expressing compatibility of duality with localization.

\subsection{Finite free models and the inner limit}
Let $A$ be an Artin local ring with nilpotent maximal ideal $\mathfrak m$.
A flat $A$-module is free, including in infinite rank: lift a basis
modulo $\mathfrak m$, use nilpotence for surjectivity, and flatness and
nilpotence for vanishing of the kernel. A continuous finite-coefficient
cochain module is a filtered union of modules of functions on finite
clopen partitions, hence is flat.

Suppose the residual cohomology of a complex of such modules is finite
and bounded. Split the residual complex into cohomology and disks,
and lift graded bases. Let $d_0$ be the lifted disk differential and
$d=d_0+\epsilon$, where $\epsilon$ has entries in $\mathfrak m$.
Use the graded decomposition into cohomology, disk sources and disk
targets. Let $i_0,p_0$ include and project the cohomology summand,
and let $h_0$ invert $d_0$ from disk targets to disk sources, zero
on the other summands. Then
\[
 p_0i_0=1,\quad h_0^2=h_0i_0=p_0h_0=0,\quad
 1-i_0p_0=d_0h_0+h_0d_0.
\]
With this normalized contraction, the finite geometric series give
\[
 \begin{aligned}
 d_C&=p_0\epsilon(1+h_0\epsilon)^{-1}i_0,&
 i&=(1+h_0\epsilon)^{-1}i_0,\\
 p&=p_0(1+\epsilon h_0)^{-1},&
 h&=h_0(1+\epsilon h_0)^{-1}.
 \end{aligned}
\]
Substitution using $d^2=0$ gives $pi=1$ and $1-ip=dh+hd$.
The resulting bounded finite free model carries actual maps and
homotopies to the original cochains. Transport any fixed finite
collection of arithmetic maps by $f\mapsto p'fi$. The identity
$1-ip=dh+hd$ transports their compatibility homotopies as well,
including the second homotopies in the cup-product diagrams.

Fix one nonprincipal ultrafilter on coefficient stages. Bounded ranks
can be fixed on a set in that ultrafilter. Each finite tuple of matrix
entries then has a compatible limit at every coefficient precision,
giving matrices over $\Lambda$. Every polynomial identity between
the matrices survives. Retained homotopy inverses remain homotopy
inverses. Mere stagewise quasi-isomorphisms without those maps would
not justify a nonflat coefficient change.

The residual rank estimate needs a precise order of operations.
Over $B_i=\F_2[t]/t^i$, the quadratic character $\chi_h$ is literally
trivial. Thus an actual smaller-support representation exists, with
old support, two places over $r_i$, and the fixed number $j$ of
derivative labels. Cancel only units over $B_i$.
The index-two pro-two generator bound controls its residual $H^1$.
It does not by itself control $H^2$. For that use the totally imaginary
global Euler characteristic $\chi=-2$ and vanishing above degree two.
Let $v=\omega(2N)$ and let $s$ count rational primes in the smaller
ramification support, including primes dividing $\operatorname{disc}K$,
the moving prime and $j$ derivative labels. Since $K=\Q(\sqrt k)$
and two is already counted in $2N$, the condition $\omega(k)\le A$
gives $s\le v+A+j+1$. This adds the quadratic field's ramification
to the rational generator count, not new Selmer local conditions.
The local contributions
give the explicit bound
\[
 N(E,A,j)=8(s+1)+6+48(v+1+j)
          \le56v+8A+56j+70.                                    \tag{5.1}
\]
For each fixed active support and family, retain the actual comparison
cones, maps and homotopies in a padded tuple whose ranks are bounded
over coefficient stages, and pass to $\F_2[[t]]$. The ranks of this
whole retained tuple may grow with the number of active cones.
The bound $N(E,A,j)$ concerns the reduced generic model, not that
whole tuple uniformly in the active support. Only now invert $t$.
The active local determinants are nonzero because their scalar
exponents are odd; their cones become acyclic. This proves the generic
residual bound for the original complex without inventing a fixed-$r$
infinite character. The bound is independent of the number and sizes
of active primes and of $v_2(r_i-1)$.

\subsection{Evaluation and the outer valuation ring}
For fixed $K$ let $\Gamma_K=\prod_iG_K$ and include one simultaneous
complex conjugation. Evaluation of retained cochain maps at arbitrary
sequences $g=(g_i),h=(h_i)$ gives
\[
 e^1_gd=(\rho(g)-1)e^0,\qquad
 e^2_{g,h}d=\rho(g)e^1_h-e^1_{gh}+e^1_g.                          \tag{5.2}
\]
One fixed ultrafilter supplies these identities for all sequences.
A cycle therefore defines an abstract crossed homomorphism of the
product group. We do not assert that its restriction to one original
Galois group is a continuous pointwise limit.

Evaluation remains injective after field extension. Residually, take
representatives $V$ for $H^1$ and coefficient space $W$. Evaluate
$(c,w)$ by $b_c(g)-(\rho(g)-1)w$. The intersection of all kernels is
$(0,W^\Gamma)$. A finite list of at most $\dim V+\dim W$ evaluations
already has that intersection. The comparison to $[W\to W^d]$ is an
isomorphism on $H^0$ and injective on $H^1$. Its cone has no residual
cohomology in degrees at most zero. Cancel unit disks there so that
the free cone starts in degree one; the same is true after every
coefficient extension. The long exact sequence gives the claimed
injection. For Selmer complexes this applies because the local maps
are isomorphisms on $H^0$ and injections on $H^1$.

Now allow $K_j,h_j$ to vary with $\omega(k_j)\le A$.
For copies $\mathscr O_j$ of $\mathscr O$, form
\[
 \mathscr O_\infty=
 \left(\prod_{\mathcal U}\mathscr O_j\right)/
 \{(x_j):v_2(x_j)\to_{\mathcal U}\infty\}.                         \tag{5.3}
\]
A nonzero element has valuation one fixed nonnegative integer on a
set in $\mathcal U$ and is that power of two times a unit.
Every nonzero ideal has an element of least valuation and is
principal. Thus (5.3) is a DVR with uniformizer two and characteristic-zero
fraction field. Completeness is unnecessary.

Integral minimal-model matrices and (5.2) pass to (5.3).
So do the retained duality maps, inverse homotopies and
duality-to-localization compatibility homotopies. After minimal
cancellation an integral homotopy equivalence is degreewise invertible:
its residual matrices are invertible, and Nakayama applies.
These inverse matrices and all compatibility identities pass to the
same outer quotient. Thus perfect duality and exact orthogonality
are preserved there; they are not inferred merely from isotropy.
Approximate cycles suffice: if $v_2(dw_j)\to\infty$, the degree-two
term in (5.2) tends to zero, so $w_\infty$ is an actual crossed
homomorphism. The finite detecting lists have bounded length by
(5.1); retain their cone contractions to preserve injection.
One must not replace approximate cycles by unexplained exact cycles.

\subsection{One new prime detects two classes}
The exact kernel consists of sequences trivial on the full Tate
module, all relevant ring-class fields and quadratic characters.
The scalar element of Section~\ref{sec:arithmetic} is central modulo
this kernel, though not in the kernel itself.
For a cocycle zero on the kernel,
\[
 (u^4-1)b(g)=(g-1)b(z)
\]
makes it a coboundary over characteristic zero. The same fixed scalar
survives the outer quotient. To eliminate residual invariants, use moving-prime
inertia: its scalar difference is $t$, invertible in $\F_2((t))$.
These are two separate arguments.

On the characteristic-zero tests used for rank reduction, the fixed
finite matrix-span witness gives absolute irreducibility of each
rank-two summand. When the scalar character is not quadratic, moving
inertia has distinct characters on the two summands and simultaneous
conjugation exchanges them. The induced rank-four representation is
therefore absolutely irreducible with scalar endomorphisms. The outer
fraction field is one such characteristic-zero test; the quadratic
character exception at an inner test is treated separately below.

For two independent cohomology classes on these tests, the span of
their joint values on the exact kernel projects onto both copies of
this absolutely irreducible representation. A proper submodule with
both projections onto would be the graph of a scalar isomorphism,
contradicting independence. Thus the span is the full sum. If the
determinant of their two finite localizations vanished on every actual
kernel element, additivity and polarization would make it vanish on
the entire span, impossible. Hence one actual element supplies a
rank-two localization on these characteristic-zero tests.

The polarization identity itself is also valid in characteristic two,
but no residual irreducibility or residual rank-two localization is
asserted. Here $E[2]$ is trivial of rank two, so its scalar twist is
isotypic and the induced residual representation has multiplicity two.
At the residual test we use only scalar-inertia vanishing of $H^0$,
the bounded residual cohomology dimensions, and finite detecting lists
from intersections of linear evaluation kernels. Those lists need not
consist of one rank-two localization. The rank-two selection for the
uniform valuation bound occurs only after passage to the
characteristic-zero fraction field of the outer DVR.

Chebotarev is applied at finite stages, not to the product group.
Each of the finitely many original basis cochains in degrees zero
through two factors through a finite quotient. Include their quotient,
the required Tate precision, ring-class and quadratic fields and old
diagrams in one compositum. The element just found, multiplied by
conjugation, is a single actual automorphism of that compositum.
Its restrictions agree on intersections. Choose the new prime and
a place above it to realize that representative.

This also resolves the precision issue caused by localization.
Choose the new prime to match all original basis-cochain evaluations
at a common retained precision $n(i)\to\infty$. Taking the same inner
ultrafilter gives an exact equality of $\Lambda$-matrices
\[
                 L_j=E_{\sigma_j^2}.                            \tag{5.4}
\]
Only afterward compose with the inclusion $I_j$ from an
$\mathscr O$-minimal model:
$L_jI_j=E_{\sigma_j^2}I_j$.
The denominators in $I_j$ may depend on $j$; equality (5.4) makes this
irrelevant. Long $t$-jets are never treated as approximations in the
two-adic valuation of $\mathscr O$, where $t$ is a unit.

\subsection{Switches and the varying-companion estimate}
At a new inert prime the finite and singular planes are free of rank
two, with cross pairing
$\beta(x,y)=\varepsilon\langle x,Jy\rangle$.
Suppose the old actual class $A$ spans generic $H^1$ and has
$f(A)=v\ne0$. Its derivative $B$ satisfies $f(B)=0$, $s(B)=Jv$.
The strict $H^1$ vanishes. The transported duality and localization
identities identify the relaxed localization image with the kernel
of its dual localization map, hence with its exact annihilator
\cite[Sections 6.3.3--6.3.4]{nekovar}. This image
contains $(v,0)$ and $(0,Jv)$, so these span
it. Hence the new generic $H^1$ is spanned by $B$.

The determinant comparison is exact even if $v$ is divisible.
In the fraction field choose $x$ with $v\wedge x$ of unit volume,
and similarly $y$ with $Jv\wedge y$ of unit volume.
The finite and singular determinant exact sequences have the same
strict factor, so their class-functional coordinates have equal
valuations. All strict torsion remains in that common factor.
Requiring $x$ integral would be wrong when $v$ is divisible.

The starting generic kernel has rank one. Indeed, at a known simple
zero the ordinary class is nonzero. Equation (4.5), with its bounded
nonzero multiplier, makes the retained central class nonzero.
Over $\Lambda_{(t)}$ cancel the degree-zero and degree-three disks.
If the generic kernel of the resulting two-term free complex vanished,
the actual lifted cycle would vanish, contradicting its central value.
Semicontinuity bounds that kernel by the central dimension one.

Suppose no uniform lower bound exists. Choose instances with
$v_2u_j\to-\infty$, $\omega(k_j)\le A$, and write their minimal
completed models as $[\mathscr O_j^n\xrightarrow{D_j}\mathscr O_j^n]$,
fixing the bounded $n$ on an outer-large set. In each generic
one-dimensional kernel choose a primitive vector $v_j$.
Its outer image is nonzero, so the outer kernel dimension $s$ is
at least one. If $s>1$, one new derivative prime gives rank-two
localization including that vector. The relaxed localization image
is self-annihilating and contains the whole finite plane, hence equals
that plane. The strict kernel has dimension $s-2$ and the singular
switch adds no localized vector, so the new outer dimension is $s-2$.

The actual inner class is a nonzero scalar multiple of $v_j$.
Even if this scalar tends to zero, its generic localization is nonzero
on an outer-large set. Exact switch preservation therefore keeps
each inner determinant valuation. The new inner kernels again have
primitive vectors. Thus their outer dimension remains at least one;
a proposed drop to zero is already a contradiction.
There can be only finitely many drops. For each fixed number of them
(5.1) bounds the models; no bound uniform in an unbounded number of
switches is needed.

At terminal outer dimension one a maximal boundary minor $\Delta$
is nonzero. Its inner minors have one finite valuation $b$ on an
outer-large set. Determinant-one Schur operations in a square
two-term rank-one-kernel complex give
$u=\pm\alpha\beta/\Delta$, where $\alpha$ is an integral surviving
cycle coordinate and $\beta$ has valuation at least $-a$.
Consequently $v_2u_j\ge-a-b$, the required contradiction. We obtain
\[
                     v_2u_{\rm prim}\ge-C(E,A).                 \tag{5.5}
\]
Here (5.5) is asserted for the known simple-zero ring instances
used to obtain the varying-companion estimate. Their nonzero actual
central class supplied the rank-one start. It is not a uniform
valuation assertion at the unknown corank-one base.
Bounded matrix size alone would not prove it: $[2^j]$ is a
counterexample to that inference.

At a height-one $\mathfrak p\ne(2)$ the residue field has
characteristic zero. If the specialized scalar character is not
quadratic, the same argument applies. If it is quadratic, the induced
representation splits into two nonisomorphic rank-two signs.
When both occur, one evaluation detects one class from each.
When only one occurs with dimension greater than one, a first switch
preserves total dimension but introduces the other sign: the local
symmetric hyperbolic plane has exactly the finite and singular axes
as isotropic lines. A second switch then drops dimension by two.
This quadratic exception is necessary, for example at a specialization
where the moving character becomes $-1$.

Use the separate integral functional representatives constructed in
Section~\ref{sec:arithmetic}. The terminal minimal complex has one
term in each of degrees one and two and zero differential.
Its class-functional coordinate is integral over
$\Lambda_{\mathfrak p}$. Exact switch preservation transports
integrality back. Intersecting these height-one localizations gives
$u_{\rm prim}\in\Lambda[1/2]$. This regularity statement also applies
at the unknown full-corank-one base: if its generic class is zero
the coordinate is zero; otherwise its nonzero generic cycle and
the central dimension bound give the rank-one start for the
height-one argument. It does not require a uniform outer valuation
bound at that base. Combining (5.5), on the known simple-zero
instances, with (4.6) and
$\sum_{p\mid k}t_p\le2A$ proves (3.5). No fixed-field class number is
absorbed into this varying-field bound.

\subsection{Actual augmentation and the fixed-field count}
We now use the final $K$ in the remaining branch of Section~\ref{sec:rc},
where the intrinsic odd empty configuration is not a unit.
There are no derivative labels in this
last binary family. At augmentation every graph prime is unramified,
and the actual support triangle gives
\[
 C_{\rm old}\longrightarrow(C_b)_0
       \longrightarrow\bigoplus_{q\in Q}Q_q\longrightarrow C_{\rm old}[1],
                                                                    \tag{5.6}
\]
with $C_{\rm old}$ an actual subcomplex of an upper triangular model.
Each two-place local determinant modulo two has $t$-order at most
four by the same odd scalar exponents.

Let $d_K=\dim_{\F_2}\Sel_2(E/K)$,
$t_K=\operatorname{length}E(K)[2^\infty]$ and
$h_K=\operatorname{length}H^2(C_K)_{\rm tors}$.
The last number is finite without finite $\Sha$ at the unknown base.
After central specialization the two moving-prime singular blocks
have total determinant valuation at most $2c_E$.
The long exact sequence and duality give
\[
 \operatorname{rank}_{\rm total}C_{\rm old}\le2d_K+12,\quad
 m\le d_K+5,\quad
 h_1\le h_K+2c_E+2t_K.                                          \tag{5.7}
\]
These bounds depend on fixed $K$, not on $b$ or the graph size.
The $j$-dependent bound (5.1) is not used with $j=b$.

A homotopy equivalence of minimal free complexes over a local ring
is degreewise invertible, by reduction and Nakayama. Cancel unit
disks and add them back as needed to lift (5.6) upstairs. The
augmentation kernel is in the Jacobson radical, so invertible
augmented matrices lift to invertible matrices.
The block containing the local quotients has
\[
 B_0=A=\operatorname{diag}(A_q),\qquad T_0=0
\]
in a differential $\left(\begin{smallmatrix}P&Q\\T&B\end{smallmatrix}\right)$.
Over $\mathscr O[G]$, cancel $B$. The surviving Schur differential
$P-QB^{-1}T$ augments to $P_0$.

Write $\det A_q=t^{d_q}v_q+2w_q$, $d_q\le4$, with $v_q$ a power-series
unit, and put
$C_q=\operatorname{adj}(A_q)/(t^{d_q}v_q)$.
For a character $x$, $B_xC=1+2F$ with both $C$ and $F$ having no
Laurent terms below degree $-4$. Therefore
\[
             B_x^{-1}=C\sum_{j=0}^{M-1}(-2F)^j\pmod {2^M}        \tag{5.8}
\]
has pole at most $4M$, independent of the number of blocks.
The adjacent surviving differentials are unchanged by $d^2=0$.
Every entry of the remaining old-rank complex has the same pole bound.

\subsection{Determinant clearing and the missing-address contradiction}
In integral central bases decompose the old complex into $m$ Smith
disks and one free generator in each of degrees one and two.
In each differential group all disk sources into $E_i$ and all their
targets into $B_{i+1}$. Cancel the entire block
$A_i:E_i\to B_{i+1}$, once per degree, and set
\[
              \Delta=\prod_i\det A_i.
\]
A Schur complement can change the complementary part of that same
differential. It does not change the adjacent surviving differentials,
so it does not change another grouped $A_j$.
Thus $\Delta$ has degree $m$ in the original differential entries and
$v_2\Delta_0(0)\le h_1$.

The projected class, projected functional and determinant comparison
each use at most one denominator $\Delta$. For the integral
unused-prime multiplier $\eta$, there is $Q\in\mathscr O[G]$ with
\[
 \widetilde u_x=2^{-a}\eta_xQ_x/\Delta_x^3
\]
when $\Delta_x\ne0$. The integral unit from the cancelled local-block
determinant is absorbed in $Q$. Put
\[
 f=2^a\Delta^4,\qquad d=\eta\Delta Q.
\]
Then $f_x\widetilde u_x=d_x$ at every character, including
$\Delta_x=0$, where both sides vanish. Here
\[
 c=v_2f_0(0)\le c_{\det}:=a+4(h_K+2c_E+2t_K),\qquad
 L_T(M)=16(d_K+5)M                                               \tag{5.9}
\]
bounds the negative Laurent support of $f\bmod2^M$.
There is no required pole bound for the complicated numerator $d$.
The common functional bound alone would not prove (5.9); the actual
augmentation and grouped-disk calculation are essential.

For a coefficient of any element of $\mathscr O[G]$, its sign-value
function on $\F_2^b$ expands as
\[
 A(x)=a_\varnothing+
        \sum_{I\ne\varnothing}2^{|I|}a_I\prod_{i\in I}x_i .
\]
Its low bit is constant and its bit of index $j\ge1$ has Boolean
degree at most $2^{j-1}$. One can see this by extracting the
$Y^{2^j}$ coefficient of $(1+Y)^{A(x)}$ modulo two:
the first possible contribution of $x_I$ has degree at least
$2^{|I|}$ in $Y$, and $|I|\le2^{|I|-1}$.
Thus $q$ coefficient equalities modulo $2^M$ have sum of Boolean
degrees at most $q(2^{M-1}-1)$.

If that sum is less than $b$, multiply the complementary indicators
of the equations. Its degree is less than $b$, so its sum over the
Boolean cube is zero: every monomial misses a variable.
The support of this indicator is the simultaneous solution set.
It has even cardinality and contains zero, hence another address.

Apply this to the coefficients of $f$ from $-L_T(M)$ through zero
and the single constant coefficient of $d$. It suffices to choose
\[
 M>\max(c_{\det},B+2c_{\det}),\qquad
 b>(L_T(M)+2)(2^M-1),                                           \tag{5.10}
\]
where $B$ bounds the nonzero-address central valuations.
These choices precede the graph and its theta precision.
At the resulting nonzero address, negative coefficients of $f_x$
vanish modulo $2^M$, and its constant coefficient agrees with
$f_0(0)$ and has valuation $c$.
Positive coefficients may have smaller valuation. Nevertheless
$v_{\mathscr O}(f_x)\le c$, so integrality of
$d_x=f_x\widetilde u_x$ gives
$v_{\mathscr O}(\widetilde u_x)\ge-c$.
This estimate comes from the displayed product at the selected
nonzero address; it does not apply (5.5) to the unknown base.
Both $\widetilde u_x$ and $\widetilde u_0$ lie in $\Lambda[1/2]$
by the height-one argument above, so neither has negative powers
of $t$. In particular, since $f_0\in\Lambda$,
$[t^0]d_0=f_0(0)\widetilde u_0(0)$.
The convergent constant-term convolution yields
\[
 v_2\bigl(\widetilde u_x(0)-\widetilde u_0(0)\bigr)\ge M-2c>B.     \tag{5.11}
\]

By (3.9), (3.4), (4.6) and the unit comparison (4.7), every nonzero address has
$v_2\widetilde u_x(0)\le B$ for one $B$ fixed before $b$; the finite
$\Sha$ terms have negative sign and may be discarded for this upper
bound. If the base derivative vanished, Section~\ref{sec:arithmetic}
would give $\widetilde u_0(0)=0$, contradicting (5.11).
Hence $L'(E/K,1)\ne0$. The complementary value $L(E^{(k)},1)$ is
nonzero, and
$L(E/K,s)=L(E,s)L(E^{(k)},s)$, so $E$ has analytic rank one.
This proves (RC), and hence Theorem~\ref{hyp:RC}, using the classical
applications and finite constructions specified above.

% End 03b-limits.tex


% Begin 04-computation-value.tex
\section{Exact precision computation and polynomial compilation}\label{sec:compiler}
Encode a signed integer by
\[
 \operatorname{enc}(n)=\begin{cases}2n,&n\ge0,\\-2n-1,&n<0.\end{cases}
\]
The five natural inputs decode in the fixed order $(n,\kappa,a,l,j)$.
The decision predicate rejects unless
$n\notin\Cset$, $\kappa\ne0$, $a,l\ge1$, and $j=na$.
For nonzero multiples of $P_0$, $z=-x/y$ is well-defined. Set
\[
 z([l]P_0)=u_l+v_ls,\qquad\epsilon_l=|u_l|+|v_l|,
 \qquad r(a,j)=\tfrac12\Tr_{K_0/\Q}\frac{z([j]P_0)}{z([a]P_0)}.
\]
With the fixed constant $A_{\rm ht}^{(0)}=H_6(S_{\rm ht})$, its final
tests are
\[
 |r(a,j)-n|<\epsilon_l,\qquad
 4\epsilon_l<4^{-A_{\rm ht}^{(0)}|\kappa F(n)|^{24}}.
\]
These are strict tests on exactly computable rationals and integers.

The reference program S21 uses exact quadratic-field arithmetic and capped
integer operations. The ideal register intermediate representation S22
computes the same predicate on every five-natural tuple. Its independent
all-input argument covers normalized signed arithmetic, origin and
zero-ordinate cases, equal-abscissa addition, and every early rejection.
Nine exact polynomial identities verify the nonexceptional projective
formulas and preservation of $Y^2=X^3+AXZ^4$. Scalar loops decrease by
halving; the precision loop divides by four; factorial iteration is bounded.
This is a proof about the ideal intermediate representation, rather than
an assertion that a budgeted simulator always finishes.

For example, if $\epsilon=U/V>0$ with natural $U,V$ and integer $B\ge0$,
\[
 4\epsilon<4^{-B}\quad\Longleftrightarrow\quad
 4^B\le\left\lfloor\frac{V-1}{4U}\right\rfloor.
\]
This permits a finite cap. For positive $L$, $\sigma(x)=\min(x,L+1)$
commutes with saturated addition and multiplication. Powers and factorials
can be computed with this saturation. Capping an exponent is safe for
base at least two because an exponent above $L$ already makes the value
exceed $L$. Base zero, base one and exponent zero are separate cases.
This proves equality with the fixed-height expression without requiring
its full evaluation when the available precision is too small.

\subsection{The natural polynomial and the power interface}
The saved precision polynomial has $1\,305\,811$ natural variables,
the first five being inputs. Its SHA256 is
\begin{center}\small\ttfamily
f3e163f5c60b972c98b48647274f956dd7d78428e27ddc0543574a5f18120b8c.
\end{center}
Source S23 specializes a fixed universal natural halting polynomial at
the code of a wrapper. The wrapper inverts four Cantor pairs and then
runs the unchanged precision machine. Its pairing equations are
\[
 2t_i=(c_i+t_{i+1})(c_i+t_{i+1}+1)+2t_{i+1}\quad(0\le i<4),
 \qquad t_4=c_4.
\]
Their triangular intervals prove unique natural unpairing. The new
polynomial sums the square of the old halting polynomial and the four
squared differences above.

The universal compiler packs a finite run of a fixed $m$-state
increment/decrement/branch program into base-$B$ integers. Acceptance
means reaching the unique final halt with all registers zero; the fixed
interpreter includes a finite cleanup stage. Choose
\[
 c>0,\quad K=2^c,\quad B=4K>m,\quad T=B^q,\quad
 F_q=1+B+\cdots+B^q,\quad D=(K-1)F_q,
\]
and require every input to be below $K$. For natural integers, say
$X$ is masked by $Y$ when every set binary bit of $X$ occurs in $Y$.
Require each register history $R_j$ to be masked by $D$ and satisfy
$R_j<T$. Each state history $S_i$ is masked by $F_q$, and
\[
 \sum_{i=0}^{m-1}S_i=F_q,\qquad S_{\rm halt}=T.
\]
Thus register digits are below $K$, terminal register digits are zero,
and state digits are zero or one. Since $B>m$, the state sum has no
carries and selects exactly one state at each time.

The nonzero mask $N_j$ is determined by
\[
 K N_j=(R_j+D)\mathbin{\mathrm{AND}}(K F_q).
\]
Indeed each digit $r+K-1<2K<B$ has its $K$-bit set exactly when
$r>0$. For a decrement/branch instruction at state $i$ testing
register $j$, impose
\[
 V_i=S_i\mathbin{\mathrm{AND}}N_j,\qquad S_i=V_i+W_i.
\]
These are the nonzero and zero branch histories. Form the packed
increment indicator $A_j$ from states incrementing $j$, the decrement
indicator $C_j$ from the appropriate $V_i$, and each incoming-state
indicator $I_i$ from the states or branches directed to $i$.
They have zero-one digits, since there is just one active instruction
and branch at each time. The transition equations are
\[
 R_j+BC_j=\operatorname{initial}_j+BR_j+BA_j,\qquad
 S_i=\mathbf1_{i=0}+BI_i.
\]
Every register coefficient on either side is at most $K<B$;
uniqueness of base-$B$ digits gives the initial conditions and each
exact transition. The possible extra final slot vanishes because
halt has no outgoing instruction and the terminal registers are zero.
For $q=0$ the same equations require the initial state already to be
halt and the inputs to vanish. Conversely, any finite accepting run
fits after choosing $c$ sufficiently large, and supplies all these
masks and equations.

\subsection{The literal alpha and power formulas}
Define $U_n=\alpha_b(n)$ by
\[
 U_0=0,\qquad U_1=1,\qquad U_{n+2}=bU_{n+1}-U_n\quad(n\ge0).
\]
For $b\ge3$ this recurrence gives
$0\le U_n<U_{n+1}$ and $U_n\ge n$. For $n\ge1$, the matrix
\[
 M=\begin{pmatrix}b&-1\\1&0\end{pmatrix},\qquad
 M^n=\begin{pmatrix}U_{n+1}&-U_n\\U_n&-U_{n-1}\end{pmatrix}
\]
has determinant one. Thus adjacent terms are coprime and
$U_n^2-bU_nU_{n-1}+U_{n-1}^2=1$.
Parameter congruence follows directly from the recurrence:
$b\equiv b'\pmod d$ implies
$\alpha_b(n)\equiv\alpha_{b'}(n)\pmod d$.
In particular $\alpha_2(n)=n$. This is parameter congruence,
not merely congruence between indices.

Every natural solution of $x^2+y^2=1+bxy$, $b>2$, is an unordered
adjacent pair. To prove this, order it as $X>Y$. If $Y=0$, $X=1$.
Otherwise $Z=bY-X$ satisfies $XZ=Y^2-1$, hence $0\le Z<Y$,
and $(Y,Z)$ satisfies the same equation. Descent reaches $(1,0)$;
reversing the steps generates exactly the adjacent pairs.
No ordering assumption on the original pair is needed.

For $k>0$, reduction of $M^k$ modulo $U_k$ is the invertible scalar
$U_{k+1}I$. Dividing an index by $k$ and using strict growth proves
$U_k\mid U_m$ if and only if $k\mid m$. More precisely, for $m>0$,
\[
                    U_k^2\mid U_m\quad\Longleftrightarrow\quad
                         kU_k\mid m.                           \tag{C.1}
\]
For $k=0$ both sides are false. For $k>0$ write $m=kq$ and expand
$M^{kq}=(U_kM-U_{k-1}I)^q$ modulo $U_k^2$. Its lower-left entry is
$qU_k(-U_{k-1})^{q-1}$, so coprimality proves (C.1), including $U_k=1$.

For $m>0$ put $v=U_{m+1}-U_{m-1}=\operatorname{tr}(M^m)$.
Cayley--Hamilton gives $M^{2m}\equiv-I\pmod v$.
Extend the recurrence by $U_{-j}=-U_j$. Every $n\ge0$ has
$n=2lm+\varepsilon j$ with $l\ge0$, $0\le j\le m$ and
$\varepsilon\in\{1,-1\}$; then
$U_n\equiv(-1)^l\varepsilon U_j\pmod v$.
The value sign need not equal the index sign. If $b\ge4$,
$v=bU_m-2U_{m-1}>2U_m$.

The actual alpha formula has natural witnesses $r,s,t,u,v,w,x,y$
and the following fifteen conditions:
\[
\begin{gathered}
 b>3,\quad u^2+t^2=1+but,\quad s^2+r^2=1+bsr,\quad r<s,\quad
 u^2\mid s,\quad v+2r=bs,\\
 w\equiv b\pmod v,\quad w\equiv2\pmod u,\quad w>2,\quad
 x^2+y^2=1+wxy,\\
 2X<u,\quad2X<v,\quad x\equiv X\pmod v,\quad
 2N<u,\quad x\equiv N\pmod u.
\end{gathered}                                                   \tag{C.2}
\]
A strict inequality $a<b$ is $b=a+1+d$; divisibility is $b=ad$;
congruence is $A+Mu=B+Mv$, with fresh natural witnesses.
At zero modulus the last equation says equality, so there is no
exception hidden in this replacement.

For soundness, the unordered Pell result gives $u=U_k$, $k>0$,
and $x=\alpha_w(n)$. The ordered pair gives $s=U_m,r=U_{m-1}$,
$m>0$. By (C.1), $u\mid m$. Parameter congruence gives
$x\equiv U_n\pmod v$ and $x\equiv n\pmod u$.
Reflect $n$ to $j\in[0,m]$. Since $2X<v$ and $2U_j<v$,
the congruence $X\equiv\pm U_j\pmod v$ forces $X=U_j$,
with the negative case forcing both zero.
Then $j\le X$, $2j<u$ and $2N<u$ force $N=j$ from
$N\equiv\pm j\pmod u$. Thus $X=U_N$.

Conversely choose an odd $u=U_k>2U_N$. Such terms are unbounded,
since adjacent terms are coprime and the sequence grows.
Set $t=U_{k+1}$, $m=ku$, $s=U_m$, $r=U_{m-1}$ and $v=bs-2r$.
Equation (C.1) gives $u^2\mid s$, and the required strict bounds hold.
The moduli $u,v$ are coprime: a common divisor is odd, divides $s$
and $2r$, hence divides both adjacent terms. CRT gives
$w\equiv2\pmod u$, $w\equiv b\pmod v$ with $w>2$.
Take $x=\alpha_w(N)$, $y=\alpha_w(N+1)$.
Parameter congruence verifies every remaining condition.
This includes $N=X=0$.

The literal power formula separates $Q=0,R=0,P=1$ and
$Q=0,R>0,P=0$. For $Q\ge1$ it imposes
\[
 a_0=\alpha_{Q+4}(R+1),\quad b=a_0+Q^2+2,\quad
 m+Q^2+1=bQ,\quad P<m,
\]
\[
 a_1=\alpha_b(R),\quad a_2=\alpha_b(R+1),\quad
 P+ba_1\equiv Qa_1+a_2\pmod m.                                  \tag{C.3}
\]
Induction gives $a_0\ge(Q+3)^R$, so $m>Q^R$ and $m>0$.
Let $E_R=Q U_R-bU_R+U_{R+1}$. The recurrence and
$bQ\equiv Q^2+1\pmod m$ give $E_0=1$ and
$E_{R+1}\equiv QE_R\pmod m$. Thus (C.3) says $P\equiv Q^R\pmod m$.
Both values lie in $[0,m)$, forcing equality. Conversely equality
supplies every witness. This includes $R=0$, $Q=1$ and $0^0=1$.

These formulas make the binomial and bit operations polynomial.
For $B=2^{n+1}$, the equation
\[
 (1+B)^n=B^k(Bq+C)+r,\quad C<B,\quad r<B^k
\]
selects exactly $\binom nk$, since all its base-$B$ coefficients
are below $B$, including zero for $k>n$.
The identity
$(1+t)^n=\prod_{i:\,\text{bit}_i(n)=1}(1+t^{2^i})$ over $\F_2$
makes odd binomial coefficients equivalent to bit inclusion.
Bit disjointness of $a,b$ is equivalent to inclusion of $a$ in $a+b$:
at the least shared bit there is no incoming carry and that bit
is cleared. Finally $x=z+a,y=z+b$ with $z,a,b$ pairwise bit-disjoint
says exactly $z=x\mathbin{\rm AND}y$.

\subsection{Interpreter, lowering and five-input specialization}
The universal program numbering is $e+1=2^w(2D+1)$.
At $w=0$ it halts. Otherwise the four base-$2^w$ digits at positions
$4p,\ldots,4p+3$ give the opcode, register and two targets at state $p$.
Opcode one increments, opcode two decrements or takes its zero branch;
every other opcode halts. Missing digits are zero, so decoding is total.

The interpreter stores a finite tape as
$0$ or $\operatorname{cons}(h,t)=2^h(2t+1)$.
Repeated halving decodes a nonempty head and tail; the tail is smaller
than the code. To update a register, save the finite reversed prefix,
replace the entry and restore the prefix. Absent entries are zero.
The reserved jump register remains zero. Copies restore their source;
aliased addition snapshots its operand before overwriting it.
Division maintains $x=qy+r$ and subtracts only when $r\ge y$;
the interpreter's divisors are at least two.
These invariants prove that one interpreter loop executes exactly one
decoded instruction, with finite work. A simulated halt enters cleanup
and clears all registers. An infinite simulation never enters cleanup.
Thus its all-zero halt is equivalent to virtual halting, without
requiring the virtual final tape itself to be zero.

The precision IR uses normalized signed pairs for integer arithmetic
and pairs over the basis $1,\sqrt2$ for field arithmetic.
Norm zero implies the field element is zero: after clearing
denominators, a primitive solution of $a^2=2b^2$ would have both
entries even. Every division is guarded by that exact test.
The Jacobian point formulas satisfy the curve equation and the affine
chord or tangent formulas after multiplying denominators.
Equal abscissas split into doubling and opposite-point branches;
zero ordinate doubles to the identity. The fixed projective point
is the same affine $P_0$. These are polynomial identities, independent
of the generator theorem.

The scalar loop halves its remaining count. The precision loop divides
its remainder by four. The capped arithmetic computes exactly
$\min(x,L+1)$; zero bases, unit bases and zero exponents are separate.
For bases at least two an exponent exceeding $L$ already exceeds the
cap. Thus the ideal IR terminates and agrees with the stated rational
predicate on all five natural inputs, including every early rejection.

Each primitive lowering uses fresh auxiliary registers, snapshots
overlapping sources, and preserves the live-register map at the end
of each finite block. Saturating decrement gives natural subtraction,
and repeated addition gives multiplication. All variable divisions
are guarded. Induction on completed loop iterations proves the same
claim for loops. Acceptance enters cleanup; rejection enters a
permanent self-loop. Hence the primitive machine halts exactly on
acceptance, although the ideal decision program is total.

For $\operatorname{pair}(a,b)=(a+b)(a+b+1)/2+b$, unpairing maintains
$\text{remainder}=\text{input}-s(s+1)/2$ and subtracts $s+1$
while $s<\text{remainder}$. At exit
$0\le\text{remainder}\le s$, so $b=\text{remainder}$, $a=s-b$.
Four repetitions put the five codes into registers zero through four.
Fresh scratch lies above the old register range and is cleared;
old work registers remain zero. Relocating every target by the prefix
length therefore preserves the precision machine.

The polynomial specialization fixes the universal program input to
this wrapper code, maps the old packed input to variable five, and
maps every old variable $i\ge2$ and private-block start by $i\mapsto i+7$.
The four pairing variables occupy five through eight, disjoint from
the five inputs. The appended four squared equations force their
unique Cantor values. Thus, for the exact saved polynomial $C$,
\[
 \forall\mathbf c\in\N^5,\qquad
 \exists\mathbf z\in\N^{1305806}\ C(\mathbf c,\mathbf z)=0
 \quad\Longleftrightarrow\quad
        \operatorname{accepts\_naturals}(\mathbf c)=\mathrm{True}. \tag{C.4}
\]
The natural compiler proof has no height, restricted-converse or
rational-model premise. Its use of the height expression concerns
only computation of that fixed integer expression.

\subsection{Translation to point codes}
Apply (C.4) using the exact signed index model. Induction on the actual
source DAG gives each node the code of its signed integer value.
Natural guards constrain source variables, not intermediate sums or
products. Raw addition receives canonical codes; globally typed
multiplication justifies the natural guards. Fixed generator equations
and scalar templates supply constants.

The power-value recipe has three signed code inputs, seventy-nine
private natural witnesses, and an output code for the value of the
power-template polynomial. That value is not independently required
to vanish. Its saved root equality is $[17751579,15]$.
The precision parent has $54056$ explicit variables and $15845$
disjoint private blocks of seventy-nine variables. No parent variable
node accesses a private coordinate. Replacing each block by its
private power-value call therefore preserves the full existential
assignment. Only the final parent root is equated to the identity.

At each primitive formula, equality is a difference, conjunction is
a sum of squares, and disjunction is a product. A topological induction,
including the globally conjoined fixed equations, proves equality of
polynomial values under substitution, not merely equisatisfiability
of an unrelated formula. This is the required link to the saved AST.

\section{A variable-shift relation for ordinary values}\label{sec:bridge}
We prove the ordinary-value deduction from the exact index model,
the compiled precision relation, the height theorem and $\mathsf{PC}$.
These are supplied by Sections~\ref{sec:model}, \ref{sec:compiler},
\ref{sec:height} and Theorem~\ref{hyp:PC}, respectively. Capital letters below denote five-coordinate point codes;
their lower-case indices are ordinary integers. The external variable
$t$ is any rational number.

\subsection{Approximation for every fixed shift}
Define $\Close(N,\mathcal K,t)$ by choosing $A,L,J$, linking all five
signed codes to precision inputs in the order $(n,\kappa,a,l,j)$, and
imposing $|t-r(a,j)|<\epsilon_l$ on those same coordinates. Put
\[
 B(n,\kappa)=H_6(S_{\rm ht})|\kappa F(n)|^{24}.
\]
The precision contract gives
\begin{align}
 \Close(\iota(n),\iota(\kappa),t)&\Longrightarrow
                  |t-n|<e^{-B(n,\kappa)},\label{eq:close-sound}\\
 \Close(\iota(n),\iota(\kappa),n)&\quad\text{for every }
                  n\notin\Cset,\ \kappa\ne0.\label{eq:close-complete}
\end{align}
The first follows from $|t-n|<2\epsilon_l<4^{-B}/2<e^{-B}$.
For the second, positive multiples of $P_0$ approach the identity
simultaneously in the compact group $E(\mathbf R)\times E(\mathbf R)$.
This follows, for instance, by partitioning the compact closure of the
cyclic subgroup into arbitrarily small neighborhoods and subtracting
two sufficiently close multiples. If the resulting positive differences
were bounded, a fixed nonzero multiple would be the identity; infinite
order excludes this. Choose $l$ first to make the positive $\epsilon_l$
small enough. Since the derivative of $[n]$ at the identity is $n$,
$r(a,na)\to n$ along positive multiples $a$ approaching it. Choose such
an $a$ and put $j=na$. This works for every already fixed nonzero shift
$\kappa$, regardless of its size.

\subsection{The six local tests}
For $N=\iota(n)$, $\mathcal K=\iota(\kappa)$, form the code for
$\kappa n$, choose $H=\iota(h)$, $U=\iota(s_0)$, and put
$V_0=H-\mathcal K$, $V_1=U-\iota(\kappa n)$.
Require $\kappa\ne0$, $H,V_0\ne O$. Define $w(Q)=1/x(Q)$ for
$Q\ne O$, $w(O)=0$. For $t\notin\Cset$ use
\begin{align*}
 D_1&=w(H),&D_2&=w(U),&D_3&=\sigma w(V_0),&D_4&=\sigma w(V_1),\\
 D_5&=z(U)/z(H)-t,&D_6&=\sigma(z(V_1)/z(V_0))-t.
\end{align*}
Here $\sigma$ is conjugation in $K_0$. Choose $\eta,q\in\Q$ and impose
\begin{align*}
 \eta^2&=1,&R_0&=\eta2730^6F(t),&a_0&=1-R_0q^2\ne0,\\
 R&=R_0a_0^2,&C_j&=D_j/F(t),&W&=\prod_{j=1}^6(1+RC_j^2)\ne0,\\
 \Theta&=\Tr_{K_0/\Q}(W^{-1})\ne0,&B_{\rm pole}&=R/\Theta^2,
 &\Phi(B_{\rm pole})&.
\end{align*}
Call this relation $\Test(N,\mathcal K,t)$. All quotients are guarded
polynomially. There is one common denominator $F(t)$, one product $W$,
one trace and one pole instance.

\begin{lemma}[Divisibility]\label{lem:test}
If $\Test(\iota(n),\iota(\kappa),t)$ holds, then
$M_{S_{\rm ht}}(t)\mid\kappa F(n)$.
\end{lemma}
\begin{proof}
Let $p\mid M_{S_{\rm ht}}(t)$. Exactly one $c\in\Cset$ has
$\delta=v_p(t-c)>0$ odd. The contact differences are units, so
$v_p(F(t))=\delta$; the field is unramified at $p$. Put $b=v_p(a_0)$.
If $q=0$, then $b=0$. Otherwise $\delta+2v_p(q)$ is odd and
nonzero, so $b=\min(0,\delta+2v_p(q))\le0$. In either case
$\alpha=v_p(R)=\delta+2b$ is odd. At each place above $p$,
\[
 v(1+RC_j^2)=\min(0,\alpha+2(v(D_j)-\delta)).
\]
Let $\lambda=\max_{v\mid p}v(W)\le0$. Then
$v_p(\Theta)\ge-\lambda\ge0$. The pole input has odd valuation
$\alpha-2v_p(\Theta)$, which is positive by Lemma~\ref{lem:pole}.
Therefore $\alpha>0$ and $\lambda> -\alpha/2$.

At a place attaining $\lambda$, each nonpositive factor valuation is
at least $\lambda$. Consequently
\[
 v(D_j)\ge\delta-3\alpha/4=\delta/4-3b/2>0
 \quad(1\le j\le6).
\]
Integral valuations are at least one, at the same place for all tests.
The first four put $H,U$ in one formal group and $V_0,V_1$ in the
conjugate formal group. At the identity we set both $z(O)=0$ and
$\log(O)=0$; thus the cases $U=O$ and $V_1=O$ obey the same ratio
identities below. Their denominators are nonzero by the guards on
$H$ and $V_0$. For a nonzero point in either formal group,
\[
 \log(Q)=z(Q)(1+e_Q),\qquad v(e_Q)\ge1.
\]
The invariant differential is integral, and integration loses at most
$v_p(r)$ in degree $r$; for $r\ge2$,
$(r-1)v(z)-v_p(r)\ge1$.
The relation $[s_0]H=[h]U$ gives
$\log(U)/\log(H)=s_0/h$. The fifth test makes this integral and
congruent to $t$ modulo $p$. Likewise the sixth gives
$(s_0-\kappa n)/(h-\kappa)\equiv t$.
Multiplying by the integer denominators and subtracting yields
$v_p(\kappa(n-t))\ge1$. Since $t\equiv c$, we get
$p\mid\kappa(n-c)$, hence $p\mid\kappa F(n)$. The squarefree
product divides as asserted. No unit assumption on $\kappa$ was used.
\end{proof}

\begin{lemma}[One suitable shift]\label{lem:test-complete}
For every $n\in\Z\setminus\Cset$ there is
$\kappa\ne0$ with $\Test(\iota(n),\iota(\kappa),n)$.
\end{lemma}
\begin{proof}
Choose the sign so $R_0=\eta2730^6F(n)>0$ and let $S'$ be its finite
prime support. Each sufficiently small local formal neighborhood has
finite index in the compact local group. A positive common multiple
of the quotient orders gives $\kappa$ with $[\kappa]P_0$ in every
chosen neighborhood. Take
\[
 H=[2\kappa]P_0,\quad U=[2\kappa n]P_0,\quad
 V_0=[\kappa]P_0,\quad V_1=[\kappa n]P_0.
\]
The formal parameter and the derivative of $[n]$ make each $C_j$
integral above $S'$ after sufficiently shrinking the neighborhoods.
Infinite order preserves the nonzero guards.

Choose rational $q$ integral at $S'$ and arbitrarily close to
$1/\sqrt{R_0}$ in $\mathbf R$, without equality. Denominators
$R_0m+1$ with nearest integer numerators give such choices, adjusting
by one if equality occurs. At $S'$, $a_0$ is a unit and
$W\equiv1$, $\Theta\equiv2$ modulo $p^{v_p(R_0)}$.
Integral trace justifies this also at the ramified place two.
The pole input has valuation at least four at two and at least six
at the other primes of $S_0$. Outside $S'$, its valuation is
$2v_p(a_0)-2v_p(\Theta)$ and is even, so its reduced denominator is
a square. In both real embeddings, $R$ tends to zero positively,
$W\to1$ and $\Theta\to2$. All guards hold and
$0<B_{\rm pole}<1/2$ for a sufficiently close choice.
Theorem~\ref{hyp:PC} supplies the pole witnesses.
\end{proof}

\subsection{Exact ordinary value}
Define $\Val(N,t)$ by membership of $N$ and the disjunction
\[
 \bigvee_{c\in\Cset}(N=\iota(c)\wedge t=c)
 \quad\vee\quad
 \exists\mathcal K\,
 (\Close(N,\mathcal K,t)\wedge\Test(N,\mathcal K,t)).
\]
The same five coordinates $\mathcal K$ are used in both calls.
In an accepting main branch $n\notin\Cset$, $\kappa\ne0$, and
Lemma~\ref{lem:test} and the height estimate imply
\[
 h(t)\le H_6(S_{\rm ht})|\kappa F(n)|^{24}=B(n,\kappa).
\]
For $t=a/b\ne n$ in reduced form, $|a-nb|\ge1$, hence
\[
 |t-n|\ge1/b\ge e^{-h(t)}\ge e^{-B(n,\kappa)},
\]
contradicting~\eqref{eq:close-sound}. Thus $t=n$.
Conversely choose $\kappa$ by Lemma~\ref{lem:test-complete}, then use
\eqref{eq:close-complete} for that same shift. Private witness blocks
are disjoint, so the two calls can be satisfied simultaneously.
The six remaining integers have exactly their displayed branches.

Therefore $\Val(\iota(n),t)$ is exactly the graph $t=n$.
Existentially quantifying $N$ gives ordinary $\Z$ in $\Q$.
S01 and S16 multiply the exception and main-branch polynomials and
conjoin membership by a sum of squares. The literal substitutions in
Section~\ref{sec:syntax} and Appendix~\ref{app:literal} identify this
polynomial with the saved root $P_*$. This proves the zero-set part of
Lemma~\ref{lem:assembly}.

% End 04-computation-value.tex


% Begin 05-bounds-appendices.tex
\section{The cutoff and polynomial size}\label{sec:size}
The literal constants in S03 define the catalogue product $H_{\rm cat}$,
its maximum $H_{\max}$, $N_s=P_sH_{\rm cat}$ with
$P_s=\prod_{p\le324}p$, the least formal multiple $M_0$, and $c_*=4M_0$.
They take
\[
 R=\max(D,H_{\max},c_*,10^{12}),\quad
 A_{\rm ht}=2\cdot10^{13}\max(1,C_{\rm cm},Q_R^{17}),\quad
 T=(48A_{\rm ht})^{13}.
\]
The pole clearing constant is $\ell_0=12P_s^{2916}H_{\rm cat}^{732}$.
The symbol $F_0=324!$ is distinct from the six-contact polynomial $F(u)$.

\begin{proposition}[Literal cutoff]\label{prop:cutoff}
The fixed prescription satisfies
\[
 R=D,\qquad T=(96\cdot10^{13})^{13}D^{2210G_0^2D^3}
            <2^{\,2^{\,2^{9013}}}.
\]
This is a statement about the prescribed integer, independently of its
semantic role in excluding false multiplication triples.
\end{proposition}
\begin{proof}
We give the bounds of S25. Exact small calculations give
$10^5<W<2^{512}$ and $2^{2000}<F_0<2^{2250}$.
For $q\ge10^5$, the unconditional progression estimate
\cite[Theorem 1.2, (1.10)--(1.11), Corollary 6.17 with $A=1$]{bmor}
applies for $x\ge\exp(0.03\sqrt q(\log q)^3)$, including the possible
exceptional-zero case. At $X=2^{W^3}$ its error is at most
$X/(160\log X)$ and the hypothesis follows from
$0.03\sqrt W(\log W)^3<(\log2)W^3$.
Hence $\vartheta(X;W,a)>X/(2W)$, giving more than
$X/(2W^4)>22$ primes in each unit class. No prime at most 324
belongs to such a class. Thus
\[
 H_{\max}<X,\quad H_{\rm cat}\le X^{22W},\quad
 \log_2N_s<2^{2054}.
\]
For every odd prime dividing $N_s$, take the actual reduction order of
$P_0$ at each place above it, and let their product be $m_0$.
All preceding positive multiples stay in the dyadic formal subgroup,
since $P_0$ starts there and the subgroup is closed under addition.
A nonidentity point there has negative abscissa valuation, so its
reciprocal is defined; the constructor accepts the identity directly.
Thus every earlier search step is defined without a nontorsion premise.
The point has nonsingular reduction also at the bad places over seven.
Each order is at most $p^3$ and there are at most two places, giving
$m_0\le N_s^6$. Multiplication by $N_s^4$ raises the formal depth enough
that the rational coefficients of the reciprocal abscissa are divisible
by $N_s^4$, with denominator prime to $N_s$.
At two use the valuation with uniformizer $\sqrt2$, so $v(2)=2$.
The initial valuations are $v(x)=-6$, $v(y)=-9$, $v(z)=3$,
and $v([2]z)\ge v(z)+2$ at this depth. Therefore the literal search
succeeds by $N_s^4m_0$ and $M_0\le N_s^{10}$.
This uses actual finite orders, not numerical upper bounds incorrectly
treated as multiples of those orders.

Consequently $c_*<2^{2^{2059}}$, whereas
$2^{2^{8008}}<D<2^{2^{9010}}$, so $R=D$.
Next $C_{\rm cm}<2^{50000D^2}$ and
\[
 \log_2Q_D^{17}=170G_0^2D^3\log_2D>50000D^2.
\]
The maximum is therefore $Q_D^{17}$, giving the identity for $T$.
Finally $G_0<2^{37}$, $E_0<D^3$, $\log_2Q_D<D^5$,
$\log_2A_{\rm ht}<D^7$, and $\log_2T<D^8<2^{2^{9013}}$.
\end{proof}

\subsection{Quantitative syntax bounds}
Use coefficient $\ell^1$ norm for integer polynomials and the weighted
pair norm $\|a+bs\|=\|a\|_1+2\|b\|_1$. They are submultiplicative.
The division-polynomial recurrences give
\[
 \deg f_n\le(n+1)^3,\qquad\|f_n\|\le2^{100(n+1)^3}.
\]
For the even recurrence the exponent slack is
$4m^3-24m-27\ge9$ for $m\ge3$; for the odd recurrence it is at least
ten for $m\ge2$. This absorbs the fixed factors and addition.

The following facts have been checked for the actual saved expression:
the absolute values of all literal fixed constants and all occurring
division-polynomial indices are at most $T$; the specified syntax-height
traversal has root height at most 256; each model recipe has at most
$2^{20}$ constraints and each composition recipe has at most $2^{20}$
terms. The check accounts for all 23 templates, seven recipes, private
witness blocks, 173 scalar constant nodes and 50 finite constructors. Its exact
artifact identity and independent replay are recorded in the source map.
Together with the constructor bounds below, this discharges the
quantitative part of $\mathsf B$. The all-input arithmetic interpretation of the same constituents is the
separate substitution induction of Section~\ref{sec:compiler}.

Put $B_*=T+1$. Coordinates of $\iota(n)$ for
$|n|\le T$ have absolute numerator and positive denominator at most $2^{K_*}$, with
$K_*=2^{30}(T+1)^3$. Evaluate the division polynomials at $P_0$ and use
the fixed coordinate formulas; at most eight quadratic-field operations
are needed, each taking an available height bound $H$ to at most $3H^4$.
There are at most $8(T+1)^3$ guard factors. The resulting bounds for
each finite constructor are
\[
 \deg\le B_*^4,\qquad\log_2\max(1,\|\cdot\|_1)\le B_*^8.
\]

Give variables syntax height zero, constants height one, and add one
for each arithmetic operation or finite constructor. A template call
adds one plus the callee's intrinsic height and the maximum argument
height. The substitution inequalities
\[
 \deg P(Q)\le\deg P\max(1,\max_i\deg Q_i),\qquad
 \|P(Q)\|_1\le\|P\|_1\max(1,\max_i\|Q_i\|_1)^{\deg P}
\]
give degree and logarithmic norm at most $B_*^{8h}$ at positive height
$h$, with the bounded sums in the actual recipes treated as in S26.
This general call rule applies to polynomial arguments. In each outer
composition recipe the arguments are coordinate variables, so the call
is only a renaming, even when coordinates repeat. For any such map
$\sigma$,
\[
 \deg P(Y_{\sigma(i)})\le\deg P,\qquad
 \|P(Y_{\sigma(i)})\|_1
 =\sum_\beta\left|\sum_{\alpha\mapsto\beta}c_\alpha\right|
 \le\sum_\alpha|c_\alpha|=\|P\|_1.
\]
Thus these calls retain the callee's majorant height before charging
the sum-of-squares or product connective. The seven recipe heights are
$228,233,235,251,252,254,256$. A rule charging an extra unit for each
of these mere renamings would describe a looser syntax-depth bound,
not this proved degree-and-norm majorant.
For the model wrappers use the safe height bound four for the
generator equations and root differences; the callee contribution
dominates this bound. The literal height-two claim for those equations
in S26 is not used.
This height is a majorant for the compact constructor syntax, not for
an expanded monomial tree. The height-256 bound yields
\[
 \deg P_*\le(T+1)^{2048},\qquad
 |[X^\alpha]P_*|\le2^{(T+1)^{2048}}.
\]

Since $T$ is an integer, $T+1\le U=2^{2^{2^{9013}}}$. Hence
\[
 (T+1)^{2048}\le U^{2048}=2^{2^{2^{9013}+11}}
                     <2^{2^{2^{9014}}}.
\]
As $9014<65536=\tw4$, the degree is less than $\tw7$ and the
coefficient magnitude less than $\tw8$. The root recipe declares
$292\,890\,231\,904$ witnesses. Together with the preceding estimates,
these complete the bounds for the same $P_*$ in
Theorem~\ref{thm:target}.

\section{Consequences of the integer definition}
The proof of Theorem~\ref{thm:target} combines the following results.
\begin{center}\small
\begin{tabular}{p{.20\textwidth}p{.43\textwidth}p{.27\textwidth}}
\toprule
Result & Mathematical inputs & Proof location\\\midrule
Index group & Two-isogeny descent, reductions, finite height search &
Section~\ref{sec:model}\\
Six-contact bound & Quaternionic model, tame pullback, compatible-isogeny
height theorem, theta comparison & Section~\ref{sec:height}\\
Pole completeness & Simultaneous descent, restricted converse, forward
rank-one theorem & Section~\ref{sec:pole}\\
Multiplication & Index group, height bound, selected primes,
pole completeness, finite guard & Section~\ref{sec:mult}\\
Precision relation & Decision program, lowering, power and packed histories &
Section~\ref{sec:compiler}\\
Ordinary value & Shared-shift Close and Test, height and rational separation &
Section~\ref{sec:bridge}\\
Bounded $P_*$ & Literal substitutions and quantitative syntax &
Sections~\ref{sec:syntax}, \ref{sec:size}; Appendix~\ref{app:literal}\\\bottomrule
\end{tabular}
\end{center}

\subsection{Three corollaries}\label{sec:four-corollaries}
We apply the exact equality between the rational zero projection of
$P_*$ and ordinary $\Z$ proved in Theorem~\ref{thm:target}.

\begin{corollary}[Negation of Mazur's finite-component conjecture]
\label{cor:mazur}
Mazur's
finite-connected-components conjecture for the real-topological closure
of rational points is false.
\end{corollary}
\begin{proof}
Mazur's Conjecture~3 in \cite[p.~165]{mazur1995} asserts that the closure
of the rational points of any variety over $\Q$ in its real locus has
finitely many connected components. In that source a variety means a
reduced scheme of finite type; irreducibility is not required.
Put $r=292\,890\,231\,904$ and
\[
 V=\bigl(\{P_*(t,w_1,\ldots,w_r)=0\}\subset\mathbf A^{r+1}_{\Q}\bigr)_{\rm red},
 \qquad C=\overline{V(\Q)}\subset V(\mathbf R).
\]
Reduction changes neither the rational nor the real points. Let
$\pi:V(\mathbf R)\to\mathbf R$ be projection to $t$.
Theorem~\ref{thm:target} gives $\pi(V(\Q))=\Z$. Since $\Z$ is closed in $\mathbf R$
and $\pi$ is continuous, $C\subset\pi^{-1}(\Z)$, so $\pi(C)\subset\Z$.
Every integer already occurs on $V(\Q)$, hence $\pi(C)=\Z$.
The image of any connected component of $C$ is a connected subset of
$\Z$, therefore a singleton. Points projecting to distinct integers
belong to distinct components. Thus $C$ has infinitely many connected
components, contradicting the stated conjecture.
\end{proof}

\begin{corollary}[A universal rational Diophantine polynomial]
\label{cor:rational-universal}
There exist $m\ge1$ and a fixed polynomial
$U\in\Z[e,n,x_1,\ldots,x_m]$ such that, for every $e,n\in\N$,
\[
 \exists\mathbf x\in\Q^m\ U(e,n,\mathbf x)=0
 \quad\Longleftrightarrow\quad
 \text{machine $e$ halts on input $n$}.
\]
Here $e$ and $n$ are ordinary integer parameters, with $\N$ embedded
in $\Q$, and are not elliptic-curve coordinates. In particular, every
computably enumerable subset of $\N$ is represented by a suitable
fixed value of $e$ in this one polynomial.
\end{corollary}
\begin{proof}
Fix an effective enumeration of machines. The Davis--Putnam--Robinson--Matiyasevich theorem, \cite[Basic Theorem, p.~3]{matiyasevich1971},
applied to the universal halting relation gives a fixed
$D\in\Z[e,n,z_1,\ldots,z_s]$, with $s\ge1$, such that, for $e,n\in\N$,
\[
 \text{machine $e$ halts on $n$}
 \quad\Longleftrightarrow\quad
 \exists\mathbf z\in\Z^s\ D(e,n,\mathbf z)=0.
\]
To use integer witnesses in this statement, replace each nonnegative
witness by a sum of four integer squares. In the positive-integer
convention of the cited source, apply the theorem to the predicate
that machine $a-1$ halts on $b-1$, substitute $a=e+1$, $b=n+1$,
and use one plus that sum for each witness; see also
\cite[p.~2]{matiyasevich1971}. A dummy witness ensures $s\ge1$ if needed.
For each $z_i$ take a separate block $\mathbf w_i$ of $r$ rational
variables, where $r$ is the witness count of $P_*$. Define
\[
 U(e,n,\mathbf z,\mathbf w_1,\ldots,\mathbf w_s)
   =D(e,n,\mathbf z)^2+\sum_{i=1}^s P_*(z_i,\mathbf w_i)^2.
\]
A rational zero makes every summand zero. Theorem~\ref{thm:target}
then forces every $z_i$ to be an integer and yields the displayed
halting condition. Conversely, integer witnesses for $D$ have separate
rational witness blocks for $P_*$, giving a zero of $U$. All blocks
are private and the polynomial is fixed independently of $e,n$.
No variable, degree or coefficient bound proved for $P_*$ is asserted
for this new polynomial $U$.
\end{proof}

\begin{corollary}[The many-one degree of $\operatorname{HTP}(\Q)$]
\label{cor:htp-degree}
\[
 \operatorname{HTP}(\Q)\equiv_m 0',
 \qquad 0'=\{e\in\N:\text{machine $e$ halts on input $e$}\}.
\]
Here $\operatorname{HTP}(\Q)$ is the set of codes of integer-coefficient
polynomials, of arbitrary finite arity, having a rational zero, under a
fixed effective coding.
\end{corollary}
\begin{proof}
Given $e$, specialize the fixed polynomial of
Corollary~\ref{cor:rational-universal} at its ordinary integer parameters
to obtain $U(e,e,\mathbf x)$. Computing its coefficient list and code is
effective, and
\[
 e\in0'\quad\Longleftrightarrow\quad
 \ulcorner U(e,e,\mathbf x)\urcorner\in\operatorname{HTP}(\Q).
\]
This proves $0'\le_m\operatorname{HTP}(\Q)$.

Conversely, from the code of a polynomial $f$ construct a machine that
ignores its input and enumerates all rational tuples of the appropriate
arity. It evaluates $f$ using exact rational arithmetic and halts as soon
as it finds a zero. Its index $j(f)$ is computable from the finite code
of $f$, and
\[
 f\in\operatorname{HTP}(\Q)\quad\Longleftrightarrow\quad j(f)\in0',
\]
where the left side uses the code of $f$. Thus
$\operatorname{HTP}(\Q)\le_m0'$. Invalid polynomial codes, if present in
the coding convention, are sent to a fixed nonhalting index. Using
rational-coefficient polynomial codes gives the same result, since
clearing their finitely many denominators is effective and preserves
rational zeros. The two reductions prove the assertion.
\end{proof}

More generally, an integer Diophantine problem translates to a rational
one by adjoining $P_*$ for every integer witness, with private witness
blocks. Natural witnesses can first be expressed as sums of four integer
squares. Bounds for the resulting polynomial require their own composition
calculation; the bounds for $P_*$ do not automatically transfer.

For example, a machine that enumerates ZFC proofs and halts upon finding
a contradiction has a rational Diophantine halting condition by
Corollary~\ref{cor:rational-universal}. This equivalence does not assert
that a contradiction exists, and no additional explicit polynomial is
specified here.

\appendix
\section{Dependency order in the arithmetic argument}\label{app:obligations}
The even coefficient bound uses stationary-base selection as well as
integrality. The odd minimum uses a fixed nonzero witness and the prior
genus-index lower bound. Support removal is applied only to the actual
unramified specialization. The functional bound and determinant clearing
are proved separately, and the same corrected coordinate is used at every
address. These constructions are given in Sections~\ref{sec:rc}--\ref{sec:limits}
and Appendices~\ref{app:even-source}--\ref{app:endpoints}.

The conversion to ordinary integers uses the canonical model and total
multiplication, followed by the exact computation relation. Only the
shared-shift separation argument in Section~\ref{sec:bridge} identifies
the external rational variable with the ordinary value of its integer
index. The literal substitution and size calculation then apply to the
same expression $P_*$.

\section{Source identities and scope}\label{app:sources}
The companion \texttt{SOURCE\_MAP.json} records exact paths, hashes and
locators, distinguishing mathematical and construction sources, artifacts,
and version records. The base collection contains 83 sources; parity, construction
and uniformity supplements are in separate addenda. The principal local
sources are listed below; full paths
are in the source map.
\begin{longtable}{p{.07\textwidth}p{.58\textwidth}p{.25\textwidth}}
\toprule ID & Source basename & Use\\\midrule\endhead
S01 & \texttt{ordinary-integer-polynomial.js.txt} & Outer construction\\
S02 & \texttt{rational-circuit-translator.js.txt} & Substitution semantics\\
S03 & \texttt{finite-polynomial-expressions.js.txt} & Fixed constructors\\
S04 & \texttt{variable-shift-polynomial-library.js.txt} & Close and Test\\
S05 & \texttt{elliptic-generator-proof.txt} & Generator certificate\\
S06 & \texttt{explicit-elliptic-data.txt} & Rank, torsion, codes\\
S07 & \texttt{verify-elliptic-generator.js.txt} & Exact finite algorithm\\
S08 & \texttt{six-contact-height-proof.txt} & Contacts and constants\\
S09 & \texttt{z-in-q-height-independent-audit.txt} & Height proof\\
S10 & \texttt{explicit-pole-predicate.txt} & Pole polynomial\\
S11 & \texttt{z-in-q-pole-completeness-audit.txt} & Simultaneous descent\\
S12 & \texttt{parity-frozen-4eaee875.tex} & Raw local matrices\\
S13 & \nolinkurl{multiplication-certificate-correctness-proof.txt} & Multiplication proof\\
S14 & \texttt{twenty-two-prime-pattern-proof.txt} & Selected primes\\
S15 & \texttt{z-in-q-variable-shift-bridge.txt} & Ordinary-value graph\\
S16 & \texttt{z-in-q-ordinary-integer-assembly-proof.txt} & Outer composition\\
S17 & \texttt{z-in-q-restricted-converse-assembly.txt} & RC assembly inputs\\
S18 & \texttt{ring-limits-frozen-4eaee875.tex} & Finite models and limits\\
S19 & \texttt{paired-functional-and-uniformity-proof.txt} & Paired functional\\
S20 & \texttt{determinant-transfer-audit.txt} & Clearing and transfer\\
S21 & \texttt{variable-shift-precision.py.txt} & Exact decision program\\
S22 & \texttt{precision-register.js.txt} & Ideal IR and lowering\\
S23 & \texttt{specialize-precision.js.txt} & Wrapper specialization\\
S24 & \texttt{integer-halting-correctness-proof.txt} & Natural compiler\\
S25 & \texttt{z-in-q-explicit-cutoff-bound.txt} & Literal cutoff\\
S26 & \texttt{z-in-q-polynomial-size-bound-proof.txt} & Syntax/norm argument\\
S27 & \texttt{z-in-q-numerical-polynomial-bounds.txt} & Tower comparison\\
S28 & \texttt{determinant-switch-audit.txt} & Switches and outer rank\\
S29 & \texttt{ring-determinants-frozen-4eaee875.tex} & Ring data and derivatives\\
S30 & \texttt{preliminaries-frozen-4eaee875.tex} & Selmer coranks and lengths\\
S31 & \texttt{fixed-certificate-constants.js.txt} & Exact least-multiple search\\
S32 & \texttt{pointwise-frozen-4eaee875.tex} & Distinct point indices\\
\bottomrule
\end{longtable}
The supplement manifest also records the revised fixed-degree CM proof,
the complete natural compiler proof, actual arithmetic and analytic
producer arguments, finite-Artin and determinant-count clarifications,
the fixed BFH witness application, exact substitution checks, and the
new fixed-form, weighted-theta, operator and normalization derivations.
Primary theorem applications remain explicit in the mathematical text.
The version records identify the source files used in the construction.

% End 05-bounds-appendices.tex


% Begin 06-literal-construction.tex
\section{Literal finite-constructor specification}\label{app:literal}
The supplementary file \nolinkurl{ordinary-integer-definition.json}
is part of this specification. It is a finite directed acyclic description,
not an expanded monomial list. The following rules determine one polynomial
from its bytes. Diagnostic labels and intended-theorem strings add no equations.

The root is \texttt{integers}. Variable zero is the external input and
variables one through $292890231904$ are its rational witnesses.
A nonnegative point reference $r$ means $X_r,\ldots,X_{r+4}$.
References $-1,-2,-3$ respectively mean
\[
 (1,0,0,0,0),\quad(0,g_0,-145,g_1,g_2),\quad
                      (0,g_0,-145,-g_1,-g_2).
\]
Every model recipe conjoins
$8g_0-1641=0$, $4g_1+15465=0$, $32g_2-87483=0$.

\subsection{Scalar nodes and fixed expressions}
Read nodes in their stored topological order.
A \texttt{const} is its decimal integer; \texttt{var} is its local variable;
\texttt{add} and \texttt{mul} are ordinary polynomial operations.
For a \texttt{call} to a template with $p$ parameters, substitute the
specified node polynomials for variables $0,\ldots,p-1$, and $X_{w+j}$
for variable $p+j$, where $w$ is the witness start. The value is the
callee's entire root polynomial.

A \texttt{fixed\_integer} evaluates a finite expression using literals,
references to preceding definitions, addition, multiplication,
nonnegative powers, maximum, factorial, the catalogue product or maximum,
and the least formal multiple. The catalogue uses the first 22 primes
in increasing order in each unit class. The searches are deterministic
and terminate by Section~\ref{sec:size}. They produce constants, not
quantified witnesses.

A \texttt{psi\_component} uses the division-polynomial recurrence in
Section~\ref{sec:syntax}, pair multiplication
$(a,b)(c,d)=(ac+2bd,ad+bc)$ and $A=(-11243,7950)$.
It returns component zero or one; every index is fixed.
The positive fixed point codes use
\[
 x_n=x(P_0)-\frac{\psi_{n+1}(P_0)\psi_{n-1}(P_0)}{\psi_n(P_0)^2},
 \qquad y_n=\frac{\psi_{2n}(P_0)}{2\psi_n(P_0)^4}.
\]
Nontorsion makes these fixed denominators nonzero.
Negative indices negate both ordinate components.

A \texttt{finite\_guard\_product} ranges over
$(a,b,c)\in[-T,T]^3\cap\Z^3$ with $ab\ne c$, with fixed $T$.
Write the fifteen coordinates of the three codes as reduced
$A_j/B_j$, $B_j>0$, and put $D=\prod_jB_j$. Its factor is
\[
                 \sum_{j=1}^{15}(DX_j-(D/B_j)A_j)^2.
\]
Multiply these factors, with empty product one. This is an integer
polynomial with precisely the finite false-triple zero set used above.

\subsection{Recipe constraints and witness scope}
A model-evaluation recipe sums squares of the generator equations,
listed natural calls, scalar-constant calls, add/multiply/power-value calls,
and the five coordinate differences of \texttt{rootEquality}.
A source-variable row with a non-null natural witness start contributes
exactly that natural call. Operation kinds zero, one and two select
addition, multiplication and power value. Source-node and provenance
fields add no equations.

A sum-of-squares composition sums squares of its callee polynomials;
a product composition multiplies them. Parameter \texttt{[var,i]} means
$X_i$. Each private block uses its recorded start and length.
The checked ranges are disjoint from parameters and other private blocks.
These rules include globally conjoined fixed equations in active callees.

The following parameter and witness counts are extracted from the exact
specification, without expanding its polynomial.
\begin{longtable}{p{.47\textwidth}rrr}
\toprule Template & Parameters & Witnesses & Root node\\\midrule\endhead
\nolinkurl{member} & 5 & 7 & 149\\
\nolinkurl{raw_add} & 15 & 6 & 308\\
\nolinkurl{raw_z} & 7 & 1 & 68\\
\nolinkurl{raw_inverse_x} & 7 & 0 & 54\\
\nolinkurl{log25} & 4 & 0 & 229\\
\nolinkurl{pole} & 1 & 11 & 106\\
\nolinkurl{test_group} & 44 & 162 & 515\\
\nolinkurl{slope} & 22 & 1,298 & 2681\\
\nolinkurl{height_certificate} & 36 & 35,106 & 260\\
\nolinkurl{multiply.scalar} & 10 & 1 & 136\\
\nolinkurl{multiply.guard} & 15 & 1 & 22\\
\nolinkurl{multiply} & 15 & 35,134 & 83\\
\nolinkurl{natural} & 5 & 140,604 & 79\\
\nolinkurl{value.signed_encode} & 10 & 140,654 & 138\\
\nolinkurl{value.abs} & 2 & 4 & 23\\
\nolinkurl{value.close_coordinates} & 16 & 53 & 271\\
\nolinkurl{value.test} & 11 & 35,282 & 775\\
\nolinkurl{value.exceptions} & 6 & 406 & 487\\
\nolinkurl{scalar.integer_constant_0} & 10 & 1 & 136\\
\nolinkurl{scalar.integer_constant_1} & 10 & 1 & 136\\
\nolinkurl{scalar.integer_constant_2} & 10 & 1 & 136\\
\nolinkurl{scalar.integer_constant_3} & 10 & 1 & 136\\
\nolinkurl{scalar.integer_constant_4} & 10 & 1 & 136\\
\bottomrule\end{longtable}

\begin{longtable}{lrrr}
\toprule Recipe & Parameters & Witnesses & Total variables\\\midrule\endhead
\nolinkurl{power_value} & 20 & 17,786,698 & 17,786,718\\
\nolinkurl{precision} & 25 & 292,889,492,836 & 292,889,492,861\\
\nolinkurl{close} & 11 & 292,890,196,199 & 292,890,196,210\\
\nolinkurl{main} & 6 & 292,890,231,486 & 292,890,231,492\\
\nolinkurl{value_union} & 6 & 292,890,231,892 & 292,890,231,898\\
\nolinkurl{value} & 6 & 292,890,231,899 & 292,890,231,905\\
\nolinkurl{integers} & 1 & 292,890,231,904 & 292,890,231,905\\
\bottomrule\end{longtable}


The five scalar-constant templates use the values in
\texttt{library.fixedConstants}. Large program literals are retained
byte-for-byte in the supplementary specification. Replacing them by
unspecified sufficiently large constants would change this definition.

\section{Primary theorem applications}
\begin{longtable}{p{.26\textwidth}p{.65\textwidth}}
\toprule Input & Exact application boundary\\\midrule\endhead
Selmer parity \cite{dd} & Rational two-isogeny, $p=2$, full Selmer corank;
no finite-$\Sha$ hypothesis.\\
Kato \cite{kato} & Fixed weight-two form and symbol, integral transfer and
fixed-valuation smoothing; rational divisibility is Theorem 13.4(2),
with oddness, rational irreducibility and the unipotent condition.\\
BFH and Hoffstein--Luo \cite{bfh,hl} & Actual signed-prime twist with
trivial character and arbitrary finite splitting set. Only the
nonvanishing clause has the four-prime bound before the extra prime.\\
Waldspurger and Ramsey \cite{waldspurger,ramsey} & Full good-Hecke
eigenspace, same dyadic squareclass, and algebraic vectors by flat
field extension. No arbitrary-ring base change in weight $3/2$.\\
CM norms and heights \cite{ccl,cst,cornutjetchev} & Split level,
conductors coprime to level and the new prime, no extra units,
unaveraged sums and one fixed parametrization.\\
Duality \cite{milne,nekovar} & Actual derived Kummer conditions and their
homotopies. Support change only after the coefficient representation
is unramified at the removed place.\\
Forward finiteness \cite{gz,kolyvagin} & Only at known analytic
nonvanishing/simple zeros or a known nontorsion Heegner trace.\\
Isogeny ratio \cite[proof of I.7.3]{milne} & Finite $\Sha$, dual
sum/difference maps, canonical heights and pullback measures;
no BSD leading-term assumption.\\
Height \cite{vkk,pazuki} & Compatible isogenies or same-field
maximal-order replacement with bounded degree.\\
Primes \cite{greentao,moebius,gtz,bmor} & Fixed finite-complexity system
after fixing sieve cutoff, and the unconditional explicit progression
estimate.\\\bottomrule
\end{longtable}
The source package includes the specification, natural precision polynomial
and construction sources in \nolinkurl{SUPPLEMENT_MANIFEST.json}.
The larger proof corpus is identified by hashes and stable locators in
\nolinkurl{SOURCE_MAP.json}. A source identity is not a proof verdict.

% End 06-literal-construction.tex


% Begin 07-even-source.tex
\section{Fixed-form arithmetic and the even coefficient source}\label{app:even-source}
This appendix supplies the constructions used in (3.1)--(3.3).
Throughout, $E=E_\ell$ is fixed, with the lattice, parametrization and
period convention of Section~\ref{sec:rc}. Constants may depend on these
data and on the prescribed finite list of local types. They do not depend
on the number or sizes of the later good conductor primes.

\subsection{One retained class, including discriminant \texorpdfstring{$4h$}{4h}}
Fix the weight-two newform of $E$ and an integral modular-symbol
combination $\xi$ whose projections on both period lines are nonzero.
Such a combination exists because the symbol span is not the union of
the kernels of these two nonzero rational functionals
\cite[Sections 13.6--13.7]{kato}. A fixed lattice comparison sends its
cohomological lattice into $T_2E$ after multiplication by $2^{c_E}$.
This is a comparison for the original form, before any character or
trace level is chosen.

Kato's classes are formed in $j_*$ cohomology at finite coefficients.
The Leray edge map identifies this group with the classes whose
restrictions to inertia away from two vanish. Restricting a
corestriction is a Mackey sum of transfers of upper inertia
restrictions. Each summand is zero for the retained class.
Consequently even a ramified transfer preserves the integral condition.
Tensoring by a character line is performed over a field on which that
character is trivial. There is no division by an inertia index.

For a finite binary family with group $G$, use the common odd
modulus $F$ and retained-smoothed class $Z_F$ defined in
Section~\ref{sec:rc}. The coefficient map is the integral map on group bases
\[
 [a]\longmapsto\chi_0(a)[a|_\Gamma][\psi(a)].
\]
Inverse-action Shapiro gives the quadratic twist and inverse cyclotomic
action in this convention. In particular this is not an averaged
character projector. At a constituent $\chi$ with odd conductor
$d\mid F$, the actual classes
\[
 y_m=\operatorname{Cor}_{\Q(\mu_{dm})/\Q(\mu_m)}
                       (z_{dm}\otimes e_\chi)
\]
satisfy the twisted norm relation. Transitivity of the same transfers
gives
\[
 (Z_F)_\chi=\prod_{q\mid F/d}
 \left(1-\frac{a_q\chi(q)}q\sigma_q^{-1}
                  +\frac1q\sigma_q^{-2}\right)y_{2^\infty}.
 \tag{E.1}
\]
The factors are present in the classes before any divisibility theorem
is applied.

For an odd signed squarefree parameter $h$, write its fundamental
discriminant as $\Delta_h$ and
\[
 |\Delta_h|=2^a d,\qquad \chi_h=\chi_2\chi_{\rm odd}.
\]
Here $a=0$ or $2$; allowing the finite list $a\le3$ does not change the
argument. At cyclotomic levels $n\ge\max(a,2)$ use the actual map
\[
 [a_0]\longmapsto\chi_2(a_0)\chi_{\rm odd}(a_0)[a_0|_{\Gamma_n}].
 \tag{E.2}
\]
Its coefficients are signs. A real-sign quotient is integral. If the
dyadic character restricts nontrivially to $\Gamma$, the resulting
change is the integral automorphism
\[
                      1+t\longmapsto-(1+t),\qquad t\longmapsto-2-t.
\]
It must be retained in the module action. To compare with the Euler
system, define $y_m$ using $|\Delta_h|m$ in place of $dm$.
Kato's norm factor is a product over newly introduced primes
\cite[Proposition 8.12]{kato}. Increasing an already present power of
two therefore contributes no factor. Factoring the transfer at
$m=2^n$ through $\Q(\mu_{d2^n})$ proves that this $y_{2^n}$ is exactly
the evaluation of (E.2) on a cofinal tail. This proves the
discriminant-$4h$ extension on actual classes.

Choose smoothing integers greater than one with
\[
 c,d_s\equiv1\pmod{3N_{\rm odd}F},\qquad
 c,d_s\equiv1+2^e\pmod{2^{e+1}},
\]
where $e\ge\max\{3,v_2(N)\}$ is fixed by the old dyadic data. Their sizes may depend on
$F$, but the central weight-two, degree-one smoothing factor has valuation
exactly $2e$ \cite[Sections 4.2.4, 6.6 and 9.7]{kato}.
For the signed fundamental discriminant, the quadratic Gauss identity is
\[
 \tau(\chi_h)^2=\chi_h(-1)|\Delta_h|=\Delta_h.
\]
Its dyadic valuation is $a/2$; the
ratio $\sqrt{|\Delta_h|}/\sqrt{|h|}$ also belongs to this finite dyadic
list. Neither contributes a term growing with $\omega(h)$.

\subsection{Rational divisibility and the integral determinant}
Apply Kato's Theorem 13.4(2), not its stronger part (3).
The coefficient prime two is allowed in part (2), which concerns
height-one primes away from two. Oddness gives two one-dimensional
real eigenspaces. Purity gives the required good Frobenius weight.
Non-CM open image, intersected with a quadratic kernel, gives rational
irreducibility and a nontrivial unipotent with one-dimensional fixed
space. The fixed $\xi$ has nonzero period projections. Explicit
reciprocity and Theorem 13.5(2) give infinitely many nonzero cyclotomic
specializations in each sign, hence generic nonvanishing.
No residual irreducibility at two is used.

This theorem initially controls $H^2_0$, the kernel of localization at
two. Choose a fixed integral factor $F_E$ to pay for local Iwasawa
$H^2$ at two and the finitely many bad-prime determinants.
Its central value is nonzero: local control and
$H^2(\Q_v,V_h)=0$, by local duality and finiteness of local elliptic
torsion, exclude a factor $t$. Take the product over the finite allowed
local and dyadic types before varying $h$.

At an active good prime rational inertia is $-1$, so its rational
$j_*$-to-full cone vanishes. At an unused good prime this cone is
\[
 [\,H^1(I_q,V_h)\otimes A
       \xrightarrow{\,1-\sigma_q^{-1}\operatorname{Fr}_q\,}
                 H^1(I_q,V_h)\otimes A\,],\qquad A=\Lambda[1/2],
\]
in degrees one and two. The map is injective, including after
localization at a prime dividing its determinant. The added
degree-two length is the valuation of the Euler factor already in
(E.1). Thus divisibility for the retained $F_EZ_F$ is the required
divisibility for full global cohomology. At a bad prime one uses
$H^1(I_q,V_h)$ itself, not an unjustified replacement by $V_h^{I_q}(-1)$.

Let $C^+$ be the fiber of ordinary global cochains to ordinary real
cochains. Real $H^0$ is included. Over a finite Artin quotient,
continuous cochains are filtered unions of finite free function
modules, hence flat and therefore free. A residual splitting into
finite cohomology and disks lifts by the nilpotent perturbation
series to a finite free model with actual homotopies. Minimal models
at successive quotients are related by invertible matrices:
a homotopy equivalence reduces to an isomorphism of minimal residual
complexes, and Nakayama applies. Lifting these matrices and comparison
homotopies makes the models compatible. Their acyclic, termwise
surjective comparison cones have acyclic inverse limit.
This proves the needed perfection and coefficient change.

For clarity, the residual dimension calculation has two separate
arithmetic inputs. In the real cyclotomic tower
$K_n=\Q(\zeta_{2^{n+2}})^+$ the ordinary two-class group vanishes.
In a quadratic step only the prime above two ramifies. A possible
Hilbert two-class extension would give a finite two-group whose
abelianization is generated by inertia of order two. Its Frattini
quotient has one generator, so the group is cyclic and equals that
inertia. This excludes the extension. Separately the unit
$1+\zeta+\zeta^{-1}$ has norm $-1$; its signature has augmentation
one in $\F_2[C_{2^n}]$, hence is a unit. Its conjugates span all
signatures. The base uses $-1$. Kummer theory and the Brauer invariant
relation then give, for a scalar coefficient,
\[
 \dim H^1(C_n^+)=2^n+s_n-1,\quad
 \dim H^2(C_n^+)=s_n-1,\quad H^i=0\ (i\ne1,2).
\]
For each fixed support $s_n$ is bounded, since a fixed odd prime has
open decomposition group in the cyclotomic tower. Universal
coefficients over $\F_2[[t]]$ leave generic rank one only in degree
one. The two trivial composition factors of $E[2]$ leave rank two.
Minimal cancellation over $B=\mathscr O[G]$ therefore gives
\[
                        C^+\otimes B\simeq B^2[-1].
 \tag{E.3}
\]
This uses the local ring $B$, not semisimplicity of $\F_2[G]$.

Choose a generator of the fixed real invariant lattice and its
connecting class $a$. Since real $H^1$ is killed by two,
$2^{c_E+1}F_EZ_F$ has an integral positive lift $b$.
For a basis $\theta$ of the original inverse determinant line, define
\[
                         u^{\rm ev}\theta=2^{-c_E}a\wedge b.
 \tag{E.4}
\]
Equation (E.3) proves $2^{c_E}u^{\rm ev}\in\mathscr O[G]$.
Rational divisibility gives $u^{\rm ev}\in A[G]$.
Fourier inversion is used only for this latter membership, with no
bound asserted for its denominator. In the original group basis,
\[
              A\cap\mathscr O=\Lambda,\qquad
                  2^{c_E}u^{\rm ev}\in\Lambda[G].                 \tag{E.5}
\]
Indeed an element of $A$ has bounded powers of two as denominators
in its power-series coefficients; membership in $\mathscr O$ makes
their minimum valuation nonnegative. The coefficientwise intersection
retains $c_E$, without a term $v_2|G|$.

\subsection{Central saturation and the numerical lower bound}
At a corank-zero constituent put
\[
 g=\ell_{\Z_2}E^h(\Q)[2^\infty],\quad
 g_v=\ell_{\Z_2}E^h(\Q_v)[2^\infty],\quad
 r_i=\ell_{\Z_2}H^i(\mathbf R,T_h)\quad(i=1,2).
\]
Let $i$ be the index length of the image of global $H^1$ in the
free rank-one quotient $L_2=H^1(\Q_2,T_h)/H_f^1(\Q_2,T_h)$.
Orthogonal Kummer duality, with its real terms, gives
\[
 \ell H^2_{\rm glob}=s(h)-i+\sum_v g_v+r_2-g.                     \tag{E.6}
\]
Write $K=\ker(H^1_{\rm glob}\to H^1_{\mathbf R})$, let $k$ be its
torsion length and $m$ its free index. If $P=H^1(C_h^+)$, $p_0$ is
its torsion length and $a_s$ the saturation index of the real image
in $P/P_{\rm tors}$, exactness gives
\[
 \ell\operatorname{coker}(H^1_{\rm glob}\to H^1_{\mathbf R})
       =r_1-g+k-m,\qquad k=p_0+a_s.
\]
For the numerical calculation suppose also $L_*(h)\ne0$.
The wedge in (E.4) adds $a_s+1+v_{\rm free}(z)-m$.
Conversion to the inverse determinant adds $p_0-\ell H^2(C_h^+)$.
Substitution of these equalities cancels $m$ and $a_s$ exactly:
\[
 v_2u_h^{\rm ev}(0)
 =v_2\exp^*z-s(h)-\sum_v g_v+2g+1-r_1-e_2,                      \tag{E.7}
\]
where $e_2$ is the valuation of $\exp^*$ on a generator of $L_2$.
Primitivity of the real image after inverting $t$ has not been
specialized to $t=0$.

Explicit reciprocity for the single unnormalized character sum gives
\[
 v_2\exp^*z
 =v_2L_*(h)+\sum_{q\ {\rm unused}}v_2\#E^h(\F_q)+\epsilon_{\rm rec}.
\]
At unused good primes $g_q=v_2\#E^h(\F_q)$, so these terms cancel.
At an active good prime inertia invariants in $E^h[2^\infty]$ are
$E[2]$, of length two for our curve. All remaining local terms have
fixed bounds, including global torsion by its injection at two.
Thus
\[
 v_2u_h^{\rm ev}(0)=v_2L_*(h)-2\omega(h)-s(h)+\epsilon_h,\qquad
 |\epsilon_h|\le C_{\rm loc}.                                  \tag{E.8}
\]
The central rational positive complex has only rank-two $H^1$, so
this is specialization of (E.4). Corank-zero duality also shows
$u_h^{\rm ev}(0)=0$ if and only if $L_*(h)=0$.
At $L_*(h)\ne0$, the forward rank-zero theorem supplies the
corank-zero and finite-$\Sha$ hypotheses just used.
Equations (E.5) and (E.8), already with $G=1$, yield
\[
          v_2L_*(h)\ge2\omega(h)+s(h)-c_E-C_{\rm loc}.             \tag{E.9}
\]
No binary-family dimension is needed for this lower bound.

\subsection{The actual even form and its phase}
Fix a sign $\sigma$ and unit classes at a finite support $S_0$
containing two and every prime dividing $N$. Let $D_0>0$ be squarefree
in these prescribed classes, with $L_*(\sigma D_0)\ne0$.
In (E.10)--(E.11), $D$ ranges over positive squarefree parameters in
these same unit classes at every prime of $S_0$.
Fix also a finite extra exclusion set $B$ with
$\gcd(D_0,\prod_{p\in B}p)=1$.
Waldspurger's nonzero local functional at $D_0$ supplies a global
holomorphic cuspidal weight-$3/2$ vector in the appropriate packet
\cite[Proposition 18]{waldspurger}; the holomorphic archimedean
member and its unary-orthogonal realization are the ones in
\cite[Theorem 1 and Proposition 1]{waldspurgerICM}.

Here is how one obtains the needed scalar character without losing
that coefficient. At an old prime choose a local vector $v$ with
Whittaker functional $\ell_{D_0}(v)\ne0$. The additive projection
\[
 P_Rv=\int_{p^{-R}\Z_p}e_p(-D_0b)\pi(n(b))v\,db
\]
retains this functional. Under the normalized unit action $M(u)$,
its value vanishes unless $u^2\equiv1\pmod{p^R}$.
Average against a finite two-power character $\alpha$ satisfying
$\alpha(-1)$ equal to the central sign. For large $R$ the surviving
neighborhoods of $1$ and $-1$ have equal nonzero contributions.
At two these are neighborhoods
$\pm1+2^{R-1}\Z_2$, not four independently chosen phases.
The average is a unit-character vector; smoothness supplies lower
unipotent invariance at sufficiently deep level.
Together these give a scalar $K_0$ character. Finite additive
projections impose the required squareclass support and preserve
the $D_0$ coefficient. Away from $S_0$ keep the spherical vector.
Its unramified coefficient $c_p^0$ is one at squarefree valuations
zero and one \cite[Section I.3, pp.~377--378]{waldspurger}.
This includes primes dividing $D_0$, so this choice preserves the
nonzero reference coefficient.

Let the resulting character be $\psi$ and put
$\eta=\psi\chi_{-\sigma}$. Both have two-power order.
Work first in the full good-Hecke eigenspace at this $S_0$-supported
level, retaining all old multiplicities. The packet identification is
\[
 \phi_0=f_E\otimes\eta,\qquad
 \chi_0=\psi\chi_{-1},\qquad
 \phi_0\otimes\chi_0^{-1}=f_E\otimes\chi_\sigma.
\]
For every vector $f$ in this whole space, Waldspurger's Corollary 2
and its same-$\Q_2$-squareclass removal of the auxiliary condition give
\[
 a_f(D)^2=\psi(D/D_0)
          \frac{L_*(\sigma D)}{L_*(\sigma D_0)}a_f(D_0)^2.
 \tag{E.10}
\]
Polarization makes the two coefficient functionals proportional on
the whole space, including its old multiplicities.

The algebraic model is the finite-dimensional space of sections on
$X_1(4N_{\rm aux})$ in \cite[Definition 3.1 and Theorem 4.2]{ramsey},
where $N_{\rm aux}$ is sufficiently divisible and supported on $S_0$.
This auxiliary level includes the depths just chosen and need not
equal the conductor $N$ of $E$.
Good-Hecke and support conditions are algebraic linear conditions.
Only flat extension of the coefficient field is used for $H^0$.
No arbitrary-ring base-change assertion at weight $3/2$ is required.
Choose an algebraic vector $g$ and normalize $a(D_0)=1$.

For an automorphism acting on the two-power characters by an odd
power $e$, put $\nu=\eta^{(e-1)/2}$. The vectors $g^\tau$ and
$g\otimes\nu$ have the same character and good eigenvalues at a
common $S_0$-supported level. The full-space proportionality in
(E.10) therefore gives
$\tau(a(D))=\nu(D/D_0)a(D)$.
Choose compatible local square roots of $D/D_0$ and let $z_D$ be
their $\eta$-value. Then $\tau(z_D)=z_D^e$, so
\[
 a(D)=z_D r_D,\quad r_D\in\Q,\quad
 r_D^2=\frac{L_*(\sigma D)}{L_*(\sigma D_0)}.                      \tag{E.11}
\]
In particular the nonzero coefficient valuations are integers.

Only now impose an extra finite coprimality sieve $B$.
Finite Fourier translations remove coefficients divisible by a
prime of $B$, without changing the desired coefficients or $D_0$.
Comparison (E.11) was proved before this enlargement.
The good recurrence is henceforth asserted only for
$p\notin S_1=S_0\cup B$:
\[
 a(p^2n)=\{\eta(p)a_p(E)-\psi(p)(-n/p)\}a(n)
                       -p\eta(p)^2a(n/p^2).                    \tag{E.12}
\]
Its square-shell multipliers are integers after removal of the
$\eta$ phase. Combining (E.9) and (E.11) gives a fixed lower bound
for $v_2a(D)-\omega(D)$ and hence an attained finite minimum.
Restricting later fresh primes by further exclusions does not
authorize choosing another minimum.

% End 07-even-source.tex


% Begin 08-odd-source.tex
\section{Actual odd theta forms and their operators}\label{app:odd-source}
We now construct the forms in (3.8), with the same $h_*$, $D_w$, $L$
and actual Kummer stages as in Section~\ref{sec:rc}.
Write $m=|h_*|$. An admissible $D$ is positive, squarefree, prime to
the fixed support and in the prescribed unit classes. Then
$K_D=\Q(\sqrt{Dh_*})$ is imaginary, split at $2N$, has no extra
units, and $Dh_*$ is an odd fundamental discriminant.
The conductor-one point is the unaveraged sum
\[
 P_D=\sum_{[\mathfrak a]\in\operatorname{Pic}(\mathcal O_{K_D})}
       \chi_{h_*}(\Norm\mathfrak a)\,
                         \pi(x_D^{[\mathfrak a]}).               \tag{O.1}
\]
It is distinct from the conductor-$|h|$ ring-class point used in
Section~\ref{sec:limits}.

\subsection{A fixed witness and one compatible sequence}
The rational two-isogeny parity theorem gives root number $-1$ for
$E$ under the full-corank-one assumption. The two clauses of
Bump--Friedberg--Hoffstein \cite[pp.~543--544]{bfh}, with all old
primes in their finite splitting set, supply first a negative
$h_*$ with nonzero central value and then a positive $D_w$ with
simple zero, after the primes of $h_*$ have been added.
The filters include the required squareclasses at $2N$; choose
$|h_*|>4$ to exclude extra units. The explicit unaveraged height
formula \cite[Theorem 1.1]{cst} applies with conductor one.
Its Rankin product has factors $L(E^{D_w},s)$ and $L(E^{h_*},s)$,
so $P_{D_w}$ is nontorsion. The forward rank-zero theorem makes the
unwanted $h_*$ component rank zero.

One integer $L\ne0$, fixed before $D$ varies, kills the finitely many
Atkin--Lehner cusp constants and all needed torsion and projection
denominators. For the torsion assertion, the points lie over fields
of degree at most four. At two fixed distinct good rational primes,
prime-to-residue-characteristic torsion injects into one of finitely
many groups $E(\F_{p^f})$, $f\le4$. Their orders bound all primary
parts. A further factor four clears the biquadratic projections.
Thus $Q=LP_{D_w}$ is a rational nontorsion point on $E^{D_w}$.

Here is a direct compatibility proof for its detecting stages.
Let $H$ be the kernel of the Tate representation and $c_Q$ the
Kummer cocycle. If $c_Q(H)=0$, its rational class descends to the
Tate image. A nonidentity central scalar $z$ gives
\[
                       (z-1)c_Q(g)=(g-1)c_Q(z),
\]
so the descended class is a coboundary, contrary to Kummer
injectivity. The closed translation image is a nonzero invariant
$\Z_2$-submodule. Non-CM open image \cite{serre} makes its rational
span the full Tate space, so it contains a full lattice.
Intersecting $H$ with the kernel of any fixed finite auxiliary
extension preserves that property.

Fix complex conjugation $\tau$, its matrix $C$, and an integral row
$b\ne0$ with $bC=b$. Choose $h\in H$ in these fixed kernels with
$v=c_Q(h)$ and $2b(v)\ne0$. For $\sigma=h\tau$ the cocycle identity gives
\[
                 c_Q(\sigma^2)=(1+C)v,\qquad b c_Q(\sigma^2)=2b(v).
\]
This is one actual Galois element. Its restrictions already agree on
all intersections: it fixes $\sqrt{D_w}$, acts nontrivially on
$\sqrt{h_*}$, and has Tate matrix $C$. At each sufficiently large
finite precision, Chebotarev applied to the common torsion,
division-point and fixed auxiliary quotient supplies a good
$\ell_0$, with a place realizing that element. Since
$\operatorname{Fr}_{\ell_0}^2=1$ on $T/2^{M_\theta}$, finite-field
Kummer gives
\[
 \lambda_{M_\theta}=b\kappa_{M_\theta}:
       E(\F_{\ell_0^2})\longrightarrow\Z/2^{M_\theta},\qquad
 v_2\lambda_{M_\theta}(\overline Q)=c_w<M_\theta                 \tag{O.2}
\]
for the fixed $c_w=v_2(2b(v))$. Its Frobenius sign is $+1$.
Later finite exclusions of rational primes do not change $c_w$.
Additional field prescriptions would require a compatible common
lift; arbitrary independent prescriptions are not asserted.

\subsection{Local functions and the theta sum}
Let $B$ be the quaternion algebra ramified at $\ell_0$ and infinity,
$R$ an Eichler order of level $N$, $V=B^0$, and $Q_B(x)=-x^2$.
The finite set
\[
                       X=B^\times\backslash\widehat B^\times/
                                             \widehat R^\times
\]
is the set of supersingular pairs. Put
$R_z=B\cap z\widehat Rz^{-1}$ and
$P_M(z)=\lambda_M(L\pi_{\ell_0}(z))$.
This is an actual function on those pairs. Its Hecke,
Atkin--Lehner and Frobenius identities hold modulo $2^M$.
Choose arbitrary integer lifts; they need not satisfy any
characteristic-zero eigenrelation. In this appendix $M=M_\theta$.

Define a finite Schwartz function $f=\prod f_p$ as follows.
At $p\nmid2Nh_*\ell_0$ use $\mathbf1_{R_p^0}$.
At $p^e\Vert N$ require an integral split quadratic embedding.
The cyclic level quotient selects one of its two summands; if $r$
is the eigenvalue of $x$ on this summand, give it weight $\alpha_p(r)$,
where the fixed two-power character satisfies
\[
                    \alpha_p(-1)=w_{p^e}\chi_{h_*}(p)^e.         \tag{O.3}
\]
Conjugation by $R_p^\times$ preserves the selected summand.
Scaling $x$ by a unit $u$ multiplies the weight by $\alpha_p(u)$.

At two impose $(1+x)/2\in R_2$ and $x^2$ a unit square.
When $2\mid N$, the operator $(1+x)/2$ is not itself asserted to be
idempotent. Its eigenvalues $(1+r)/2,(1-r)/2$ differ by the odd unit
$r$. Its integral spectral projectors distinguish the two summands,
including after reduction modulo two. For odd $u$,
\[
                  (1+ux)/2=u(1+x)/2+(1-u)/2.
\]
The summands are preserved and the selected eigenvalue of $x$
becomes $ur$. Selection is not by parity of the eigenvalue of
$(1+x)/2$. This proves the dyadic root-character rule at every
power of two in $N$.

For $p\mid h_*$ require $x\in M_2(\Z_p)^0$ with $v_p(Q_B(x))=1$.
Its reduction is nonzero nilpotent. With
$X_0=\left(\begin{smallmatrix}0&1\\0&0\end{smallmatrix}\right)$ put
\[
                 \epsilon_p(gX_0g^{-1})=(\det g/p).
\]
The centralizer has square determinant, so this is well-defined.
Conjugation has multiplier $(\det u/p)$ and scaling by $a$ has
multiplier $(a/p)$.

At $\ell_0$, the residue of a trace-zero element is in a
one-dimensional trace-zero line in $\F_{\ell_0^2}$.
Fix $\beta\ne0$ on this line and give $t\beta$, $t\ne0$, weight
$(t/\ell_0)$; give zero weight to zero.
Units act trivially on the residue field; a uniformizer acts by
Frobenius and sends $\beta$ to $-\beta$. Its multiplier is $-1$
since $\ell_0\equiv-1\pmod4$. Scaling has character $(u/\ell_0)$.
Also $-\beta^2$ is a square, and $(m/\ell_0)=1$, so the extracted
indices have $(D/\ell_0)=1$.

Extra fixed unit conditions are characteristic functions and
preserve these scaling rules. Altogether,
\[
             f(u^{-1}xu)=\chi_{h_*}(\operatorname{Nrd}u)^{-1}f(x).
\]
The product scaling character is even: its value at $-1$ is
\[
 \left(\prod_{p^e\Vert N}w_{p^e}\chi_{h_*}(p)^e\right)
                            \chi_{h_*}(-1)(-1)=1.
\]
The first factor is the positive functional sign of $E^{h_*}$.
Write $\eta_B$ for the resulting classical character, using the
inverse-idele convention; its value at a positive good prime is
the product of the finite scaling values.

Every finite $R_z^\times$ embeds by left multiplication in
$\operatorname{GL}_4(\Z)$. Reduction modulo three is injective
on a finite subgroup: a prime-order element $1+3^aA$ in its kernel
has nonzero first term in its prime-th power, modulo $3^{a+1}$,
or modulo $3^{a+2}$ for order three. Thus one may take
\[
             c_0=|\operatorname{GL}_4(\F_3)|
                       =24\,261\,120,\qquad v_2(c_0)=9.
\]
It is independent of $\ell_0$. Define
\[
 H_M(n)=c_0\sum_{z\in X}
  \frac{P_M(z)\chi_{h_*}(\operatorname{Nrd}z)}{|R_z^\times|}
              \sum_{\substack{x\in V(\Q)\\Q_B(x)=n}}f(z^{-1}xz).
 \tag{O.4}
\]
The conjugation rule makes this independent of representatives.
Its coefficients are integral in the fixed character ring.

\subsection{Transformation law and recurrence}
Use the Weil representation with self-dual measures and the standard
global additive character \cite[I.13, II.25, II.28, III, Theorem 6]{weil}.
Upper unipotents multiply by $e_p(bQ_B(x))$, so integral ones fix the
local function. If it is invariant under translations by a lattice
$L'_p$, its Fourier transform is supported in the dual lattice.
For sufficiently large $A_p$, $p^{A_p}Q_B((L'_p)^\vee)$ is integral.
Fourier conjugation therefore gives invariance under lower unipotents
of that depth. This argument also applies at two.

In a presentation $B=(-a,-b)$, $Q_B$ has diagonal coefficients
$a,b,ab$, with square product. Its diagonal Weil factor relative to
the ordinary theta multiplier cubed is consequently
$(\det Q_B,u)_p=1$. Factoring a matrix of sufficiently deep
$K_0$ into lower unipotent, diagonal and upper unipotent proves
the scalar character $\eta_B$. Conjugation by $z$ is an isometry,
so the level works in every summand. The Gaussian at infinity is
the holomorphic lowest-weight vector of weight $3/2$.
At any cusp its finite Weil transform is again Schwartz, and
positive definiteness gives only nonnegative rational exponents.
The series converges normally on compact subsets. This proves
holomorphy at every cusp, not only the formal transformation law.

Every nonzero index in (O.4) is divisible by $m$.
Replace $Q_B$ directly by $Q_B/m$ on the same Schwartz space.
It is integral on the support, so the same proof gives the actual form
\[
                  g_M(\tau)=\sum_n H_M(mn)e^{2\pi in\tau}.
\]
This justifies extraction without invoking an unspecified
half-integral-weight $U_m$ operation. The determinant squareclass
becomes $m$, hence
\[
                 \psi=\eta_B(m/\cdot),\qquad
                 \eta=\eta_B\chi_{h_*},\qquad
                 \psi/\eta=\chi_{-1}.                           \tag{O.5}
\]

For a good $p$ let $s=\eta_B(p)$. Lattice neighbors are indexed by
lines in $\F_p^2$. A nonzero trace-zero reduction has
$1+(-n/p)$ invariant lines; a nonzero nilpotent has one and zero
has $p+1$. The nonintegral contributions are $x=y/p$ with $y$
primitive of norm $p^2n$. In a basis reducing $y$ to $X_0$, the
lower-left entry is divisible by $p^2$, so exactly its kernel-line
neighbor makes $y/p$ integral. The two parts of the neighbor sum are
\[
 \begin{array}{ll}
 \text{primitive part:}&s^{-1}H_M(p^2n)-H_M(n),\\
 \text{integral part:}&(1+(-n/p))H_M(n)+psH_M(n/p^2).
 \end{array}
\]
The orbit-stabilizer weights in (O.4) validate the count even with
extra automorphisms. Multiply by the norm character
$\chi_{h_*}(p)$ and use the actual relation
$T_pP_M=a_p(E)P_M\pmod{2^M}$.
After extraction this gives
\[
 a_M(p^2n)=\{\eta(p)a_p(E)-\psi(p)(-n/p)\}a_M(n)
                 -p\eta(p)^2a_M(n/p^2)\pmod{2^M}.                \tag{O.6}
\]
No characteristic-zero eigenform assertion has entered.

\subsection{The unaveraged point comparison}
For this comparison assume also $(D/\ell_0)=1$ at the current stage.
Together with $(h_*/\ell_0)=-1$, this makes $\ell_0$ inert in $K_D$.
If $(D/\ell_0)\ne1$, the local Schwartz support instead gives
$a_M(D)=0$; no supersingular reduction formula is asserted there.
At norm $mD$ the vectors in (O.4) are exactly optimal embeddings
of $\mathcal O_{K_D}$ into $R_z$. The condition at two is the
integrality of the generator $(1+\sqrt{Dh_*})/2$.
There are two orientations at each $p^e\Vert N$, two at $\ell_0$,
and one elsewhere. The cyclic level quotient must lie in one
idempotent summand, since a nonzero quotient in both summands
would not be cyclic.

Apply \cite[Theorem 3.1, Corollary 3.2, Sections 4.9--4.10]{cornutjetchev}.
The base is $\Q$, the source algebra is $M_2(\Q)$, the changed set
is $\{\ell_0\}$, and the target is $B$. The parity condition is
$1+0+1$ even; $\ell_0$ is inert in $K_D$.
The level is prime to $\ell_0$ but may have prime powers.
Both fine conductors are one, so the ring-class ratio and the
local factor $N_{\ell_0}(0,0,0)$ are one.
Equivariance identifies the ideal-class orbit with reduction of
(O.1), retaining the level and $\ell_0$ orientations.

The stabilizer of an embedding is $\{\pm1\}$.
Its $R_z^\times$-orbit has size $|R_z^\times|/2$, so (O.4)
contributes a factor $1/2$. Reversing an old orientation multiplies
the weight and the norm-weighted point by the same sign (O.3).
Reversing the $\ell_0$ orientation gives $-1$ from each.
Thus all $2^{\omega(N)+1}$ choices contribute equally and
\[
 a_M(D)=c_0\,2^{\omega(N)}\zeta_M(D)
                  \lambda_M(\overline{LP_D})\pmod{2^M},         \tag{O.7}
\]
with a two-power root of unity $\zeta_M(D)$.
There is no class-number average. The character field is generated
by the fixed $\alpha_p$; the moving prime contributes quadratic
values only. Formula (O.2) gives fixed finite valuation
$c_w+9+\omega(N)$ at $D_w$ for every sufficiently advanced stage.

For a known simple-zero $D$, write $P_D=m_DG+T$ in its native
rank-one free lattice. The same $L$ kills $T$.
An integral finite-field Kummer map cannot decrease
$v_2(m_D)$, so (O.7) gives the lower bound by the native index.
The rational free line has index at most two in that native
lattice, by conjugation and norm. Absorb this and the fixed
initial convention in $C_{\rm det}$. The first inequality below
holds for a known simple-zero $D$; the visibility implication holds
for every admissible $D$:
\[
 \begin{split}
 v_2a_M(D)&\ge\min\{M,j_{\rm gen}(D)-C_{\rm det}\},\\
 v_2a_M(D)<M&\ \Longrightarrow\
  \ord_{s=1}L(E^D,s)=1,\quad
       j_{\rm gen}(D)\le v_2a_M(D)+C_{\rm det}.
 \end{split}                                                    \tag{O.8}
\]
If an odd-sign $D$ is not a simple zero, the height formula gives a
torsion genus point, killed by $L$, and its coefficient is zero
modulo $2^M$. Combined with the first inequality on simple-zero
parameters, this gives the all-$D$ truncated bound used below,
without assigning a rank-one lattice index to a torsion point.

\subsection{Integer divisions and actual trace operators}
All new primes below avoid the full coefficient support $S$,
including every late sieve and the current $\ell_0$.
Here $S$ is the actual coefficient support.  The coefficients of
$g$ vanish at all indices not prime to $S$.
For the odd source this is forced at $2,N,h_*,\ell_0$ by the
local Schwartz support and the extraction at $|h_*|$.
At every additional coefficient-sieve prime it is imposed by an
exact zero-or-retain Fourier sieve, which commutes with the good
recurrence outside the enlarged support.  This is the vanishing
used in the condition $(m,S)=1$ in the unary series below.
Later stationary prime-selection tests may enrich the common
label alphabet without resieving the intrinsic odd source;
such tests do not enlarge this coefficient-support condition.
Use the minimum-depth inequalities proved in Appendix~\ref{app:endpoints}.
The general bookkeeping assigns weight one to a split prime of
$\Q(E[2])$ and one-half to a prime with transposition Frobenius, called
simple. In our actual full-two-torsion family every prime is split,
$w(U)=|U|$, and the depth $\delta$ is an integer. Thus only matching
parity, meaning $e_U=0$, is needed here; the displayed general formula
also records where the unused simple-prime term would enter.
For a finite set $U$, put
\[
 n_U=\lceil\delta+w(U)\rceil,\quad e_U=n_U-\delta-w(U),\qquad
 B_U=2^{-n_U}\prod_{p\in U}(U_p-\eta(p)V_p)g,
\]
where $U_p(\sum a(n)q^n)=\sum a(pn)q^n$ and $V_pg(\tau)=g(p\tau)$.
These are operations on the actual form $g$. The degeneracy
transformation for $V_p$, and $U_p=U_{p^2}V_p$, put each new prime
at level exponent at most one, with character $\psi(p/\cdot)$.
The double coset for $U_{p^2}$ has representatives
$\left(\begin{smallmatrix}1&r\\0&p^2\end{smallmatrix}\right)$.
For a matrix in $K_0(p)$ the rule
$r'=(b+rd)/(a+rc)\pmod{p^2}$ permutes them, since $a+rc$ is a unit.
This verifies the classical bad-prime extraction at that level.

Write an index as $Jrm^2$ with $J\subset U$ and $r$ squarefree,
prime to $U$. Products of prime sets are understood here.
The originating squarefree parameter in the coefficient of $g$ is
$D=(U\setminus J)r$, not the squarefree part $Jr$ of the tested index.
Remove the phase $\eta(p)^s$ in (O.6) and put
\[
 b_{-1}=0,\quad b_0=1,\quad b_1=a_p-\kappa,\quad
 b_{s+1}=a_pb_s-pb_{s-1},\quad
 \kappa=(\psi/\eta)(p)(-D/p).
\]
For $p\in U\setminus J$ the factor $b_s-b_{s-1}$ is odd;
for $p\in J$ the factor $b_{s+1}-b_s$ is even.
Consequently, for the even source,
\[
 v_2B_U[Jrm^2]\ge w(r)
              +\tfrac12\#\{p\in J:p\text{ is simple}\}-e_U.      \tag{O.9}
\]
The right side is an integer on the allowed filter and is
nonnegative. In our full-two-torsion application all primes are
split and $w=\omega$.
For the odd source, fix an integer $R\ge1$ and a weight bound $W$
first. At one sufficiently advanced actual stage, simultaneously for
all $U$ with $w(U)\le W$ and every Fourier index, (O.9) holds with
its right side replaced by its minimum with $R$. Thus the untruncated
bound on a prescribed bounded range of squarefree weights follows by
choosing $R$ and then the stage. For the unary-removal and two-trace
tests below, choose $R\ge2$ after fixing the operator budget.
For a split prime $a_p\equiv p+1\pmod4$. The recurrence gives the
same residue on each divided shell,
\[
        (b_{s+1}-b_s)/2\equiv
     c_p(D):=\{a_p-(\psi/\eta)(p)(-D/p)-1\}/2\pmod2.              \tag{O.10}
\]

Choose $i\in\{0,1\}$ with $\eta\chi_{-4}^i$ odd and let
\[
 \Theta_J=\sum_{\substack{m>0\\(m,S)=1}}
                      \eta(m)\chi_{-4}(m)^i m q^{Jm^2}.
\]
Poisson summation gives this genuine unary weight-$3/2$ form.
Its scalar character is $\eta\chi_{-4}^i(-J/\cdot)$ and every
$p\mid J$ enters the level at exponent one.
In matching parity the shell recurrence gives
$B_U[Jm^2]-B_U[J]\eta(m)\in2\mathcal O_{\rm coef}$.
Thus
\[
                 H_U=\frac{B_U-\sum_JB_U[J]\Theta_J}{2}          \tag{O.11}
\]
is integral; in opposite parity take $H_U=B_U$.
This is division by the integer two, not reduction modulo the
maximal ideal of a ramified coefficient ring. For odd-stage
forms, the coarse bound and bounded-weight stabilization proved
below pay every fixed denominator before choosing the stage.

Multiply (O.11) and the undivided corrections by
$\theta=1+2\sum_{m>0}q^{m^2}$. They are actual weight-two forms.
For a fixed stage and $U$, choose a finite complete coefficient-field
extension containing all constituents and retain their old
multiplicities and diamond translates in the finite-dimensional
characteristic-zero space. Use modular Galois representations on
cuspidal constituents and the two character summands on Eisenstein
constituents. Holomorphy alone does not exclude the latter.
Diamond translates remain integral:
their correction coefficients have the form
\[
 \tfrac12 B_U[J]\{\psi_U(u)-\eta(u)\chi_{-4}(u)^i(-J/u)\},
\]
whose character ratio is quadratic. Good Hecke operations preserve
integrality by
$(T_qf)[n]=f[qn]+q(\langle q\rangle f)[n/q]$.
Finitely many Fourier coordinates inject this space into a finite
coordinate space. The module generated by these translates is
therefore finite over the coefficient DVR, and closed.
Chebotarev and continuity make scalar trace and determinant
operators limits of good Hecke and diamond operations, so they
preserve this module. No integral eigenspace projector is used.

\paragraph{Reading an added vertex.}
For the actual full-two-torsion family, all primes are split,
$\delta\in\mathbb Z$, and $n_U=\delta+|U|$.  Put $F_U=\theta H_U$
and let $q\notin S\cup U$ be a fresh prime, where $S$ includes
every late sieve and the current moving theta prime.
Write $u=\prod_{p\in U}p$ and let $\mathcal A(D)$ denote the
normalized valuation symbol on the tested range: it is one exactly when
$v_2a(D)=\delta+\omega(D)$ on the filter. For an odd source, that
range has a fixed weight bound chosen before the valid actual stage,
as specified before (F.7). No finite-stage value at larger weights is
asserted; the compatible limit tables used later are a separate object.
Expanding the operators, every term containing a $V_p$ has zero
coefficient at $q$, and no unary index $Jm^2$ with $J\subset U$
equals $q$.  Therefore
\[
 B_U[q]=2^{-n_U}a(uq),\qquad
 H_U[q]=2^{-n_U-1}a(uq).
\]
At coefficient one the diamond term in the good Hecke formula
vanishes.  Since $F_U\equiv H_U\pmod{2\mathcal O_{\rm coef}}$,
\[
 \overline{(t(\operatorname{Fr}_q)F_U)[1]}
   =\overline{H_U[q]}
   =\overline{\,2^{-\delta-|U|-1}a(uq)\,}
   =\mathcal A(uq)\in\mathbb F_2.
\]
Here residues lie in the coefficient residue field, with
$\mathbb F_2$ its prime subfield.  Every permitted phase is a
two-power root of unity and reduces to one.  For an odd source,
fix the finite operator budget $|U|\le W$ and $R\ge2$ first;
the stage choice in (O.9) may be taken with
$M_\theta\ge\max(0,\delta+W)+R$.
Thus $M_\theta-n_U-1\ge1$, so the normalized finite-stage error
is still divisible by the integer two.  The same reading holds
with value zero off the fixed or moving coefficient support.

The diamond correction above is a sum of integral multiples of
$\theta\Theta_J$; its fresh coefficient $q$ is zero after reduction.
It introduces neither an extra phase nor a denominator in this
reading.  All readings and subsequent omissions use this same
$g,\eta,\psi,\delta$ and actual stage; only the set $U$ changes.

A fresh trace kills every unary correction after reduction, by its
squareclass support. The exact two-by-two matrix identity then gives,
on the resulting trace orbit,
\[
 d(u)t(v)=t(v),\qquad
 t(u)t(v)=t(uv)+t(uv^{-1}),\qquad t(1)=0.
\]
Two fresh traces vanish by (O.9), with the unary division paid
first. Hence $t(gu^2)=t(g)$. The relevant trace functions factor
through the elementary quadratic quotient unramified outside
$S\cup U$. This says nothing about abelianity of the
characteristic-zero Galois representations.

At $p\in U$ a new constituent has conductor one and primitive
quadratic determinant; an old constituent is unramified there.
For the former, the conductor formula excludes nonzero monodromy:
a special representation either has unramified determinant or
conductor at least two. On cuspidal constituents, local compatibility
\cite[Theorem A]{carayol} gives one unramified line and one
quadratically ramified line. The coefficient prime is two and
$p$ is odd. The base $\Q$ has odd degree, so the extra
even-base-degree condition in that theorem is not required.
For an Eisenstein constituent the same description follows directly
from its two character summands. In either case, in a suitable basis,
\[
 \rho(\tau)=\operatorname{diag}(1,-1),\quad
 \rho(\phi)=\operatorname{diag}(\lambda,\mu),\quad U_p^{(2)}=\lambda.
\]
A direct matrix calculation proves
\[
 U_p^{(2)}\{t(g\tau\phi)-t(g\phi)\}
                   =d(\phi)\{t(g\tau)-t(g)\}.                   \tag{O.12}
\]
Old constituents have zero inertia difference on both sides.
Thus the identity holds with every old multiplicity retained.

Apply (O.12) at coefficient one.  Put $V=U\setminus\{p\}$,
$v=\prod_{r\in V}r$, so $n_U=n_V+1$.  For a fresh prime $r$ set
\[
 C_p(vr)=\frac{a_p-(\psi/\eta)(p)(-vr/p)-1}{2}\in\Z,
 \qquad c_p(vr)=C_p(vr)\bmod2.
\]
The two surviving operator terms at $pr$ and the good recurrence give
\[
 H_U[pr]\equiv
   \eta(p)C_p(vr)\,2^{-n_V-1}a(vr)\pmod{2\mathcal O_{\rm coef}}.
\]
For the even source this is an equality before reduction.  For the
odd source the error is in $2^{M_\theta-n_U-1}\mathcal O_{\rm coef}$,
which is contained in $2\mathcal O_{\rm coef}$ at our chosen stage.
Choose fresh $q,q_\tau,r_0,r_1$ representing respectively
$g,g\tau,g\phi,g\tau\phi$ in one sufficiently precise joint quotient
for the $U$ and $V$ forms and all retained characters.
The omitted form uses the same $g,\eta,\psi,\delta$ and stage,
and its level omits the entire rational prime $p$.
Thus the two $V$ symbols at $vr_0,vr_1$ agree.  Their Legendre
signs at $p$ are opposite, so $c_p(vr_0)+c_p(vr_1)=1$ in $\F_2$.
The coefficient-one reading and the reduced trace identity now give
\[
             \mathcal A_U(Uq_\tau)+\mathcal A_U(Uq)=\mathcal A_V(Vr_0)\qquad\text{in }\F_2.
\]
Here the subscript records the operator set, not a new coefficient
source or minimum.  The factor $\eta(p)$ reduces to one; no factor
two is cancelled in (O.12).  Retained quadratic labels multiply
under $g\mapsto g\phi$, including all external incidences.
This is the omission rule.  A split vertex with trivial omitted
label has trace $t(1)=0$; contracting a primary cluster to it proves
isolation.  These rules concern the actual arithmetic symbols. Their simultaneous
prime realization and finite graph budget are treated next.

% End 08-odd-source.tex


% Begin 09-normalization-endpoints.tex
\section{Nonvanishing, normalization and the final family}\label{app:endpoints}
We give the choice order underlying (3.3)--(3.9). This also separates
the two valuation arguments: the even upper transfer uses the common
Kato scalar, while the final contradiction uses the common Heegner
coordinate of Section~\ref{sec:limits}.

\subsection{An absolute bound of five new primes}
The theorem of Hoffstein--Luo \cite[p.~436]{hl} supplies fundamental
discriminants with nonzero central value and at most four prime
factors. Its finite splitting set is arbitrary. The bound four is
absolute; no uniform size threshold for those discriminants is needed.

Let $A/\Q$ be any fixed modular elliptic curve for this application,
$S$ contain two, three and all its bad primes, and $B$ be any finite
exclusion set. Fix a nonempty signed unit filter of even root number.
By CRT and Dirichlet choose one signed prime $q_0$ in its classes,
outside $S\cup B$. Apply the theorem to the actual curve $A^{q_0}$,
putting
\[
                 S'=S\cup B\cup\{|q_0|\}\cup
                               \{p:p\mid N_{A^{q_0}}\}
\]
in its splitting set. The theorem gives a fundamental $d$ with
$\chi_d(p)=1$ for all $p\in S'$ and at most four prime factors.
The functional equation forces $d>0$, since $A^{q_0}$ has even
root number and every conductor prime splits. Splitting at two
gives $d\equiv1\pmod8$. Thus $q_0d$ has all original unit
classes, is prime to $B$, is squarefree, and has at most five
prime factors. If $q_0d\equiv3$ or $7\pmod8$, its field
discriminant is $4q_0d$; only $d$ was required to be fundamental
in the theorem. This is compatible with the coefficient
construction (E.2).

For the rough companion the curve itself varies:
write $a$ for the bounded dummy product and take $A=E^{ha}$.
Choose $q_0$ so $k_0=aq_0$ is negative, $1$ modulo eight and
a square at every prime of $2Nh$, with the required exclusions.
Then $k_0$ is odd fundamental, prime to $N_{E^h}$, and
\[
          w(E^{hk_0})=w(E^h)\chi_{k_0}(-N_{E^h})=+1.
\]
This proves the sign directly for the actual application.
Apply Hoffstein--Luo to $E^{hk_0}$ with every conductor,
old, $h$, dummy and forbidden prime in its splitting set.
The resulting $k'=k_0d$ adds at most five new primes.
The constants and sizes in the analytic theorem may depend on
this varying curve and the whole finite set; the bound five does not.

\subsection{The even upper transfer}
For one fixed even source write (E.11) as
\[
             v_2L_*(\sigma D)=2v_2a(D)+\kappa_e.
\]
The lower bound (E.9) and a nonzero reference coefficient imply
that the nonempty set of integers
$v_2a(D)-\omega(D)$ has a finite attained minimum $\delta_e$.
A unit of the normalized symbol means precisely equality to this
minimum. It therefore gives
\[
 L_*(\sigma D)\ne0,\qquad
 v_2L_*(\sigma D)=2\omega(D)+C_{e,\rm unit},
 \quad C_{e,\rm unit}=2\delta_e+\kappa_e.                         \tag{F.1}
\]
This constant belongs to the fixed source, not to a curve later
introduced solely to obtain nonvanishing.

For the negative types used in the rough companion, the unsigned
index is $7$ modulo eight. Its total $(-1/D)$ label is nonzero, so
every unit seed is nonempty. A seed of least positive size has no
dependence on internal edge bits: a nonzero edge derivative
would be a smaller nonempty unit. Its labels therefore suffice.
If a list with the same total label contains the first $r-1$
seed anchors and a nonempty remainder, contract that remainder
into the last anchor. Its total label is forced to be the
remaining seed label. This proves partitionability regardless
of how many old $h$-primes are in the remainder.

For a partitionable $d$ already known to have Selmer corank zero,
the one-symbol address construction keeps every prime of $d$
stationary and gives arbitrarily large cubes $d_x$, $d_0=d$,
with unit symbol at each $x\ne0$. The scalar in (E.5) is common
to that cube. By (E.8) and (F.1),
\[
      v_2u_x^{\rm ev}(0)\le B_e:=C_{e,\rm unit}+C_{\rm loc}
                                      \quad(x\ne0).
\]
Choose $M-c_E>B_e$. The constant coefficient of
$2^{c_E}u^{\rm ev}\in\Lambda[G_b]$ has sign-value bits of total
Boolean degree at most $2^{M-1}-1$. For a larger cube dimension,
the Boolean comparison in Section~\ref{sec:limits} gives
$x\ne0$ with
$v_2(u_x^{\rm ev}(0)-u_0^{\rm ev}(0))\ge M-c_E$.
If the base were zero or had valuation greater than $B_e$, this
would contradict the upper bound at $x$. Therefore
\[
       v_2L_*(d)-2\omega(d)-s(d)
                       \le C_{e,\rm unit}+2C_{\rm loc}.          \tag{F.2}
\]
The required dimension may depend on $d$; this bound does not.
It precedes the odd minimum and uses no Heegner transfer.

Prepare these sources and nonempty seeds for every total type at
the fixed support. There are finitely many. A bounded supply of
dummy anchors and reserves works for all of them. For each $h$
realize that supply by distinct primes avoiding the entire fixed
support and $h$. Put every dummy and every other already fixed
forbidden prime in the exclusion set of the preceding
nonvanishing construction. It gives $k'$ with
$\omega(k')\le A$, where $A$ is the fixed dummy count plus five.
For the total type that actually occurs, $hk'$ is partitionable.
Nonvanishing gives corank zero by the forward theorem, so (F.2)
applies. Taking the maximum over the finitely many prepared
constants proves
\[
 v_2L_*(hk')\le2\omega(h)+2\omega(k')+s(hk')+C_E.                 \tag{F.3}
\]
No new even minimum is chosen after $h$ or its exclusion set.
Those excluded primes are not added as relaxed places in a
Selmer complex.

\subsection{The odd minimum is chosen after its lower bound}
The arithmetic result (3.5) and the two native-lattice height
comparisons (3.4) give, using (F.3),
\[
 2j_{\rm gen}(h)\ge2\omega(h)+s(h)-C_{\rm gen},\qquad
 C_{\rm gen}=2A+C_R+C_E-v_2L_*(h_*)+C_H.                         \tag{F.4}
\]
The variable $s(hk')$ cancels. This estimate precedes the theta
prime, precision, odd minimum, final field and binary dimension.

Combining (F.4) with (O.8) gives on all admissible $D$
\[
            v_2a_M(D)\ge\min\{M,\omega(D)-C_{\rm coarse}\},
 \qquad C_{\rm coarse}=C_{\rm gen}/2+C_{\rm det}.                 \tag{F.5}
\]
At visible coefficients (O.8) first gives a simple zero, allowing
(F.4); at the other odd-sign parameters the coefficient is zero
modulo $2^M$. This avoids a circular use of the future minimum.

Fix one nonprincipal ultrafilter on the actual stages with
$M=M_\theta\to\infty$. Consider the integers $r$ for which a
bounded-weight sequence of visible $D_M$ has
$v_2a_M(D_M)-\omega(D_M)=r$ on a set in the ultrafilter.
The fixed witness $D_w$ and (O.2),(O.7) make this set nonempty.
Equation (F.5) bounds it below. Let its attained minimum be
$\delta_o$. For every fixed $W$, on a set in the same ultrafilter,
\[
 v_2a_M(D)\ge\omega(D)+\delta_o
             \quad\hbox{for every visible }D,\ \omega(D)\le W.
 \tag{F.6}
\]
Otherwise choose a violating $D_M$ at each bad stage. One of the
finitely many integers between $-C_{\rm coarse}$ and
$\delta_o-1$ would occur on an ultrafilter-large set, a contradiction.
Invisible coefficients satisfy the required truncated inequality
once $M>W+\delta_o$.

For a fixed finite operator budget $|U|\le W_U$, put
$N_0=\max(0,\delta_o+W_U)$ and fix $R\ge2$.
Choose an integer $T\ge W_U+1$ with
$T-C_{\rm coarse}\ge N_0+R$, then choose a stage in the
ultrafilter-large set where (F.6) holds through weight $T$, with
$M\ge N_0+R$ and $M>T+\delta_o$.
For the index $Jrm^2$ put $D=(U\setminus J)r$.
The integral shell recurrence has an error in $2^M\mathcal O$.
If $\omega(D)>T$, (F.5) gives
\[
 v_2\bigl(2^{n_U}B_U[Jrm^2]\bigr)\ge N_0+R.
\]
If $\omega(D)\le T$, (F.6) for visible coefficients and
$v_2a_M(D)\ge M$ for invisible ones instead give
\[
 v_2\bigl(2^{n_U}B_U[Jrm^2]\bigr)
       \ge\min\{M,n_U+\omega(r)\}.
\]
Consequently $v_2B_U[Jrm^2]\ge\min\{R,\omega(r)\}$ for every
Fourier index and every $U$ in the budget at this one stage.
At $r=1$ the normalized originating coefficient
$2^{|J|-n_U}a_M(U\setminus J)$ is integral, so (O.10) proves
the unary congruence modulo the integer two.
The additional division in (O.11) leaves error valuation at least
$R-1\ge1$.  In particular $M-n_U-1\ge1$ for every reading and
first-shell omission in the budget.
It does not select one stage for infinitely many budgets.

More precisely, fix a squarefree-output cutoff $L$. On a set of
stages in the same ultrafilter, (F.6) holds through $L$ and
$M_\theta>\delta_o+L$. Define $\mathcal A_{M_\theta,L}(D)$ only
for positive squarefree $D$ with $\omega(D)\le L$: it is zero off the full coefficient filter,
one when $v_2a_{M_\theta}(D)=\delta_o+\omega(D)$ on that filter,
and zero otherwise. This is exhaustive on the tested range:
visible coefficients satisfy (F.6), while invisible ones have
valuation at least $M_\theta>\delta_o+\omega(D)$.
In particular a unit in this range is visible and satisfies
\[
 \ord_{s=1}L(E^D,s)=1,\qquad
        j_{\rm gen}(D)\le\omega(D)+C_o,\quad C_o=\delta_o+C_{\rm det}.
 \tag{F.7}
\]
For the preceding operator budget take $L\ge W_U+1$ and include
its valid-stage requirements; all fresh-vertex and omission outputs
then lie in this range. The constant $C_o$ is independent of $L$
and of the stage. No value is assigned by this definition at
unbounded weights of one finite stage. No minimum is changed when
a vertex is omitted or a stationary condition is later added.

\subsection{Nonempty rational seed and a common finite stage}
Choose the final negative even unit $k$ as in
Section~\ref{sec:rc}, with fresh ramification outside the old odd
label and torsion fields. Fix $K=\Q(\sqrt k)$ before the binary
dimension. The forward theorem gives corank zero for $E^k$;
the sum-and-difference isogeny, on compact Selmer modules tensored
with $\Q_2$, gives full corank one for $E/K$.
The intrinsic odd symbol omits the primes of $k$.
Later $\ell_0$ avoids them. Fresh ramification implies that every
old odd label still lifts to a $K$-split label.
Additional stationary tests can have relations; use their actual
common alphabet and its separate surjective projections.

The final $E^k$ source is transported from the fixed negative source
in which $k$ was a unit. For $(D,k)=1$, its coefficient comparison gives
\[
 v_2a_{\rm new}(D)=v_2a_{\rm old}(kD)-v_2a_{\rm old}(k)
                                                    \ge\omega(D).
\]
Thus normalization at index one has this specific inherited lower
bound; nonvanishing at index one alone would not prove it.
A unit has equality and hence
$v_2L_*(h_xk)\le2\omega(h_x)+C_{e,\rm final}$, where the conversion
from the fixed $E^k$ period convention changes only a fixed constant.
Its $K$-split seed must be nonempty. Suppose otherwise. The integral
coefficient valuations then improve the preceding lower bound to
$\omega(D)+1$ on every nonempty product of allowed $K$-split primes.
For a fresh such $p$, apply the trace identity to
$f=\theta H_{\{p\}}$:
\[
 U_p^2t(\phi)=U_p^3+d(\phi)U_p,\qquad
 (t(\phi)f)[p^2]=f[p^3]+(d(\phi)f)[p].
\]
Choose a fresh $K$-split $p'$ matching $\phi$ in the finite Galois
quotient and all fixed rational characters. The product $pp'$ is in
the prescribed square filter even when $p$ alone is not.
The left coefficient is a unit times $(a_p-1)a(pp')/4$, hence zero
modulo two by the improved bound. The determinant term at $p$
vanishes because $B_{\{p\}}$ and its $\Theta_{\{p\}}$ correction
have the same character and cancel there; $\Theta_1$ has no such
coefficient. The remaining $p^3$ coefficient is a unit times
\[
                         (a_p-1)(a_p-p-1)/4.
\]
Here $a_p=a_p(E^k)$ is even, so $a_p-1$ is odd. Its vanishing
modulo two forces $a_p(E^k)\equiv p+1\pmod8$ for every such $p$.
Choose such a $p$ with $p\equiv1\pmod4$ and $(\ell/p)=-1$.
Freshness of $k$ permits these simultaneous conditions.
The three rational two-torsion points have no rational halves:
their necessary square differences include
$-\ell,\ell,-4\ell$. The two-primary reduction group is exactly
of order four. This contradicts the asserted congruence.
This proves the required nonempty even seed. If the intrinsic odd
symbol already has a unit at the empty configuration, its detection
formula gives a simple zero for $E$ and the converse is finished.
Otherwise take a least nonempty odd unit seed. Thus the remaining
two-symbol address argument has two nonempty seeds.

Here the finite graph argument has a definite arithmetic order.
The moving theta prime contributes only a bounded number of quadratic
label coordinates. The root characters at the old level were fixed
before $M_\theta$; growing Kummer precision introduces no growing-order
character on the variable primes. Thus the elementary label groups,
their projections and their finite addition tables have bounded size.
Choose their abstract type on the same ultrafilter as $\delta_o$.
For each fixed weight cutoff $W$, restrict first to the
ultrafilter-large set of stages where the preceding budget and
precision conditions hold for every $|U|\le W$, with $R=2$.
On those stages the coefficient reading, factorization, contraction
and isolation rules hold on all configurations of at most $W+1$
vertices, with the stated nonempty-output restriction.  They therefore define actual symbol truth tables on
the compatible typed graph data.  There are finitely many such
tables; select the one occurring on a set in the same ultrafilter.
Two selected tables agree on their common restriction by taking
the finite intersection of their valid-stage sets. Their union defines
the odd limit symbol on compatible finite typed graphs of every size;
it is not a total valuation symbol at any one finite stage.
Every seed or contraction identity uses only finitely many such graphs
and is inherited at a cutoff containing all of them.
The bounded-weight witness remains a witness in these tables.
No stage common to all cutoffs is selected.
After fixing $k$, include its tests and choose the common projection
data on the same filter, without redefining the intrinsic odd minimum.
Stabilize the finite alphabet and its two projections in this way. Repeated packets of total label
zero contain anchors for both seeds. Product-label contractions
include all external incidences. Grounded Laplacian minors
provide both the singular and varying address networks, including
the singular-matrix case. Take the product of the two highest
forms in the ordinary polynomial ring before Boolean reduction.
An exponent two assigns one singular program to each factor;
an exponent one assigns a varying program to one factor.
The total degree forces the intended allocation and leaves
coefficient one. Thus one finite graph gives units at every
nonzero address for both symbols. Every contracted output is
nonempty.


For completeness the address networks can be given by explicit matrices
over $\F_2$. At a target $x_*\ne0$, choose a dual basis $\lambda_i$
with $\lambda_i(x_*)=1$. Use three primary sites per group, with
basis $(e_i,u_i,v_i)$ and $e_i$ the all-one physical vector.
Give the $e_i$ orthonormal pairing and on the remaining coordinates use
\[
                    H(T,B)=\begin{pmatrix}0&T\\T^{\mathsf t}&B\end{pmatrix}.
\]
For the full grounded matrix $Q$, the columns $h_i=Qe_i$ are supported
on group $i$. Attach them to terminals activated by $\lambda_i$.
The row-sum identities then hold on every active group restriction.
To attach a further column $c$ to a pair activated by $\alpha,\beta$,
also attach it to an $\alpha+\beta$ terminal when $\alpha\ne\beta$.
Its total added column is zero at every address, and that balancing
terminal is absent at the target.

Cauchy--Binet applied to incidence matrices gives the constant
coefficient $\det Q$ and pair-contraction coefficient
$c_a^{\mathsf t}\operatorname{adj}(Q)c_b$.
The nonzero minors are precisely forests with the indicated deleted
roots. This proves the formula even for singular $Q$.
Cycles have zero iterated difference: contract a spanning path and
the closing edge becomes internal, changing neither product labels
nor external incidences. A component of primary auxiliaries alone
contracts to an isolated neutral vertex and vanishes.

For the singular program take $T=I+P$, with $P$ cyclic, and let $B$
have one diagonal entry one. A proper group restriction has invertible
$T$; on all groups the radical is $\sum_i u_i$ and has dimension one.
Thus $\operatorname{adj}(Q)=rr^{\mathsf t}$ there.
A column dual to $u_1$ has pairing one with $r$, while the anchors
have pairing zero. The program has constant zero and exactly the
chosen pair difference in degree one at $x_*$; elsewhere its constant
is one. The same computation covers a single group.

For the varying program take $T=I+L$, where $L$ is the consecutive
lower shift, and $B_\varepsilon=\varepsilon E_{bb}$.
All restrictions are invertible. The change in the upper-left inverse
block is measured by $((T_J^{-1})_{b1})^2$, which is one exactly when
the full path is present. Hence the difference of the two programs
has the chosen degree-one difference only at $x_*$ and degree at
least two elsewhere; both have constant one everywhere.

Now expand the product of the two final evaluations and sum over
the binary choices of the varying programs. A term using only a
varying program's constants cancels. A singular program costs at
least one degree in each factor at its target, and a surviving
varying program costs at least one. The sum of these minimum costs
is the total available degree of the two top forms. Thus every
surviving term has exactly the intended cost at the intended
address; higher and wrong-address contractions vanish.
The remaining coefficient is the selected nonzero coefficient
of the ordinary product of top forms, hence one.
This proves that one choice works at all nonzero addresses.
It also explains why Boolean reduction must follow, rather than
precede, the top-form product.

For each chosen $b$, this construction has a finite budget $W_b$.
Enlarge it to bound every graph product and squarefree output in its
finite auxiliary tests, as well as every operator set. Contractions
only decrease the number of vertices. Thus every final $h_x$ lies in
the tested range of (F.7); this enlargement changes no fixed constant.
Only now take one actual theta stage in the finite intersection
of the selected table sets through $W_b+1$, the valid-stage set
for the full budget $|U|\le W_b$ with $R=2$, and the required
precision and finite exclusion sets.  In particular (F.6) holds
through the threshold $T$ used for that budget.
Realize its prescribed forward residue bits by distinct primes using CRT and Dirichlet;
quadratic reciprocity supplies the reverse bits.
There is no uniform bound on $W_b$, no intersection over all
budgets, and no product-surjectivity assumption on the alphabets.
The resulting one family satisfies (3.9), with constants fixed
before $b$. The theta prime $\ell_0$ is a coefficient prime and
an exclusion, not an active relaxed Selmer prime or a moving
arithmetic prime $r_i$.

\subsection{One coordinate in the final contradiction}
For this fixed $K$, (3.4), (F.7), and the even bound just proved give
\[
 2j_{\rm rc}(h_x,k)\le4\omega(h_x)+C_J,\qquad
 C_J=2C_o+C_{e,\rm final}-v_2L_*(h_*)+C_{H,\rm final}
\]
at every $x\ne0$.
The finite support comparison uses the same class, duality map,
bases and coefficient stages. Its integral correction
\[
 \eta=\prod_q(1+v_q(1+g_q)),\qquad l_q/e_q^2=1+2v_q,
 \qquad \widetilde u=\eta u_{\rm rel}
\]
is formed before characters. The actual integral determinant
triangle gives $\widetilde u_x=\epsilon_xu_{{\rm prim},x}$ with
$\epsilon_x\in\Lambda^\times$, rather than independently chosen
units in $\mathscr O^\times$. Therefore the central identity is
\[
 \begin{split}
 v_2\widetilde u_x(0)={}&2j_{\rm rc}(h_x,k)-4\omega(h_x)
       -s(h_x)-s(h_xk)\\
       &-\sum_{p\mid k}t_p+\epsilon'_x,\qquad
 0\le t_p\le2,\quad |\epsilon'_x|\le C_{\rm cal}.
 \end{split}                                                    \tag{F.8}
\]
The nonnegative finite $\Sha$ terms give one upper bound
$v_2\widetilde u_x(0)\le B$ for all $x\ne0$, independent of $b$.
At the unknown base no finiteness of $\Sha$ is assumed.
The retained central class and corank-one argument instead give
$\widetilde u_0(0)=0$ if the base derivative vanishes.

The fixed-field construction of Section~\ref{sec:limits} has
\[
 \begin{gathered}
 c_{\rm det}=a+4(h_K+2c_E+2t_K),\quad
 L_T(M)=16(d_K+5)M,\\
 f_x\widetilde u_x=d_x,\qquad
 c=v_2f_0(0)\le c_{\rm det}.
 \end{gathered}
\]
Choose $M>\max(c_{\rm det},B+2c_{\rm det},0)$, then
$b>(L_T(M)+2)(2^M-1)$ large enough for the address construction,
then its finite graph budget and actual theta stage $M_\theta$.
The residual comparison tests coefficients of $f$ from
$-L_T(M)$ through zero and only the constant coefficient of $d$.
It gives $x\ne0$ with
\[
       v_2(\widetilde u_x(0)-\widetilde u_0(0))\ge M-2c>B.
\]
The estimate uses both $v_{\mathscr O}(f_x)\le c$ and
$\widetilde u_x\in\Lambda[1/2]$; neither a bound on the whole
numerator nor a finite-Artin inverse of $t$ is asserted.
This contradicts (F.8) under the supposed base zero.
The Rankin factorization and $L(E^k,1)\ne0$ give the asserted
simple zero of $E$.

This completes the proof of the restricted converse,
Theorem~\ref{hyp:RC}, from the displayed classical inputs and
the arithmetic, coefficient and determinant constructions.
The conversion from that converse to pole witnesses, the integer
model, the natural compiler and finally the ordinary rational value
is the separate argument in the main text.

% End 09-normalization-endpoints.tex


% Begin references.tex
\begin{thebibliography}{99}
\bibitem{babagranath} S. Baba and H. Granath,
\emph{Genus 2 curves with quaternionic multiplication}, 2005 preprint,
Sections 2.1, 3 and Proposition 10.
\url{https://webdoc.sub.gwdg.de/ebook/serien/e/mpi_mathematik/2005/18.pdf}.
\bibitem{buzzard} K. Buzzard, \emph{Integral models of certain Shimura
curves}, 1996, corrected 2002, Theorem 2.1 and Corollary 2.3.
\url{https://www.ma.imperial.ac.uk/~buzzard/maths/research/papers/shimura.pdf}.
\bibitem{linyang} Y.-H. Lin and Y. Yang, \emph{Quaternionic loci in
Siegel's modular threefold}, arXiv:1807.00466v1, 2018, equations (2), (14).
\url{https://arxiv.org/abs/1807.00466v1}.
\bibitem{vkk} R. von K\"anel and A. Kret, \emph{Integral points on coarse
Hilbert moduli schemes}, arXiv:2307.06944v1, 2023, equations (7.1),
(7.13), (7.17), Proposition 7.13 and the maximal-order argument in
Section 7.3.3. \url{https://arxiv.org/abs/2307.06944v1}.
\bibitem{lny} K. Lauter, M. Naehrig and T. Yang, \emph{Hilbert theta
series and invariants of genus 2 curves}, Journal of Number Theory,
2015, Section 2.2, doi:10.1016/j.jnt.2015.02.020.
\url{https://people.math.wisc.edu/~tonghaiyang/2015-LNYFinal-JNT.pdf}.
\bibitem{pazuki} F. Pazuki, \emph{Theta height and Faltings height},
Bulletin de la Soci\'et\'e Math\'ematique de France, 2012,
Sections 2.2--2.3 and Corollary 1.3.
\url{https://www.numdam.org/item/10.24033/bsmf.2623.pdf}.
\bibitem{poonenstoll} B. Poonen and M. Stoll, \emph{The Cassels--Tate
pairing on polarized abelian varieties}, Annals of Mathematics 150
(1999), 1109--1149. \url{https://math.mit.edu/~poonen/papers/sha.pdf}.
\bibitem{greentao} B. Green and T. Tao, \emph{Linear equations in primes},
Annals of Mathematics 171 (2010), 1753--1850.
\url{https://annals.math.princeton.edu/2010/171-3/p08}.
\bibitem{moebius} B. Green and T. Tao, \emph{The M\"obius function is
strongly orthogonal to nilsequences}, Annals of Mathematics 175 (2012).
\url{https://annals.math.princeton.edu/2012/175-2/p03}.
\bibitem{gtz} B. Green, T. Tao and T. Ziegler, \emph{An inverse theorem
for the Gowers $U^{s+1}[N]$-norm}, arXiv:1009.3998v5, 2026 corrected version.
\url{https://arxiv.org/abs/1009.3998v5}; author erratum:
\url{https://people.maths.ox.ac.uk/greenbj/papers/GI(s)-erratum.pdf}.
\bibitem{bmor} M. A. Bennett, G. Martin, K. O'Bryant and A. Rechnitzer,
\emph{Explicit bounds for primes in arithmetic progressions},
Illinois Journal of Mathematics 62 (2018), 427--532,
arXiv:1802.00085v3, Theorem 1.2.
\url{https://arxiv.org/abs/1802.00085v3}.
\bibitem{milne} J. S. Milne, \emph{Arithmetic duality theorems},
second edition, 2006, Chapter I.
\url{https://www.jmilne.org/math/Books/ADTnot.pdf}.
\bibitem{nekovar} J. Nekov\'a\v r, \emph{Selmer complexes},
Ast\'erisque 310, 2006, Proposition 5.2.4, Theorems 6.3.4 and 7.8.8, and Sections 5.7 and 6.4.
\url{https://www.numdam.org/item/AST_2006__310__1_0/}.
\bibitem{afp} Archive of Formal Proofs, \emph{DPRM\_Theorem},
\texttt{Alpha\_Sequence}, \texttt{Exponential\_Relation}, and
\texttt{Exponentiation}. Exact source versions are pinned in the local
corpus; current browse pages are not substitutes for those pins.
\url{https://isa-afp.org/entries/DPRM_Theorem.html}.
\bibitem{silverman} J. H. Silverman, \emph{The arithmetic of elliptic
curves}, second edition, Springer, 2009, Chapters IV and VII.
\bibitem{dd} T. Dokchitser and V. Dokchitser, \emph{Root numbers and parity
of ranks of elliptic curves}, Journal f\"ur die reine und angewandte
Mathematik 658 (2011), 39--64, Theorem 1.8.
\url{https://arxiv.org/abs/0906.1815v1}.
\bibitem{gz} B. Gross and D. Zagier, \emph{Heegner points and derivatives
of $L$-series}, Inventiones Mathematicae 84 (1986), 225--320.
\bibitem{kolyvagin} V. Kolyvagin, \emph{Euler systems}, in
\emph{The Grothendieck Festschrift}, Volume II, 1990.
Preprint MPI/90-69, Theorems 1 and 5.
\url{https://archive.mpim-bonn.mpg.de/id/eprint/325/1/preprint_1990_69.pdf}.
\bibitem{kato} K. Kato, \emph{$p$-adic Hodge theory and values of zeta
functions of modular forms}, Ast\'erisque 295 (2004), 117--290.
\url{https://www.numdam.org/article/AST_2004__295__117_0.pdf}.
\bibitem{bfh} D. Bump, S. Friedberg and J. Hoffstein,
\emph{Nonvanishing theorems for $L$-functions of modular forms and their
derivatives}, Inventiones Mathematicae 102 (1990), 543--618.
\url{https://wstein.org/papers/bib/bump-friedberg-hoffstein-nonvanishing.pdf}.
\bibitem{hl} J. Hoffstein and W. Luo, \emph{Nonvanishing of $L$-series and
the combinatorial sieve}, Mathematical Research Letters 4 (1997), 435--444.
\url{https://intlpress.com/site/pub/files/_fulltext/journals/mrl/1997/0004/0003/MRL-1997-0004-0003-a012.pdf}.
\bibitem{desjardins} J. Desjardins, \emph{Root number of twists of an
elliptic curve}, Journal de Th\'eorie des Nombres de Bordeaux
32 (2020), 73--101. \url{https://jtnb.centre-mersenne.org/item/10.5802/jtnb.1112.pdf}.
\bibitem{cst} L. Cai, J. Shu and Y. Tian, \emph{Explicit Gross--Zagier
and Waldspurger formulae}, Algebra \& Number Theory 8 (2014).
\url{https://arxiv.org/abs/1408.1733}.
\bibitem{ccl} L. Cai, Y. Chen and Y. Liu, \emph{Heegner points on modular
curves}, Transactions of the American Mathematical Society 370 (2018),
3721--3743. \url{https://archive.ymsc.tsinghua.edu.cn/pacm_download/632/12220-Heegner_points_on_modular_curves.pdf}.
\bibitem{cornutjetchev} C. Cornut and D. Jetchev,
\emph{Liftings of reduction maps for quaternion algebras}, 2010 preprint.
\url{https://webusers.imj-prg.fr/~christophe.cornut/papers/liftings.pdf}.
\bibitem{waldspurger} J.-L. Waldspurger, \emph{Sur les coefficients de
Fourier des formes modulaires de poids demi-entier},
Journal de Math\'ematiques Pures et Appliqu\'ees 60 (1981), 375--484,
Proposition 18 and Corollary 2.
\bibitem{ramsey} N. Ramsey, \emph{Geometric and $p$-adic modular forms
of half-integral weight}, Annales de l'Institut Fourier 56 (2006),
599--624. \url{https://www.numdam.org/item/10.5802/aif.2195.pdf}.
\bibitem{carayol} H. Carayol, \emph{Sur les repr\'esentations
$\ell$-adiques associ\'ees aux formes modulaires de Hilbert},
Annales Scientifiques de l'\'Ecole Normale Sup\'erieure 19 (1986), 409--468,
Theorem A. \url{https://www.numdam.org/item/ASENS_1986_4_19_3_409_0/}.
\bibitem{waldspurgerICM} J.-L. Waldspurger,
\emph{Correspondances de Shimura}, Proceedings of the International
Congress of Mathematicians, 1983, 525--531.
\url{https://www.imj-prg.fr/wp-content/uploads/2020/prix/waldspurger1983.pdf}.
\bibitem{weil} A. Weil, \emph{Sur certains groupes d'op\'erateurs unitaires},
Acta Mathematica 111 (1964), 143--211, I.13, II.25, II.28, III, Theorem 6.
\url{https://webhomes.maths.ed.ac.uk/~v1ranick/papers/weil2.pdf}.
\bibitem{serre} J.-P. Serre, \emph{Propri\'et\'es galoisiennes des points
d'ordre fini des courbes elliptiques}, Inventiones Mathematicae 15
(1972), 259--331.
\bibitem{mazur1995} B. Mazur,
\emph{Speculations about the topology of rational points: an up-date},
Ast\'erisque 228 (1995), 165--181, Conjecture 3, p.~165.
\url{https://www.numdam.org/article/AST_1995__228__165_0.pdf}.
\bibitem{matiyasevich1971} Yu. V. Matiyasevich,
\emph{Diophantine representation of enumerable predicates},
Mathematics of the USSR-Izvestiya 5 (1971), 1--28,
Basic Theorem, p.~3; integer-witness conversion, p.~2.
\url{https://doi.org/10.1070/IM1971v005n01ABEH001004}.
\bibitem{openaihtp} OpenAI,
\emph{Hilbert's tenth problem over the rational numbers},
OpenAI Math Release preprint, 6 October 2026.
\url{https://github.com/openai/math/blob/main/preprints/Hilberts-tenth-problem-over-the-rational-numbers-October-6-2026/main.pdf}.
\end{thebibliography}

% End references.tex

\section*{AI usage statement}
This manuscript was developed primarily by OpenAI Codex agents (Astra 6 and Sol 6.1)
and Anthropic Claude agents (Opus 5.5). Their contributions included mathematical
exploration, construction and analysis of the polynomial, proof drafting,
literature searches, computational checks, internal review, and drafting
and revision of the exposition. Reviews by multiple AI agents formed part
of this process; they do not constitute independent human peer review.
\end{document}
